\documentclass[11pt]{article}

\usepackage[T1]{fontenc}
\usepackage{lmodern}
\usepackage[letterpaper,margin=1in]{geometry}
\usepackage{microtype}
\usepackage{mathtools}
\usepackage{amssymb}
\usepackage{amsthm}
\usepackage{graphicx}
\usepackage[font=small,labelfont=bf,skip=6pt]{caption}
\usepackage{needspace}
\usepackage{colortbl}
\usepackage{tikz}
\usetikzlibrary{arrows.meta,positioning}
\usepackage{aliascnt}
\usepackage[hidelinks]{hyperref}
\usepackage[capitalise,nameinlink,noabbrev]{cleveref}
\usepackage{authblk}

\renewcommand\Authfont{\normalsize}
\renewcommand\Affilfont{\small}
\setlength{\affilsep}{0.5em}

\Urlmuskip=0mu plus 1mu\relax

\allowdisplaybreaks[2]

\theoremstyle{plain}
\newtheorem{theorem}{Theorem}[section]

\newaliascnt{lemma}{theorem}
\newtheorem{lemma}[lemma]{Lemma}
\aliascntresetthe{lemma}

\newaliascnt{proposition}{theorem}
\newtheorem{proposition}[proposition]{Proposition}
\aliascntresetthe{proposition}

\newaliascnt{corollary}{theorem}
\newtheorem{corollary}[corollary]{Corollary}
\aliascntresetthe{corollary}

\newaliascnt{claim}{theorem}
\newtheorem{claim}[claim]{Claim}
\aliascntresetthe{claim}

\theoremstyle{definition}
\newaliascnt{definition}{theorem}
\newtheorem{definition}[definition]{Definition}
\aliascntresetthe{definition}

\newaliascnt{procedure}{theorem}
\newtheorem{procedure}[procedure]{Decision Procedure}
\aliascntresetthe{procedure}

\theoremstyle{remark}
\newaliascnt{remark}{theorem}
\newtheorem{remark}[remark]{Remark}
\aliascntresetthe{remark}

\crefname{theorem}{Theorem}{Theorems}
\Crefname{theorem}{Theorem}{Theorems}
\crefname{lemma}{Lemma}{Lemmas}
\Crefname{lemma}{Lemma}{Lemmas}
\crefname{proposition}{Proposition}{Propositions}
\Crefname{proposition}{Proposition}{Propositions}
\crefname{corollary}{Corollary}{Corollaries}
\Crefname{corollary}{Corollary}{Corollaries}
\crefname{claim}{Claim}{Claims}
\Crefname{claim}{Claim}{Claims}
\crefname{definition}{Definition}{Definitions}
\Crefname{definition}{Definition}{Definitions}
\crefname{procedure}{Procedure}{Procedures}
\Crefname{procedure}{Procedure}{Procedures}
\crefname{remark}{Remark}{Remarks}
\Crefname{remark}{Remark}{Remarks}

\newcommand{\AlgNums}{\overline{\mathbb Q}}
\newcommand{\cF}{\mathcal F}
\newcommand{\WF}{W_{\cF}}
\newcommand{\CC}{\mathsf{CC}}
\newcommand{\BO}{\mathsf{BO}}
\newcommand{\TP}{\mathsf{TP}}
\newcommand{\Mal}{\mathsf{Mal}}
\newcommand{\CSP}[1]{\mathsf{\#CSP}(#1)}
\newcommand{\CSPdeg}[2]{\mathsf{\#CSP}^{#1}(#2)}
\newcommand{\WFdeg}[1]{W_{\cF}^{#1}}
\newcommand{\SuppRel}[1]{\Phi_{#1}}
\newcommand{\RowEq}[1]{\Omega_{#1}}
\newcommand{\Pow}[2]{#1^{\langle #2\rangle}}
\newcommand{\zvec}[1]{\mathbf z_{#1}}
\newcommand{\InstMonoid}[1]{\mathfrak M_{#1}}
\newcommand{\eps}{\varepsilon}
\newcommand{\emptytuple}{\epsilon}
\newcommand{\YES}{\textnormal{\textsc{yes}}}
\newcommand{\NO}{\textnormal{\textsc{no}}}
\DeclareMathOperator{\Pure}{Pure}
\DeclareMathOperator{\arity}{arity}
\DeclareMathOperator{\supp}{supp}
\DeclareMathOperator{\typeop}{type}
\DeclareMathOperator{\Row}{Row}
\DeclareMathOperator{\pr}{pr}
\DeclareMathOperator{\im}{im}
\DeclareMathOperator{\Tor}{Tor}
\DeclareMathOperator{\lcm}{lcm}
\DeclarePairedDelimiter{\abs}{\lvert}{\rvert}
\DeclarePairedDelimiter{\norm}{\lVert}{\rVert}

\begin{document}

\title{From Block Orthogonality to 
Decidability \\[0.15em] in Complex-Weighted Counting CSP}

\author[1]{Chenghua Liu}
\author[2]{Boning Meng}

\affil[1]{Institute of Software, Chinese Academy of Sciences, Beijing, China}
\affil[2]{University of Regensburg, Regensburg, Germany}
\affil[ ]{\texttt{liuch.russell@gmail.com}, \texttt{mengboning2013@gmail.com}}
\date{}

\maketitle

\begin{abstract}
In a landmark JACM paper recognized with the 2021 G{\"o}del Prize, Cai and
Chen established a complete complexity dichotomy for counting CSPs over
arbitrary finite domains with algebraic complex weights.  Its polynomial-time
side is characterized by three conditions---Block Orthogonality, Type
Partition, and preservation by a common Mal'tsev operation---quantified over
the countably infinite family $W_{\mathcal{F}}$ generated from arbitrary
$\#\mathrm{CSP}(\mathcal{F})$ instances by partial summation.  They asked
whether these infinitary conditions are decidable from the finite language
$\mathcal{F}$ alone---equivalently, whether the polynomial-time side of this
complete fixed-language classification is uniformly recognizable.  We settle
this problem by giving, for every nonempty finite domain $D$ and every finite
exactly encoded algebraic-complex language $\mathcal{F}$, a total exact
algorithm that decides all three conditions on the full unbounded family
$W_{\mathcal{F}}$.  Beyond decidability, we prove that Block Orthogonality
alone forces both Type Partition and the existence of a single Mal'tsev
operation preserving all generated support and row-equivalence relations.
Thus the three-condition characterization collapses to Block Orthogonality,
and the finite input $(D,\mathcal{F})$ determines which side of the dichotomy
applies.  The same framework decides the corresponding conditions in the
dichotomy theorem for degree-multiple counting CSP proved by Lin.
Technically, any violation at external arity greater than $|D|^2+1$ can be
compressed, by identifying variables, to a violation of external arity at
most $|D|^2+1$.  At each bounded arity, testing all generated tables reduces
to finitely many homogeneous polynomial identities.  We construct a
tensor-character oracle that decides these identities without enumerating
$W_{\mathcal{F}}$.  Scalar tags resolve collisions introduced by
tensorization.  Young's pinned constraint-function isomorphism theorem then
reduces equality of induced characters to a finite isomorphism test.  Linear
independence of distinct monoid characters reduces universal vanishing to
finitely many exact coefficient-sum tests.  For the structural collapse,
power amplification controls complex cancellation, while a generated support
detector realizes every row-equivalence relation as another generated
support; the resulting common Mal'tsev operation yields Type Partition.
\end{abstract}
\section*{Acknowledgments}
The authors used OpenAI's ChatGPT in preparing
this manuscript, including for language editing and \LaTeX{} preparation.
ChatGPT also contributed to the exploratory development of most of the
arguments. All claims and proofs were independently checked and finalized
by the authors, who take full responsibility for the content.

\section{Introduction}
\label{sec:introduction}

Counting constraint satisfaction problems provide a common language for
partition functions, weighted graph homomorphisms, and many other counting
problems.  For nonnegative weights, supports and magnitudes interact
monotonically.  Algebraic complex weights remove this monotonicity: partial
summation can create new zeros, and proportionality or orthogonality can
emerge only after exact cancellation.  These phenomena are central to the
dichotomy theorem of Cai and Chen for complex-weighted
$\#\mathrm{CSP}$~\cite{CaiChen2017}.

Let $D$ be a nonempty finite domain and let $\cF$ be a finite language of
algebraic complex-valued constraint functions.  A finite $\CSP{\cF}$
instance $I$, with variable set $V(I)$, defines a function $F_I$ of all its
variables.  After ordering those variables, let $F_I^{[t]}$ be the function
obtained by summing out every variable after the first $t$.  Following Cai and
Chen~\cite[Section~3]{CaiChen2017}, we associate with $\cF$ the countably
infinite generated family
\begin{equation}
  W_{\cF}
  =\left\{F_I^{[t]}:
    \begin{array}{l}
      I\text{ is a finite }\CSP{\cF}\text{ instance},\\
      1\le t\le \abs{V(I)}
    \end{array}
    \right\}.
  \label{eq:intro-generated-family}
\end{equation}
Variable scopes may repeat variables, and the variables may be reordered
before marginalization.  Thus $W_{\cF}$ contains every generated partial
marginal, permutation, and diagonal minor.  Both the instance size and the
arity in~\eqref{eq:intro-generated-family} are unbounded.

Cai and Chen characterize the polynomial-time side of their dichotomy by
three conditions on this entire family~\cite[Theorem~1 and
Section~3.1]{CaiChen2017}.  The first is a jointly purified
\emph{Block Orthogonality} condition.  The second is \emph{Type Partition}
for the complex row representations of generated functions.  The third asks
for one Mal'tsev operation on $D$ that preserves every generated support and
every generated row-equivalence relation.  Their theorem classifies every
fixed language, and Cai and Chen explicitly asked whether the conjunction of
the three conditions can be recognized uniformly from the finite input
language~\cite[Section~9]{CaiChen2017}.

Using the exact definitions in \cref{sec:preliminaries}, write
$\BO(\cF)$, $\TP(\cF)$, and $\Mal(\cF)$ for these three global predicates,
and put
\[
  \CC(\cF)=\BO(\cF)\wedge\TP(\cF)\wedge\Mal(\cF).
\]
Our main result resolves the uniform recognition problem and proves that the
conjunction of these three conditions is equivalent to Block Orthogonality.

\begin{theorem}[Uniform decidability]
\label{thm:main}
Let $D$ be a nonempty finite domain and let $\cF$ be a finite
constraint language whose values are algebraic complex numbers given by
exact encodings.  There is a deterministic algorithm that halts on every
input $(D,\cF)$ and decides whether $\CC(\cF)$ holds on the full generated
family $W_{\cF}$.  Moreover,
\begin{equation}
  \CC(\cF)\quad\Longleftrightarrow\quad\BO(\cF).
  \label{eq:main-equivalence}
\end{equation}
\end{theorem}

For a prescribed integer $\delta\ge1$, let
$\CSPdeg{\delta}{\cF}$ denote the restriction of $\CSP{\cF}$ to instances in
which every variable has occurrence degree divisible by $\delta$, and let
$\WFdeg{\delta}$ be the corresponding generated family.  Lin proved a
dichotomy theorem for $\CSPdeg{\delta}{\cF}$ whose tractable side is
characterized by the corresponding Block Orthogonality, Type Partition, and
common Mal'tsev conditions on $\WFdeg{\delta}$~\cite[Theorem~3]{Lin2021}.
Our decision procedure and structural-collapse argument apply to this
generated family as well.

\begin{corollary}[Uniform recognition for $\#\mathrm{CSP}^{\delta}$]
\label{cor:degree-multiple-main}
Let $\delta\ge1$.  On input a nonempty finite domain $D$ and a finite
algebraic-complex constraint language $\cF$ whose values are given by exact
encodings, there is a deterministic algorithm deciding whether the three
tractability conditions hold on the full generated family
$\WFdeg{\delta}$.  Writing these conditions as
$\BO^\delta(\cF)$, $\TP^\delta(\cF)$, and $\Mal^\delta(\cF)$, we have
\[
  \BO^\delta(\cF)\wedge\TP^\delta(\cF)\wedge\Mal^\delta(\cF)
  \quad\Longleftrightarrow\quad
  \BO^\delta(\cF).
\]
Consequently, one can decide which side of the dichotomy theorem for
$\CSPdeg{\delta}{\cF}$ proved by Lin~\cite[Theorem~3]{Lin2021} applies.
\end{corollary}

The proof of \cref{cor:degree-multiple-main} is given after the required
degree-filtered identity oracle in
\cref{subsec:degree-multiple-identities,sec:type-collapse}.

\subsection{Proof sketch}
\label{subsec:proof-sketch}

The proof has an algorithmic part, which decides global Block Orthogonality,
and a structural part, which derives Type Partition and the common Mal'tsev
condition from Block Orthogonality.  We indicate below where each step is
proved.
All computations below take place in explicitly represented algebraic number
fields; the required field constructions, conjugation, equality tests, and
root-of-unity computations are standard~\cite{Cohen1993}.  The
Skolem--Mahler--Lech theorem~\cite{Lech1953} establishes termination of two
explicit exact searches by controlling the zero sets of the associated
linear recurrences.

\paragraph{Purification and external-arity compression.}
The reduction to finitely many external arities is carried out in
\cref{sec:purification-compression}.  First,
\cref{lem:purification-invariance} proves that Block Orthogonality of a
purified table is independent of the finite ambient family in which it is
purified; \cref{cor:singleton-purification} therefore reduces the original
joint-purification condition to singleton purification.  A violation is
witnessed by two rows.  The diagonal-minor construction in
\cref{lem:diagonal-minors} identifies row coordinates carrying the same
ordered pair of values on the two witnessing rows and preserves those rows
entry by entry.  Consequently, \cref{lem:bo-arity-bound} shows that every
violation whose external arity exceeds $d^2+1$ can be compressed to a
violation in $W_{\cF}$ of external arity at most $d^2+1$, where
$d=\abs D$.  This bounds the external arity of the compressed
counterexample, not that of the original violating table or the size of an
instance presenting it.

\paragraph{Finite algebraic tests at a fixed arity.}
In \cref{sec:bo-certificates},
\cref{thm:finite-bo-certificates} describes the singleton-purified
block-orthogonal tables as an effectively enumerable finite union of
algebraic sets; the certificate branches explicitly allow entries and
support rectangles to disappear or purified magnitude blocks to merge.
By \cref{cor:bo-product-identities}, universal membership of the generated
$k$-ary tables in this union is equivalent to finitely many polynomial
identities.

The identity oracle is developed in \cref{sec:universal-identities}.
By \cref{lem:pinned-realization}, generated $k$-ary tables are exactly the
pinned partition vectors of finite $k$-labelled instances.
\Cref{lem:tensor-character} realizes every monomial in a polynomial identity
as a multiplicative character of the labelled-instance monoid.  The pinned
form of Young's constraint-function isomorphism theorem, stated as
\cref{thm:young-pinned}, compares the corresponding tensor
targets~\cite[Corollary~35]{Young2025}.  Because
tensorization may identify distinct
input functions, \cref{lem:scalar-tags} adds common positive rational tags;
\cref{lem:character-comparison} then decides equality of the monomial
characters.  Their linear independence is supplied by
\cref{lem:dedekind-independence}, yielding the universal identity oracle in
\cref{thm:universal-identity}.  Combining this oracle with the arity bound
gives \cref{proc:global-bo}, whose termination and correctness on the full
unbounded family $W_{\cF}$ are proved in
\cref{thm:global-bo-decidable}.

\paragraph{From generated supports to one Mal'tsev operation.}
The structural argument begins in \cref{sec:generated-supports}.
\Cref{lem:simultaneous-power-sums,lem:equality-free-support-realization}
use power amplification to realize every equality-free primitive-positive
combination of finitely many generated supports as the support of a
member of $W_{\cF}$, despite possible cancellation.  Assuming $\BO(\cF)$,
the rectangularity argument in \cref{lem:finite-support-maltsev} produces a
common Mal'tsev operation for each finite family of generated supports.
Since the set of operations $D^3\to D$ is finite,
\cref{thm:common-support-maltsev} extracts one operation preserving every
generated support simultaneously.

\paragraph{Realizing row equivalence and obtaining Type Partition.}
Continue under the assumption $\BO(\cF)$.  The remaining structural step
occupies
\crefrange{sec:omega-realization}{sec:type-collapse}.  For a generated table
$G$, \cref{lem:generated-entrywise-powers} realizes the required entrywise
powers inside $W_{\cF}$, and \cref{lem:finite-row-phases} supplies uniform
finite phase control.  Let $K$ be the number field generated by the input
values and their conjugates, let $\mu(K)$ be its finite group of roots of
unity, let $e_K$ be the exponent of $\mu(K)$, and put
\[
  K_\mu=\lcm(2,e_K),
  \qquad L_D=\lcm(1,\ldots,d).
\]
For each fixed $G$, choose a single suitable $t\ge0$ and set
$q=1+tL_D$.  These ingredients give the generated detector
\begin{equation*}
  C_q(\mathbf x,\mathbf y)
  =\sum_{z\in D}
      G(\mathbf x,z)^{(K_\mu-1)q}G(\mathbf y,z)^q.
\end{equation*}
Block Orthogonality keeps the detector zero on every pair outside
$\RowEq G$, while \cref{lem:nondegenerate-exponential-sums} guarantees that
one allowed exponent $q$ can be chosen so that the detector is nonzero on
every dependent pair of nonzero rows.  Thus
\cref{thm:omega-support-realization} proves
$\SuppRel{C_q}=\RowEq G$, where $\SuppRel H$ denotes the support relation of
$H$ and $\RowEq G$ the nonzero-row dependence relation of $G$.
\Cref{cor:omega-reduces-to-supports} then transfers preservation by the common
operation from generated supports to every generated row-equivalence
relation.  Finally, \cref{lem:omega-implies-type} derives Type Partition by
a coordinatewise Mal'tsev calculation, and
\cref{thm:structural-collapse} concludes
$\CC(\cF)\Longleftrightarrow\BO(\cF)$.

\paragraph{The degree-multiple generated family.}
The cyclic degree filter is proved in
\cref{subsec:degree-multiple-identities}, culminating in
\cref{lem:degree-filtered-character-comparison,thm:degree-multiple-universal-identity};
the proof of \cref{cor:degree-multiple-main} then repeats the structural
construction within $\WFdeg{\delta}$.  The pinned realization, support
amplification, entrywise powers, and row detector are finite-instance
constructions; see
\cref{lem:pinned-realization,lem:equality-free-support-realization,lem:generated-entrywise-powers,thm:omega-support-realization}.
These constructions remain in $\WFdeg{\delta}$: fresh variables retain
occurrence degrees divisible by $\delta$, while identifying boundary
variables replaces their degrees by sums of multiples of $\delta$.
Marginalization changes only which variables are retained, not their
occurrence degrees in the underlying instance.

\subsection{Related work}

\paragraph{Counting graph homomorphisms ($\#\mathrm{GH}$).}
For a symmetric $q\times q$ matrix $A$, write $[q]=\{1,\ldots,q\}$.  On an
input graph $G=(V,E)$, the problem $\#\mathrm{GH}(A)$ evaluates
\[
  Z_A(G)=\sum_{\sigma:V(G)\to[q]}
    \prod_{\{u,v\}\in E(G)}A_{\sigma(u),\sigma(v)}.
\]
The problem $\#\mathrm{GH}(A)$ is the special case of $\#\mathrm{CSP}$ with
one symmetric binary constraint and is therefore covered by the
complex-weighted $\#\mathrm{CSP}$ dichotomy theorem of Cai and
Chen~\cite{CaiChen2017}.  Dyer and Greenhill classified $\{0,1\}$-valued
matrices, Bulatov and Grohe
extended the dichotomy to nonnegative matrices, Goldberg, Grohe, Jerrum, and
Thurley handled arbitrary real matrices, and Cai, Chen, and Lu completed the
classification for algebraic complex matrices
\cite{DyerGreenhill2000,DyerGreenhill2004,BulatovGrohe2005,
GoldbergGroheJerrumThurley2010,CaiChenLu2013}.  For algebraic-complex
matrices, the dichotomy theorem of Cai, Chen, and Lu also provides a
polynomial-time recognition algorithm for its tractable
case~\cite[Theorem~1.2]{CaiChenLu2013}.

Restricting the input $G$ to be planar creates additional matchgate-tractable
cases.  Cai and Maran proved a complete dichotomy for every real symmetric
$3\times3$ matrix, with the genuinely planar tractable cases exactly those
obtained by holographic algorithms using matchgates~\cite{CaiMaran2023}.
They subsequently obtained the analogous classification for nonnegative
full-rank $4\times4$ matrices, where tensor products of matchgates give new
tractable families~\cite{CaiMaran2024}.  In a 2026 preprint, Cai, Maran, and
Young proved a dichotomy for nonnegative matrices with pairwise distinct
diagonal entries, and more generally whenever planar edge gadgets separate
every pair of target vertices~\cite{CaiMaranYoung2026}.  Their characterization
of this separation through quantum automorphism groups yields a sharp barrier
to the usual interpolation approach, and a complete dichotomy for arbitrary
nonnegative symmetric target matrices remains open~\cite{CaiMaranYoung2026}.

\paragraph{Counting CSP.}
For general constraint languages, Dyer, Goldberg, and Jerrum classified
nonnegative weighted Boolean $\#\mathrm{CSP}$
\cite{DyerGoldbergJerrum2009}; Bulatov proved the finite-domain unweighted
dichotomy, and Dyer and Richerby subsequently gave an elementary effective
criterion
\cite{Bulatov2013,DyerRicherby2013}.  For algebraic weights, Cai, Chen, and
Lu obtained an effective dichotomy in the nonnegative setting; Cai, Lu, and
Xia classified complex-weighted Boolean languages, and Cai and Chen
established the fixed-language dichotomy over arbitrary finite domains with
algebraic complex weights
\cite{CaiChenLu2016,CaiLuXia2014,CaiChen2017}.  Lin proved a dichotomy
theorem for the degree-divisibility restriction $\#\mathrm{CSP}^{\delta}$
whose tractable side is characterized by the corresponding three infinitary
conditions on $\WFdeg{\delta}$~\cite[Theorem~3]{Lin2021}.  That work left
their uniform recognition open~\cite[p.~40:3]{Lin2021};
\cref{cor:degree-multiple-main} resolves this question as well.  Existing
effective criteria either exclude complex cancellation,
restrict the domain to Boolean, or restrict the language to a single symmetric
binary constraint; none provides, from an arbitrary finite-domain
algebraic-complex input, a uniform decision procedure for the three
infinitary tractability conditions in the complex-weighted $\#\mathrm{CSP}$
dichotomy theorem of Cai and Chen~\cite[Section~3.1]{CaiChen2017}.  This is
the recognition gap closed here.  Young's constraint-function isomorphism
result for labelled instances~\cite[Corollary~35]{Young2025} supplies our
exact comparison of tensor characters.

\paragraph{Holant.}
Valiant's holographic algorithms introduced basis transformations for
tensor-network partition functions~\cite{Valiant2008}, and Cai, Lu, and Xia
formulated Holant as a general counting framework containing $\#\mathrm{CSP}$
and graph homomorphisms~\cite{CaiLuXia2009}.  On the Boolean domain, a
sequence of dichotomies progressed from symmetric and auxiliary-signature
settings to arbitrary nonnegative signatures, complex-valued
$\operatorname{Holant}^{c}$, and finally
arbitrary real-valued signatures
\cite{CaiGuoWilliams2016,LinWang2018,
Backens2021,Backens2025,ShaoCai2020}.  The full complex-valued Boolean
classification remains open~\cite{MengWangXiaZheng2025,GuanShaoShi2026}.
Meng, Wang, Xia, and Zheng obtained an
$\mathsf{FP}^{\mathsf{NP}}$-versus-$\#\mathsf P$-hard classification in the
presence of a nontrivial odd-arity signature~\cite{MengWangXiaZheng2025}.
Guan, Shao, and Shi subsequently gave polynomial-time algorithms for every
case on the $\mathsf{FP}^{\mathsf{NP}}$ side, yielding a genuine
$\mathsf{FP}$-versus-$\#\mathsf P$-hard dichotomy in this
regime~\cite{GuanShaoShi2026}.  Hence any remaining open case may be taken to
have only even-arity nonzero signatures.

Beyond the Boolean domain, Cai, Lu, and Xia classified
$\operatorname{Holant}^{*}(f)$ for a single complex-valued symmetric ternary
signature on domain size three~\cite{CaiLuXia2013}, while Cai, Guo, and
Williams obtained dichotomies for counting edge colorings and certain
associated higher-domain Holant problems~\cite{CaiGuoWilliams2014}.  Liu,
Fan, and Cai later obtained
single-signature dichotomies for real-valued symmetric ternary signatures on
domain size three and $\{0,1\}$-valued symmetric ternary signatures on domain
size four~\cite{LiuFanCai2025}.  Cai and Ihm proved a
dichotomy for $\operatorname{Holant}^{*}_{3}(\mathcal F)$, where all unary
signatures are freely available and $\mathcal F$ is an arbitrary set of
symmetric real-valued signatures~\cite{CaiIhm2025}.
Despite these advances, a general finite-domain Holant dichotomy remains
unknown~\cite{CaiIhm2025}.

\paragraph{Decision CSP.}
For the decision problem, Schaefer established the Boolean dichotomy
\cite{Schaefer1978}.  Feder and Vardi placed finite-domain CSP in a broad
logical framework and formulated the dichotomy conjecture that every fixed
finite constraint language yields either a polynomial-time problem or an
$\mathsf{NP}$-complete one~\cite{FederVardi1998}.  The universal-algebraic
approach relates CSP complexity to polymorphisms
\cite{BulatovJeavonsKrokhin2005}; Mal'tsev polymorphisms also underlie
effective counting-CSP algorithms~\cite{DyerRicherby2013}.  Bulatov and Zhuk
independently proved the full finite-domain dichotomy in
2017~\cite{Bulatov2017,Zhuk2017}; see also Zhuk's
full journal account~\cite{Zhuk2020}.  Those decision results concern
supports alone; the present recognition problem must additionally control
magnitudes, phases, and zeros created by complex cancellation.




\subsection{Organization}

\Cref{sec:preliminaries} formalizes the three tractability conditions in the
complex-weighted $\#\mathrm{CSP}$ dichotomy theorem of Cai and
Chen~\cite[Section~3.1]{CaiChen2017} and specifies the exact input model,
including the group-theoretic notation used later.
\Crefrange{sec:purification-compression}{sec:bo-algorithm} develops the
decision procedure for global Block Orthogonality; the degree filter for
$\#\mathrm{CSP}^{\delta}$ appears in
\cref{subsec:degree-multiple-identities}.
\Crefrange{sec:generated-supports}{sec:type-collapse} proves the structural
implication in~\eqref{eq:main-equivalence}.  \Cref{sec:conclusion} summarizes
the results.




\section{Preliminaries and the Three Tractability Conditions}
\label{sec:preliminaries}

This section formalizes the three tractability conditions in the
complex-weighted $\#\mathrm{CSP}$ dichotomy theorem of Cai and
Chen~\cite[Sections~2.8--2.9 and~3.1]{CaiChen2017}.

\subsection{Algebraic, group-theoretic, and character notation}

For an integer $n\ge 1$, let $[n]=\{1,\ldots,n\}$; set
$[0]=\varnothing$ and $D^0=\{\emptytuple\}$, where $\emptytuple$ is the empty
tuple.  Throughout, $D$ is a nonempty finite domain and $d=\abs D$.
We write $\AlgNums\subseteq\mathbb C$ for the field of algebraic complex
numbers.  Each input value has an exact encoding that includes its selected
embedding into $\mathbb C$, for example a square-free polynomial over
$\mathbb Z$ and a rational isolating rectangle.  Standard number-field
algorithms can therefore perform arithmetic, conjugation, equality tests,
and zero tests exactly~\cite{Cohen1993}.

We collect the elementary group notation used below.  If $L$ is a field,
then
\[
  L^\times=L\setminus\{0\}
\]
is its multiplicative group.  For $S\subseteq L^\times$, the subgroup
generated by $S$ is
\[
  \begin{aligned}
  \langle S\rangle
  =\biggl\{&\prod_{a\in T}a^{n_a}:
      T\subseteq S\text{ is finite},\\
      &n_a\in\mathbb Z\text{ for every }a\in T\biggr\}.
  \end{aligned}
\]
We write $H\le\Gamma$ when $H$ is a subgroup of the abelian group
$\Gamma$.  A group homomorphism from a multiplicative group to an additive
group satisfies $e(xy)=e(x)+e(y)$; its kernel is
$\ker e=\{x:e(x)=0\}$.  The additive group $\mathbb Z^m$ is torsion-free:
$N\mathbf v=0$ with $N\ge1$ implies $\mathbf v=0$.

An element $x$ of a group is \emph{torsion} if $x^N=1$ for some $N\ge1$.
For an abelian group $\Gamma$, its torsion elements form the subgroup
\[
  \Tor(\Gamma)=\{x\in\Gamma:x^N=1\text{ for some }N\ge1\}.
\]
If $\psi:\Gamma\to\mathbb C^\times$ is a homomorphism, we say that
$\psi$ \emph{fixes torsion pointwise} if
\begin{equation}
  \psi(\xi)=\xi
  \qquad\text{for every }\xi\in\Tor(\Gamma).
  \label{eq:fixes-torsion-definition}
\end{equation}
This differs from two other phrases used below.  An additive homomorphism
$e:\Gamma\to\mathbb Z^m$ \emph{annihilates torsion} when
$e(\xi)=0$ for every $\xi\in\Tor(\Gamma)$, whereas an integer $N$
\emph{kills torsion} in a finite subgroup $T$ when $\xi^N=1$ for every
$\xi\in T$.
Let $\mu_\infty\le\AlgNums^\times$ be the group of all algebraic roots of
unity.  Thus, for $\Gamma\le\AlgNums^\times$,
\begin{equation*}
  \Tor(\Gamma)=\Gamma\cap\mu_\infty.
\end{equation*}
If $K$ is a number field, then
$\mu(K)=K^\times\cap\mu_\infty$ is a finite cyclic group.  The
\emph{exponent} of a finite group is the least positive integer $e$ such
that $x^e=1$ for every group element $x$; equivalently, it is the least
common multiple of the element orders.  We write $\lcm$ for least common
multiple.

Because every group here is abelian, $\Tor(\Gamma)$ is normal and the
quotient $\Gamma/\Tor(\Gamma)$ is defined.  Its elements are cosets
$[x]=x\Tor(\Gamma)$; two elements determine the same coset exactly when
their quotient is torsion.  If a homomorphism
$e:\Gamma\to\mathbb Z^m$ has kernel $\Tor(\Gamma)$, it induces the
injective homomorphism
\[
  \Gamma/\Tor(\Gamma)\hookrightarrow\mathbb Z^m,
  \qquad [x]\longmapsto e(x).
\]

A \emph{monoid} is a set with an associative multiplication and an identity
element.  All monoids below are commutative; unless displayed explicitly in
additive notation, they are written multiplicatively.
A unital multiplicative \emph{character} from a monoid $M$ to a field $L$ is
a map
\[
  \chi:M\to L,
  \qquad
  \chi(1)=1,
  \qquad
  \chi(xy)=\chi(x)\chi(y).
\]
Thus a character is a numerical observable that turns multiplication in the
monoid into multiplication in the field.  Two characters are \emph{equal}
when they agree on every element of $M$; syntactically different formulas
may therefore define the same character.  Pointwise products of characters
are again characters, and the complex conjugate of a complex-valued
character is again a character.  A character on a monoid may take the value
zero.  For example, on the additive monoid $\mathbb Z_{\ge0}$, the map that sends
$0$ to $1$ and every positive integer to $0$ is a character.  By contrast,
a character on a group cannot vanish, because
$\chi(x)\chi(x^{-1})=\chi(1)=1$.

As a basic example, let $M=(\mathbb Z_{\ge0}^r,+,\mathbf 0)$ and choose
$\boldsymbol\lambda=(\lambda_1,\ldots,\lambda_r)\in(L^\times)^r$.  Then
\begin{equation*}
  \chi_{\boldsymbol\lambda}(\mathbf n)
  =\prod_{j=1}^r\lambda_j^{n_j}
\end{equation*}
is a character: addition of exponent vectors multiplies the displayed
values.  The coordinate functions of the pinned partition vectors introduced
in \cref{subsec:labelled-instances} may vanish, but they satisfy the same
multiplicativity law and hence are characters.  Their complex conjugates and
finite pointwise products are characters as well.

We will also use the following elementary principle, stated formally and
proved in \cref{lem:dedekind-independence}: distinct unital multiplicative
characters from a monoid to a characteristic-zero field are linearly
independent as functions on that monoid.  Hence, for a finite expression
\[
  \sum_i c_i\chi_i,
\]
one first partitions the $\chi_i$ into classes of equal characters.  The
expression vanishes on every monoid element if and only if, in each class,
the sum of its coefficients $c_i$ is zero.  This distinction between a
formal expression and the function it induces is the basis of the universal
identity test in \cref{thm:universal-identity}.

Finally, for an integer $\delta\ge1$, write
$C_\delta=\mathbb Z/\delta\mathbb Z$ and choose a primitive
$\delta$th root of unity $\omega_\delta$ (with $\omega_1=1$).  For fixed
$n\in\mathbb Z$, the map $s\mapsto\omega_\delta^{sn}$ is a character of the
additive group $C_\delta$.  It is the constant-one character exactly when
$\delta\mid n$; otherwise its values form a nontrivial geometric progression
whose sum is zero.  Thus
\begin{equation}
  \sum_{s\in C_\delta}\omega_\delta^{sn}
  =
  \begin{cases}
    \delta,&\delta\mid n,\\
    0,&\delta\nmid n.
  \end{cases}
  \label{eq:cyclic-character-filter}
\end{equation}

\subsection{Constraint instances and generated tables}

A weighted constraint language $\cF$ over $D$ is a finite set of typed
functions $f:D^{r_f}\to\AlgNums$ with $r_f\ge1$.  If the input
is presented as a list, we identify exact duplicate typed tables; the
generated family is unchanged.  An instance $I$ of
$\CSP{\cF}$ consists of a variable set
$V(I)=\{x_1,\ldots,x_n\}$, where $n(I)=\abs{V(I)}=n\ge1$, and a finite
multiset $\mathcal C_I$ of constraints $(f;i_1,\ldots,i_{r_f})$, where
$f\in\cF$ and $i_1,\ldots,i_{r_f}\in[n]$.  A scope is ordered and its
indices need not be distinct.  The instance function $F_I:D^n\to\AlgNums$
is
\begin{equation}
  F_I(a_1,\ldots,a_n)
  =\prod_{(f;i_1,\ldots,i_{r_f})\in\mathcal C_I}
      f(a_{i_1},\ldots,a_{i_{r_f}}),
  \label{eq:instance-function}
\end{equation}
where multiplicities are retained and the empty product equals $1$.
For $t\in[n]$, define
\begin{equation}
  F_I^{[t]}(a_1,\ldots,a_t)
  =\sum_{(a_{t+1},\ldots,a_n)\in D^{n-t}}
      F_I(a_1,\ldots,a_n).
  \label{eq:partial-marginal}
\end{equation}
Following Cai and Chen~\cite{CaiChen2017}, the generated family is
\begin{equation*}
  W_{\cF}
  =\left\{F_I^{[t]}:
    \begin{array}{l}
      I\text{ is a finite }\CSP{\cF}\text{ instance},\\
      t\in[n(I)]
    \end{array}
    \right\}.
\end{equation*}
The variable ordering in an instance is arbitrary, so
\eqref{eq:partial-marginal} retains any prescribed ordered subset after
relabeling.  Repeated scope indices permit diagonal minors and arbitrary
identifications.  Instances may contain arbitrarily many variables,
constraints, hidden variables, and retained variables.

For $\delta\ge1$, the \emph{occurrence degree} of a variable is the number
of positions in all constraint scopes occupied by that variable, counting
constraint multiplicity and repeated positions.  Following
Lin~\cite{Lin2021}, $\CSPdeg{\delta}{\cF}$ denotes the restriction of
$\CSP{\cF}$ to instances in which every variable has occurrence degree
divisible by $\delta$.  Its
generated family is
\begin{equation*}
  \WFdeg{\delta}
  =\left\{F_I^{[t]}:
    \begin{array}{l}
      I\text{ is a finite }\CSPdeg{\delta}{\cF}\text{ instance},\\
      t\in[n(I)]
    \end{array}
    \right\}.
\end{equation*}
For $\delta=1$ this is $W_{\cF}$.  Isolated variables have degree zero and
are permitted.  We allow $\cF$ to be empty; its degenerate case is handled
after \cref{def:type-maltsev}, once all three conditions and the relations
they use have been defined.

\subsection{Rows and Block Orthogonality}

For every table $G:D^r\to\AlgNums$, write $\arity(G)=r$.  Following Cai and
Chen~\cite[Sections~2.3 and~2.5]{CaiChen2017}, suppose $r\ge2$ and regard
$G$ as the
$d^{r-1}\times d$ matrix whose row indexed by
$\mathbf x\in D^{r-1}$ is
\[
  G(\mathbf x,\ast)=\bigl(G(\mathbf x,z)\bigr)_{z\in D}.
\]
The nonzero row indices carry the equivalence relation
\begin{equation}
  \mathbf x\sim_G\mathbf y
  \quad\Longleftrightarrow\quad
  \begin{gathered}
    G(\mathbf x,\ast)\text{ and }G(\mathbf y,\ast)\text{ are}\\
    \text{linearly dependent over }\mathbb C.
  \end{gathered}
  \label{eq:complex-row-equivalence}
\end{equation}
Let $S_1,\ldots,S_s$ be its equivalence classes.  For each $j\in[s]$,
let $\mathbf v_j$ be the unique vector proportional to the rows in $S_j$
whose first nonzero entry, with respect to a fixed ordering of $D$, is $1$.
The family
\begin{equation*}
  \Row(G)=\{(S_1,\mathbf v_1),\ldots,(S_s,\mathbf v_s)\}
\end{equation*}
is the \emph{row representation} of $G$.  Distinct representatives are
pairwise nonproportional, although they may be linearly dependent as a
family.  If
$G$ is the zero table, then $s=0$ and $\Row(G)=\varnothing$.

For a table or vector $H$, write $\abs H$ for its entrywise absolute value.
For a vector $\mathbf u\in\mathbb C^D$, write
$\supp(\mathbf u)=\{z\in D:u_z\ne0\}$.  The complex
relation~\eqref{eq:complex-row-equivalence} must be distinguished from the
magnitude-row relation.  The function $G$ is \emph{block-rank one} if the
nonnegative matrix $\abs G$ is block-rank one: for every two nonzero rows,
their magnitude vectors are positively proportional or their supports are
disjoint.  Equivalently, the nonzero support is a disjoint union of complete
row-by-column rectangles, and within each rectangle the magnitude rows are
proportional.

Let $\mathbf u,\mathbf v\in\mathbb C^D$ be nonzero and suppose that
$\abs{\mathbf u}$ and $\abs{\mathbf v}$ are proportional.  Their supports
coincide; call the common support $T$.  Let
$\mu_1>\cdots>\mu_h>0$ be the distinct values of $\abs{u_z}$ on $T$, and
put
\[
  T_i=\{z\in T:\abs{u_z}=\mu_i\},\qquad i\in[h].
\]
The same partition is obtained from $\mathbf v$.  The two vectors are
\emph{block-orthogonal}, written
$\mathbf u\perp_{\mathrm B}\mathbf v$, if
\begin{equation}
  \sum_{z\in T_i}u_z\overline{v_z}=0
  \qquad\text{for every }i\in[h].
  \label{eq:vector-block-orthogonality}
\end{equation}

\begin{definition}[Block-orthogonal function]
\label{def:block-orthogonal-function}
A function $G:D^r\to\AlgNums$, $r\ge2$, is \emph{block-orthogonal} if it
is block-rank one and, whenever two nonzero complex rows have proportional
magnitude vectors, those rows are either linearly dependent or
block-orthogonal in the sense of~\eqref{eq:vector-block-orthogonality}.
\end{definition}

\paragraph{A generated block-rank-one example.}
Let $D=\{0,1,2,3\}$ and let the binary function $A$ be
\[
\begin{array}{c|rrrr}
  x\backslash z&0&1&2&3\\ \hline
  0&1&2&0&0\\
  1&3&6&0&0\\
  2&0&0&1&4\\
  3&0&0&2&8
\end{array}.
\]
For the language $\cF_A=\{A\}$, the instance consisting of the single
constraint $A(x,z)$ has $F_I^{[2]}=A$, so $A\in W_{\cF_A}$.  The first two
magnitude rows are proportional, as are the last two, while rows from
different pairs have disjoint supports.  Equivalently, the support is the
disjoint union of the rectangles
\[
  \{0,1\}\times\{0,1\}
  \qquad\text{and}\qquad
  \{2,3\}\times\{2,3\}.
\]
Thus $A$ is block-rank one.  In this example every pair of
magnitude-proportional complex rows is already linearly dependent, which
establishes Block Orthogonality directly.

\paragraph{A generated block-orthogonality example.}
On the same domain, let
\[
B=
\begin{array}{c|rrrr}
  x\backslash z&0&1&2&3\\ \hline
  0&2&2&1&1\\
  1&2&-2&1&-1\\
  2&0&0&0&0\\
  3&0&0&0&0
\end{array}.
\]
Again $B\in W_{\{B\}}$ by the one-constraint instance.  Its two nonzero
rows
\[
  \mathbf u=(2,2,1,1),\qquad
  \mathbf v=(2,-2,1,-1)
\]
are not linearly dependent, but
$\abs{\mathbf u}=\abs{\mathbf v}=(2,2,1,1)$.  Their equal-magnitude blocks
are $T_1=\{0,1\}$ and $T_2=\{2,3\}$, and cancellation holds separately on
both blocks:
\[
  \sum_{z\in T_1}u_z\overline{v_z}=4-4=0,
  \qquad
  \sum_{z\in T_2}u_z\overline{v_z}=1-1=0.
\]
Thus $\mathbf u\perp_{\mathrm B}\mathbf v$ and $B$ is block-orthogonal.
The displayed table is already pure, and the legal one-element generating
set $g_1=2$, assigned to $p_1=2$, leaves it unchanged under singleton
purification.

\subsection{Joint purification}

A function is
\emph{pure} if each of its values is a nonnegative integer times a root of
unity.  We use the notions of pure functions and generating sets, together
with the Purification Lemma, from Cai and
Chen~\cite[Definitions~2--3 and Lemma~7]{CaiChen2017}.  Given a finite tuple
$\mathbf G=(G_1,\ldots,G_h)$ with $h\ge1$, let $A(\mathbf G)$ be the set of
its nonzero values.  A \emph{generating set} for $A(\mathbf G)$ is a tuple
$g_1,\ldots,g_m$, where $m\ge0$, satisfying the following three conditions:
\begin{enumerate}
  \item $g_i\in\mathbb Q(A(\mathbf G))^\times$ for every $i\in[m]$;
  \item for every nonzero
        $\mathbf k=(k_1,\ldots,k_m)\in\mathbb Z^m$,
        \[
          \prod_{i=1}^m g_i^{k_i}\notin\mu_\infty;
        \]
  \item every $a\in A(\mathbf G)$ has a unique exponent vector
        $e(a)=(e_1(a),\ldots,e_m(a))\in\mathbb Z^m$ such that
\begin{equation}
  a=\vartheta(a)\prod_{i=1}^m g_i^{e_i(a)},
  \qquad \vartheta(a)\in\mu_\infty.
  \label{eq:multiplicative-decomposition}
\end{equation}
\end{enumerate}
Here $\mathbb Q(A(\mathbf G))$ is the smallest subfield of $\AlgNums$
containing $\mathbb Q$ and all elements of $A(\mathbf G)$, and
$\mathbb Q(A(\mathbf G))^\times
=\mathbb Q(A(\mathbf G))\setminus\{0\}$ is its multiplicative group.
When $A(\mathbf G)$ is nonempty, such a generating set
exists~\cite[Lemma~7.2]{CaiChenLu2013}; see
also~\cite[Lemma~8]{CaiChen2017}.  Let $p_1,\ldots,p_m$ be the first $m$
positive primes.
For every nonzero value $a$, put
\[
  q(a)=\prod_{i=1}^m p_i^{e_i(a)}\in\mathbb Q_{>0}.
\]
The transformation displayed in the proof of the Purification
Lemma~\cite[Lemma~7]{CaiChen2017} replaces each nonzero $a$ by
$\vartheta(a)q(a)$; zeros remain zero.  Write $P_i^0$ for the resulting table
corresponding to $G_i$.
Let $N_i$ be the least positive integer such that
$N_iq(a)\in\mathbb Z_{>0}$ for every nonzero value $a$ of $G_i$, taking
$N_i=1$ when $G_i$ is zero, and put $P_i=N_iP_i^0$.  The latter table-wide
normalization clears the denominators introduced by negative exponents and
makes $P_i$ pure.  Positive table-wide scaling
changes neither support nor any clause of Block Orthogonality, so the printed
and normalized forms define the same global condition.

Fix one legal generating-set choice whenever $A(\mathbf G)$ is nonempty and
define
\[
  \Pure(\mathbf G)=(P_1,\ldots,P_h),
\]
using the pure, table-wise normalized outputs.  If every nonzero value is a
root of unity, the empty generating list is legal, and the fixed mapping uses
it.  If $A(\mathbf G)=\varnothing$, define $\Pure(\mathbf G)$ to be the
all-zero tuple.  We call any printed or normalized output obtained from a
legal generating-set choice a \emph{legal purification}.
\Cref{lem:purification-invariance} proves that the relevant structure is
independent of these choices.

\begin{definition}[Global Block Orthogonality]
\label{def:global-bo}
The language $\cF$ satisfies the \emph{Block Orthogonality condition} of Cai
and Chen~\cite[Section~3.1]{CaiChen2017}, denoted $\BO(\cF)$, if for every
nonempty finite tuple
$(G_1,\ldots,G_h)$ of pairwise distinct members of $W_{\cF}$, writing
\[
  \Pure(G_1,\ldots,G_h)=(P_1,\ldots,P_h),
\]
the table $P_i$ is block-orthogonal for every $i$ with
$\arity(G_i)\ge2$.
\end{definition}

The quantifier in \cref{def:global-bo} is joint.
\Cref{cor:singleton-purification} proves its equivalence with the
corresponding singleton condition.

\subsection{Types and derived relations}

Following Cai and Chen~\cite[Sections~2.8--2.9 and~3.1]{CaiChen2017}, let
$G:D^r\to\AlgNums$, $r\ge2$, have row representation
$\{(S_j,\mathbf v_j):j\in[s]\}$.  For $\ell\in[r-1]$ and
$\boldsymbol\alpha\in D^\ell$, define
\begin{equation*}
  \typeop_G(\boldsymbol\alpha)
  =\left\{j\in[s]:
       \text{some }\mathbf x\in S_j\text{ has prefix }
       \boldsymbol\alpha\right\}.
\end{equation*}
The map is a \emph{type-partition map} if
\begin{equation*}
  \begin{aligned}
  &\forall \ell\in[r-1]\ \forall
    \boldsymbol\alpha,\boldsymbol\beta\in D^\ell,\\
  &\qquad
    \typeop_G(\boldsymbol\alpha)=\typeop_G(\boldsymbol\beta),\\[-2pt]
  &\qquad\text{or}\qquad
    \typeop_G(\boldsymbol\alpha)
      \cap\typeop_G(\boldsymbol\beta)=\varnothing.
  \end{aligned}
\end{equation*}
Set $\typeop_G(\emptytuple)=[s]$.  If $G$ is the zero table, then $s=0$, so
every type is empty and Type Partition holds.

\paragraph{A generated Type Partition example.}
Regard $D=\{0,1,2\}$ as the additive group $\mathbb Z_3$, and define
$G:D^3\to\AlgNums$ by
\[
\begin{array}{c|rrr}
  (x_1,x_2)\backslash z&0&1&2\\ \hline
  (0,0)&1&0&0\\
  (0,1)&0&1&0\\
  (0,2)&0&0&1\\
  (1,0)&0&2&0\\
  (1,1)&0&0&2\\
  (1,2)&2&0&0\\
  (2,0)&0&0&4\\
  (2,1)&4&0&0\\
  (2,2)&0&4&0
\end{array}.
\]
Equivalently,
\begin{equation}
  G(x_1,x_2,z)\ne0
  \quad\Longleftrightarrow\quad
  z\equiv x_1+x_2\pmod 3,
  \label{eq:type-example-modular-support}
\end{equation}
and the unique nonzero entry in a row has value $2^{x_1}$.
The one-constraint instance gives $G\in W_{\{G\}}$.  Its three nonzero
complex row classes are
\[
\begin{aligned}
  S_1&=\{(0,0),(1,2),(2,1)\},& \mathbf v_1&=(1,0,0),\\
  S_2&=\{(0,1),(1,0),(2,2)\},& \mathbf v_2&=(0,1,0),\\
  S_3&=\{(0,2),(1,1),(2,0)\},& \mathbf v_3&=(0,0,1).
\end{aligned}
\]
For prefixes of length one,
\[
  \typeop_G(0)=\typeop_G(1)=\typeop_G(2)=\{1,2,3\}.
\]
Thus one type contains three row-class labels and occurs at all three
prefixes.  For prefixes of length two, if $a+b\equiv j\pmod3$ with
$j\in\{0,1,2\}$, then
\[
  \typeop_G(a,b)=\{j+1\}.
\]
Hence, at each fixed prefix length, any two types are equal or disjoint.
Thus a type records row-class labels; domain values and individual rows enter
through the labels they realize.  The modular form of this example will also let
the same table illustrate the Mal'tsev condition below.

For every $G:D^r\to\AlgNums$, define its support relation
\begin{equation*}
  \SuppRel G=\{\mathbf a\in D^r:G(\mathbf a)\ne0\}.
\end{equation*}
For $r\ge2$, define the $2(r-1)$-ary row-equivalence relation
\begin{equation*}
  \begin{aligned}
    \RowEq G
    &=\bigcup_{j=1}^s(S_j\times S_j)\\
    &=\left\{(\mathbf x,\mathbf y)\in D^{r-1}\times D^{r-1}:
      \begin{array}{l}
        G(\mathbf x,\ast),G(\mathbf y,\ast)\text{ are nonzero}\\
        \text{and linearly dependent}
      \end{array}
      \right\}.
  \end{aligned}
\end{equation*}
Both types and $\RowEq G$ are defined from the complex row classes of the
original table $G$.

An operation $m:D^3\to D$ \emph{preserves} a relation $R\subseteq D^q$ if
coordinatewise application of $m$ to any three tuples of $R$ returns a tuple
of $R$.  It is a \emph{Mal'tsev operation} if
\begin{equation*}
  m(a,b,b)=m(b,b,a)=a
  \qquad\text{for all }a,b\in D.
\end{equation*}
Set
\begin{equation*}
  \begin{aligned}
  \Lambda_{\cF}
  ={}&\{\SuppRel G:G\in W_{\cF}\}\\
     &{}\cup
       \{\RowEq G:G\in W_{\cF},\ \arity(G)\ge2\}.
  \end{aligned}
\end{equation*}
$\Lambda_{\cF}$ is a family of relations of varying arities, and preservation
is required separately for every $\SuppRel G$ and every applicable
$\RowEq G$.

\begin{definition}[Type Partition and the common Mal'tsev condition]
\label{def:type-maltsev}
The language $\cF$ satisfies $\TP(\cF)$ if every generated table of arity at
least two has a type-partition map.  It satisfies $\Mal(\cF)$ if
\begin{equation*}
  \begin{aligned}
  \exists m:D^3\to D\quad &m\text{ is a Mal'tsev operation},\\
  &\forall R\in\Lambda_{\cF},\quad m\text{ preserves }R.
  \end{aligned}
\end{equation*}
The order of these quantifiers is essential: one operation must preserve the
entire family.
\end{definition}

\paragraph{The Type Partition table as a Mal'tsev-preservation example.}
Continue with the same $D=\mathbb Z_3$ and the same generated table $G$
from the preceding Type Partition example.  Consider the affine operation
\[
  m(a,b,c)=a-b+c\pmod3.
\]
It is Mal'tsev because
\[
  m(a,b,b)=a-b+b=a,
  \qquad
  m(b,b,a)=b-b+a=a.
\]
By~\eqref{eq:type-example-modular-support}, the support relation is the
nontrivial affine relation
\[
  \SuppRel G
  =\{(x_1,x_2,z)\in D^3:z\equiv x_1+x_2\pmod3\}.
\]
Take any three tuples
$\mathbf u^{(i)}=(x_1^{(i)},x_2^{(i)},z^{(i)})\in\SuppRel G$,
$i\in\{1,2,3\}$, and let
$\mathbf u=(u_1,u_2,u_3)$ be their coordinatewise image under $m$.  Then
\[
\begin{aligned}
  u_3
  &\equiv z^{(1)}-z^{(2)}+z^{(3)}\\
  &\equiv
    \bigl(x_1^{(1)}-x_1^{(2)}+x_1^{(3)}\bigr)
    +\bigl(x_2^{(1)}-x_2^{(2)}+x_2^{(3)}\bigr)\\
  &\equiv u_1+u_2\pmod3.
\end{aligned}
\]
Thus $\mathbf u\in\SuppRel G$, so $m$ preserves the support.

The row classes displayed above also give the direct description
\[
  \RowEq G
  =
  \left\{
  \begin{array}{l}
    (x_1,x_2,y_1,y_2)\in D^4:\\
    x_1+x_2\equiv y_1+y_2\pmod3
  \end{array}
  \right\}.
\]
Indeed, two rows are proportional exactly when their unique nonzero entries
occur in the same column.  Take three tuples
\[
  \mathbf r^{(i)}
  =(x_1^{(i)},x_2^{(i)},y_1^{(i)},y_2^{(i)})\in\RowEq G,
  \qquad i\in\{1,2,3\}.
\]
Put
\[
  s_x^{(i)}=x_1^{(i)}+x_2^{(i)},
  \qquad
  s_y^{(i)}=y_1^{(i)}+y_2^{(i)}.
\]
Then $s_x^{(i)}\equiv s_y^{(i)}\pmod3$ for each $i$.  If the coordinatewise
image of the three tuples is $(X_1,X_2,Y_1,Y_2)$, linearity modulo three
gives
\[
\begin{aligned}
  X_1+X_2
  &\equiv s_x^{(1)}-s_x^{(2)}+s_x^{(3)}\\
  &\equiv s_y^{(1)}-s_y^{(2)}+s_y^{(3)}\\
  &\equiv Y_1+Y_2\pmod3.
\end{aligned}
\]
Therefore $(X_1,X_2,Y_1,Y_2)\in\RowEq G$.  The same operation $m$
preserves $\SuppRel G$ and $\RowEq G$ \emph{separately}, so this one table
simultaneously illustrates Type Partition and the single-table content of
the Mal'tsev condition.  The global predicates $\TP(\cF)$ and $\Mal(\cF)$
still quantify over all applicable tables and relations in $W_{\cF}$.

For $\delta\ge1$, define $\BO^\delta(\cF)$,
$\TP^\delta(\cF)$, and $\Mal^\delta(\cF)$ by the same definitions with
$W_{\cF}$ replaced throughout by $\WFdeg{\delta}$, and put
\[
  \CC^\delta(\cF)
  =\BO^\delta(\cF)\wedge\TP^\delta(\cF)\wedge\Mal^\delta(\cF).
\]

In both the ordinary and degree-$\delta$ settings, Type Partition and the
Mal'tsev condition are considered under the corresponding Block Orthogonality
condition.  Types, $\SuppRel G$, and $\RowEq G$ are formed from the original
generated table; purification occurs only in the corresponding Block
Orthogonality condition.  If $G:D^r\to\AlgNums$ is the zero table, then
$\SuppRel G=\varnothing$ and, when $r\ge2$,
$\RowEq G=\varnothing$; every operation preserves these empty relations
vacuously.

We now discharge the empty-language case.  If $\cF=\varnothing$, then
$F_I\equiv1$, and every generated table $G=F_I^{[t]}$ is the positive
constant table
\[
  G(a_1,\ldots,a_t)=d^{n(I)-t}.
\]
Consequently $\SuppRel G=D^t$.  When $t\ge2$, all rows are nonzero and
proportional, there is exactly one complex row class, and
$\RowEq G=D^{2(t-1)}$.  Every legal purification is again a nonzero constant
table, so Block Orthogonality holds; since $G$ itself is constant, Type
Partition holds as well.  The operation
\[
  m_0(a,b,c)=
  \begin{cases}
    c,&a=b,\\
    a,&a\ne b
  \end{cases}
\]
is Mal'tsev and preserves every full relation.  Thus $\CC(\cF)$ holds; the
same argument applies to every $\CC^\delta(\cF)$ because constraint-free
instances have degree zero.  The decision procedures accept this case
immediately, and henceforth we may assume $\cF\ne\varnothing$.

\section{Purification Invariance and Arity Compression}
\label{sec:purification-compression}

We first remove the ambient-family dependence of purification.

\begin{lemma}[Ambient invariance of purification]
\label{lem:purification-invariance}
Let $A:D^r\to\AlgNums$, where $r\ge2$, be a table that occurs as a
distinguished component of each of two finite ambient tuples.  Legally purify
the two tuples, and let $P$ and $Q$ be the two purified tables occupying the
distinguished position formerly occupied by $A$.  Either printed
purifications or their positive table-wide normalizations may be used.  Then:
\begin{enumerate}
  \item the two purified tables have the same support;
  \item disjointness and positive proportionality of every pair of magnitude
        rows agree under $P$ and $Q$;
  \item complex linear dependence of every pair of nonzero rows agrees under
        $P$ and $Q$; and
  \item for proportional magnitude rows, the two purifications induce the
        same equal-magnitude blocks, and the Hermitian sum on each such block
        vanishes under $P$ if and only if it vanishes under $Q$.
\end{enumerate}
In particular, $P$ is block-orthogonal if and only if $Q$ is
block-orthogonal.
\end{lemma}

\begin{proof}
If $A$ is the zero table, both $P$ and $Q$ are zero tables and every assertion
is immediate.  Assume henceforth that $A$ has a nonzero entry.
Let
\[
  \Gamma=\langle a:a\text{ is a nonzero entry of }A\rangle
  \le\AlgNums^\times.
\]
Thus $\Tor(\Gamma)=\Gamma\cap\mu_\infty$ consists exactly of those
multiplicative combinations of entries of $A$ that are roots of unity.
Fix one of the two ambient purifications, with
nonzero ambient values $c_1,\ldots,c_t$, legal generators
$g_1,\ldots,g_m$, and exponent vectors
$e(c_j)=(e_1(c_j),\ldots,e_m(c_j))$.  Put
\[
  H=\langle c_1,\ldots,c_t\rangle.
\]
Every nonzero entry of $A$ is one of the ambient values, so $\Gamma\le H$.
We claim that the prescribed exponent vectors define
\begin{equation}
  e\left(\prod_{j=1}^t c_j^{n_j}\right)
  =\sum_{j=1}^t n_j e(c_j)
  \qquad(n_1,\ldots,n_t\in\mathbb Z).
  \label{eq:extended-exponent-map}
\end{equation}
The point requiring proof is that this value is independent of the chosen
representation.  If
\[
  \prod_{j=1}^t c_j^{n_j}=1,
\]
then~\eqref{eq:multiplicative-decomposition} gives
\begin{equation}
  \prod_{i=1}^m
       g_i^{\sum_{j=1}^t n_j e_i(c_j)}
  =\prod_{j=1}^t\vartheta(c_j)^{-n_j}\in\mu_\infty.
  \label{eq:ambient-exponent-relation}
\end{equation}
Condition~2 in the generating-set definition above forces
\[
  \sum_{j=1}^t n_j e_i(c_j)=0
  \qquad\text{for every }i\in[m].
\]
More generally, if two products in the $c_j$ represent the same element of
$H$, their quotient represents $1$, and the same calculation shows that the
two sums in~\eqref{eq:extended-exponent-map} agree.  Hence the displayed rule
is well-defined.  It is a homomorphism because multiplication of products
adds exponent vectors, and it is the unique homomorphism taking each $c_j$
to the prescribed $e(c_j)$.  Restricting this homomorphism to $\Gamma$ gives
$e:\Gamma\to\mathbb Z^m$.

Its kernel is exactly $\Tor(\Gamma)$.  Indeed, write
$x=\prod_j c_j^{n_j}$.  If $e(x)=0$, then the decompositions of the $c_j$
give
\[
  x=\prod_{j=1}^t\vartheta(c_j)^{n_j}\in\mu_\infty.
\]
Conversely, if
$x^N=1$ for some $N\ge1$, then $N e(x)=e(x^N)=0$, and the torsion-free group
$\mathbb Z^m$ gives $e(x)=0$.  The exponent map therefore descends to an
injective homomorphism
\[
  \nu:\Gamma/\Tor(\Gamma)\hookrightarrow\mathbb Z^m,
  \qquad \nu([x])=e(x).
\]
Define, for $x\in\Gamma$,
\[
  \vartheta(x)=x\prod_{i=1}^m g_i^{-e_i(x)}.
\]
The preceding well-definedness argument shows that $\vartheta$ is a
homomorphism into $\mu_\infty$.  It fixes torsion pointwise in the sense of
\eqref{eq:fixes-torsion-definition}: if $\xi\in\Tor(\Gamma)$, then the
exponent map annihilates $\xi$, so
\[
  e(\xi)=0,
  \qquad
  \vartheta(\xi)=\xi.
\]
Thus
\[
  \psi(x)=\vartheta(x)\prod_{i=1}^m p_i^{e_i(x)}
\]
is an injective homomorphism $\Gamma\to\mathbb C^\times$ that also fixes
torsion pointwise, because $\psi(\xi)=\vartheta(\xi)=\xi$ whenever
$e(\xi)=0$.  Here the \emph{multiplicative independence} of the distinct
positive primes means
\[
  \begin{gathered}
    \prod_{i=1}^m p_i^{k_i}=1,
    \qquad k_i\in\mathbb Z,\\
    \Longrightarrow\qquad k_1=\cdots=k_m=0;
  \end{gathered}
\]
this is unique factorization in $\mathbb Q_{>0}$.  It implies injectivity:
if $\psi(x)=1$, taking absolute values forces $e(x)=0$, after which
$x\in\Tor(\Gamma)$ and $\psi(x)=\vartheta(x)=x$, so $x=1$.

We now make the denominator-clearing step explicit.  The printed
purification replaces each nonzero original value $a$ by $\psi(a)$.  If some
$e_i(a)$ is negative, the positive rational number
$\prod_i p_i^{e_i(a)}$ can have a denominator.  Since $A$ has only finitely
many values, one can choose a positive integer $N_A$ such that
\[
  \begin{gathered}
    N_A\prod_{i=1}^m p_i^{e_i(a)}\in\mathbb Z_{>0}\\
    \text{for every nonzero entry }a\text{ of }A.
  \end{gathered}
\]
The normalized purified entry is then $N_A\psi(a)$.  Thus $\psi$ is the
homomorphism \emph{before denominator clearing}; multiplying its values on
one table by $N_A$ produces the pure table.  The normalized entry map need
not be multiplicative.  Every product comparison and Hermitian sum used
below nevertheless acquires the same nonzero factor $N_A^2$, so equality and
vanishing are preserved.

The crucial consequence is
\begin{equation}
  \abs{\psi(x)}=1
  \quad\Longleftrightarrow\quad
  e(x)=0
  \quad\Longleftrightarrow\quad
  x\in\Tor(\Gamma).
  \label{eq:purification-unit-modulus}
\end{equation}
Indeed, $\vartheta(x)$ has modulus one, and unique factorization in
$\mathbb Q_{>0}$ gives $\prod_i p_i^{e_i(x)}=1$ exactly when every exponent
vanishes.

Applying this construction to the two ambient purifications gives
homomorphisms $\psi_P$ and $\psi_Q$ and positive denominator-clearing factors
$N_P$ and $N_Q$, where the factor is $1$ if the printed form is retained.
For a nonzero original value $a$ of $A$, write
\[
  P(a)=N_P\psi_P(a),
  \qquad
  Q(a)=N_Q\psi_Q(a).
\]
This notation means the entry replacing the original value $a$ in the
distinguished purified table.  The factors $N_P$ and $N_Q$ cancel from every
comparison below.

Fixing torsion survives denominator clearing in the precise covariance form
needed later.  For $\xi\in\Tor(\Gamma)$ and $x\in\Gamma$,
\begin{equation}
  \psi_P(\xi x)=\xi\psi_P(x),
  \qquad
  \psi_Q(\xi x)=\xi\psi_Q(x).
  \label{eq:purification-torsion-covariance}
\end{equation}
Consequently, if both $a$ and $\xi a$ occur as entries of $A$, then
\[
  P(\xi a)=\xi P(a),
  \qquad
  Q(\xi a)=\xi Q(a).
\]
Thus denominator clearing preserves every root-of-unity ratio between
entries, while representing a torsion entry $\xi$ itself by $N_P\xi$ or
$N_Q\xi$.

For nonzero entries $a,b,c,u$ of $A$, applying
\eqref{eq:purification-unit-modulus} to the ratios $a/b$ and $au/(bc)$ gives
\begin{equation}
  \begin{gathered}
  \abs{P(a)}=\abs{P(b)}\\
  \Longleftrightarrow\quad a/b\in\Tor(\Gamma)\\
  \Longleftrightarrow\quad \abs{Q(a)}=\abs{Q(b)}.
  \end{gathered}
  \label{eq:pure-equal-magnitude}
\end{equation}
Moreover,
\begin{equation*}
  \begin{gathered}
  \abs{P(a)P(u)}=\abs{P(b)P(c)}\\
  \Longleftrightarrow\quad au/(bc)\in\Tor(\Gamma)\\
  \Longleftrightarrow\quad \abs{Q(a)Q(u)}=\abs{Q(b)Q(c)}.
  \end{gathered}
\end{equation*}
Zeros remain zero, so support is invariant.  The second equivalence, with
the usual separate treatment of zero entries, gives disjointness and
positive proportionality of magnitude rows.  The first gives the common
partition into equal-magnitude blocks.

Ordinary complex row dependence is determined by vanishing $2\times2$
minors.
Write the two rows under comparison as
$\mathbf a=(a_z)_{z\in D}$ and $\mathbf b=(b_z)_{z\in D}$.
For every $z,w\in D$, the corresponding original minor is
\[
  \det\begin{pmatrix}a_z&a_w\\ b_z&b_w\end{pmatrix}
  =a_zb_w-a_wb_z.
\]
For nonzero corners, injectivity and multiplicativity give
\[
  \begin{gathered}
  P(a_z)P(b_w)=P(a_w)P(b_z)\\
  \Longleftrightarrow\quad a_z b_w=a_w b_z\\
  \Longleftrightarrow\quad Q(a_z)Q(b_w)=Q(a_w)Q(b_z).
  \end{gathered}
\]
The equivalence is immediate as well when a corner is zero.  Thus all minors
of a row pair vanish under $P$ exactly when they vanish under $Q$.

It remains to compare a Hermitian block sum.  By this we mean the restriction
of the usual complex inner product to an equal-magnitude block $T$:
\[
  \sum_{z\in T}P(a_z)\overline{P(b_z)}.
\]
Suppose the two magnitude rows are proportional, let $T$ be one of their
common equal-magnitude blocks, and fix $z_0\in T$.  By
\eqref{eq:pure-equal-magnitude}, there are roots of unity
$\xi_z,\eta_z\in\Tor(\Gamma)$ such that
\[
  a_z=\xi_z a_{z_0},\qquad b_z=\eta_z b_{z_0}
  \qquad(z\in T).
\]
By \eqref{eq:purification-torsion-covariance}, the underlying homomorphisms
$\psi_P$ and $\psi_Q$ preserve the displayed root-of-unity ratios; each
purified table is one common positive scalar times its underlying image.
Consequently,
\begin{align*}
  \sum_{z\in T}P(a_z)\overline{P(b_z)}
  &=P(a_{z_0})\overline{P(b_{z_0})}
    \sum_{z\in T}\xi_z\overline{\eta_z},\\
  \sum_{z\in T}Q(a_z)\overline{Q(b_z)}
  &=Q(a_{z_0})\overline{Q(b_{z_0})}
    \sum_{z\in T}\xi_z\overline{\eta_z}.
\end{align*}
The two prefactors are nonzero, so the two block sums vanish simultaneously.
This proves all four assertions.
\end{proof}

\begin{corollary}[Singleton purification suffices]
\label{cor:singleton-purification}
The joint condition $\BO(\cF)$ holds if and only if every
$G\in W_{\cF}$ of arity at least two has a block-orthogonal singleton
purification.
\end{corollary}

\begin{proof}
If the joint condition holds, apply it to the one-element tuple $(G)$ for
each $G\in W_{\cF}$.

Conversely, assume every $G\in W_{\cF}$ of arity at least two has a
block-orthogonal singleton purification.  Fix an arbitrary nonempty finite
tuple
\[
  \mathbf G=(G_1,\ldots,G_h)
\]
from $W_{\cF}$ and legally purify it jointly:
\[
  \Pure(\mathbf G)=(J_1,\ldots,J_h).
\]
For each $i$ with $\arity(G_i)\ge2$, also let $S_i$ be a legal singleton
purification of $G_i$.  Apply \cref{lem:purification-invariance} with
distinguished table $A=G_i$ and with the two ambient tuples
\[
  (G_i)
  \qquad\text{and}\qquad
  (G_1,\ldots,G_h).
\]
The two distinguished outputs in that lemma are precisely $S_i$ and $J_i$.
The nonzero values in the other tables may enlarge the ambient multiplicative
group and change the exponent coordinates, the assigned primes, and hence the
numerical entries of $J_i$.  The intrinsic subgroup remains
\[
  \Tor(\Gamma_i),
  \qquad
  \Gamma_i=\langle a:a\text{ is a nonzero entry of }G_i\rangle,
\]
because, if $H$ is the larger ambient group, then
$\Tor(H)\cap\Gamma_i=\Tor(\Gamma_i)$.  Hence
\cref{lem:purification-invariance} applies to the two ambient tuples and gives
that $S_i$ is block-orthogonal if and only if $J_i$ is
block-orthogonal.
By assumption the left side holds, so every jointly purified component of
arity at least two is block-orthogonal.  As $\mathbf G$ was arbitrary, the
joint condition $\BO(\cF)$ holds.
\end{proof}

We next establish a bound on the external arity of a counterexample.

\begin{lemma}[External diagonal minors]
\label{lem:diagonal-minors}
Let $G\in W_{\cF}$ have arity $r$, let $1\le s\le r$, and let
$\pi:[r]\twoheadrightarrow[s]$ be a surjection.  Then
\[
  H(y_1,\ldots,y_s)
  =G(y_{\pi(1)},\ldots,y_{\pi(r)})
\]
belongs to $W_{\cF}$.
\end{lemma}

\begin{proof}
Choose a base-language instance presenting $G$ as a partial marginal.
Replace its retained variable in position $i$ by $y_{\pi(i)}$, while keeping
all previously marginalized variables distinct and retaining the same sum.
Repeated variables in scopes are legal, so the modified object is again an
instance.  Its marginal is exactly $H$, including every cancellation already
present in $G$.
\end{proof}

To see the instance construction explicitly, suppose a presentation of a
five-ary $G$ has retained variables $x_1,\ldots,x_5$, hidden variables
$z_1,z_2$, and constraint multiset
\begin{equation*}
\begin{aligned}
  \mathcal C_I=\{&f(x_1,z_1,x_3),
                  g(z_1,x_2,x_5),\\
                 &h(x_4,z_2,x_3),
                  \ell(z_2,x_5)\}.
\end{aligned}
\end{equation*}
For $\pi=(1,2,1,3,2)$, construct $I_\pi$ by keeping every constraint and
hidden variable and replacing only retained-variable incidences:
\[
\begin{aligned}
  \mathcal C_{I_\pi}=\{&f(y_1,z_1,y_1),
                        g(z_1,y_2,y_2),\\
                       &h(y_3,z_2,y_1),
                        \ell(z_2,y_2)\}.
\end{aligned}
\]

\begin{figure}[!htbp]
  \centering
  \resizebox{0.94\linewidth}{!}{%
  \begin{tikzpicture}[
      x=1cm,y=1cm,
      var/.style={circle,draw,minimum size=5.8mm,inner sep=0pt,
                  font=\scriptsize},
      hidden/.style={var,dashed},
      factor/.style={draw,rounded corners=1pt,minimum width=6mm,
                     minimum height=5.5mm,inner sep=1pt,
                     double,double distance=.45pt,font=\scriptsize},
      edge/.style={line width=.42pt},
      title/.style={font=\small\bfseries},
      scopes/.style={font=\scriptsize,align=center}
    ]
    \begin{scope}[shift={(2.1,0)}]
      \node[title] at (1.2,2.2) {original instance $I$};
      \node[factor] (f) at (0,1) {$f$};
      \node[factor] (g) at (2.4,1) {$g$};
      \node[factor] (h) at (0,-1) {$h$};
      \node[factor] (l) at (2.4,-1) {$\ell$};

      \node[hidden] (z1) at (1.2,1) {$z_1$};
      \node[hidden] (z2) at (1.2,-1) {$z_2$};
      \node[var] (x3) at (0,0) {$x_3$};
      \node[var] (x5) at (2.4,0) {$x_5$};
      \node[var] (x1) at (-1.15,1.62) {$x_1$};
      \node[var] (x2) at (3.55,1.62) {$x_2$};
      \node[var] (x4) at (-1.15,-1.62) {$x_4$};

      \draw[edge] (f)--(x1) (f)--(z1) (f)--(x3);
      \draw[edge] (g)--(z1) (g)--(x2) (g)--(x5);
      \draw[edge] (h)--(x4) (h)--(z2) (h)--(x3);
      \draw[edge] (l)--(z2) (l)--(x5);

      \node[scopes] at (1.2,-2.25)
        {$f(x_1,z_1,x_3),\ g(z_1,x_2,x_5)$\\
         $h(x_4,z_2,x_3),\ \ell(z_2,x_5)$};
    \end{scope}

    \draw[-{Latex[length=2mm]},line width=.55pt]
      (6.05,0)--(7.65,0)
      node[midway,above=2pt,font=\scriptsize,align=center]
      {$x_1,x_3\mapsto y_1$\\$x_2,x_5\mapsto y_2$\\$x_4\mapsto y_3$};

    \begin{scope}[shift={(9.5,0)}]
      \node[title] at (1.2,2.2) {constructed instance $I_\pi$};
      \node[factor] (fp) at (0,1) {$f$};
      \node[factor] (gp) at (2.4,1) {$g$};
      \node[factor] (hp) at (0,-1) {$h$};
      \node[factor] (lp) at (2.4,-1) {$\ell$};

      \node[hidden] (z1p) at (1.2,1) {$z_1$};
      \node[hidden] (z2p) at (1.2,-1) {$z_2$};
      \node[var] (y1) at (0,0) {$y_1$};
      \node[var] (y2) at (2.4,0) {$y_2$};
      \node[var] (y3) at (-1.15,-1.62) {$y_3$};

      \draw[edge] (fp)--(z1p);
      \draw[edge] (gp)--(z1p);
      \draw[edge] (hp)--(z2p);
      \draw[edge] (lp)--(z2p);
      \draw[edge] (fp) to[bend left=17] (y1);
      \draw[edge] (fp) to[bend right=17] (y1);
      \draw[edge] (hp)--(y1);
      \draw[edge] (gp) to[bend left=17] (y2);
      \draw[edge] (gp) to[bend right=17] (y2);
      \draw[edge] (lp)--(y2);
      \draw[edge] (hp)--(y3);

      \node[scopes] at (1.2,-2.25)
        {$f(y_1,z_1,y_1),\ g(z_1,y_2,y_2)$\\
         $h(y_3,z_2,y_1),\ \ell(z_2,y_2)$};
    \end{scope}
  \end{tikzpicture}%
  }
  \caption{Construction in \cref{lem:diagonal-minors}.  Circles are
  variables, boxes are constraints, and dashed circles are hidden variables.
  Constraint nodes and hidden variables are unchanged; only incidences of
  retained variables are relabelled.  The parallel edges at $f$ and $g$ in
  $I_\pi$ record repeated occurrences created by identification.  Summing
  over the unchanged $z_1,z_2$ gives
  $F_{I_\pi}^{[3]}(y_1,y_2,y_3)
   =G(y_1,y_2,y_1,y_3,y_2)=H(y_1,y_2,y_3)$.}
  \label{fig:diagonal-minor-example}
\end{figure}

\begin{lemma}[Bounded external arity]
\label{lem:bo-arity-bound}
If the singleton Block Orthogonality condition fails on $W_{\cF}$, then it
fails for some member of $W_{\cF}$ of arity at most
\[
  R_{\mathrm{BO}}=d^2+1.
\]
\end{lemma}

\begin{proof}
Let $G:D^r\to\AlgNums$, where $r\ge2$, be a counterexample.  Its failure is
witnessed by two
row indices $\mathbf x,\mathbf y\in D^{r-1}$: their purified magnitude rows
overlap but are not proportional, or they are proportional while the complex
rows are neither dependent nor block-orthogonal.  On the row-coordinate
positions $[r-1]$, put
\[
  i\equiv j
  \quad\Longleftrightarrow\quad
  (x_i,y_i)=(x_j,y_j).
\]
There are at most $d^2$ equivalence classes, one for each possible ordered
pair in $D^2$.  For a concrete example with $D=\{0,1\}$ (so $d=2$), take
\[
\begin{aligned}
  \mathbf x&=(0,0,1,0,1,1,0,1),\\
  \mathbf y&=(0,1,0,0,1,0,1,1).
\end{aligned}
\]
\Cref{tab:pair-class-example} displays the four equivalence classes.  Identify
positions in each class and keep the final column variable separate.  By
\cref{lem:diagonal-minors}, the resulting table $H$ belongs to $W_{\cF}$ and
has arity at most $d^2+1$.

\begin{table}[!htbp]
  \centering
  \small
  \setlength{\tabcolsep}{3.6pt}
  \renewcommand{\arraystretch}{1.18}
  \begin{tabular}{c|cccc}
    class
      &$C_{00}\,(\bullet)$
      &$C_{01}\,(\circ)$
      &$C_{10}\,(\blacktriangle)$
      &$C_{11}\,(\diamond)$\\ \hline
    positions &$\{1,4\}$&$\{2,7\}$&$\{3,6\}$&$\{5,8\}$\\
    new variable &$u_{00}$&$u_{01}$&$u_{10}$&$u_{11}$\\
    $\mathbf x$-value &$0$&$0$&$1$&$1$\\
    $\mathbf y$-value &$0$&$1$&$0$&$1$
  \end{tabular}
  \caption{The position classes for the displayed Boolean rows.  Each
  black-and-white marker identifies one class of positions carrying the same
  ordered pair $(x_i,y_i)$.}
  \label{tab:pair-class-example}
\end{table}

In this example the identified table is
\[
\begin{aligned}
  H(u_{00},u_{01},u_{10},u_{11},z)
   =G(&u_{00},u_{01},u_{10},u_{00},\\[-2pt]
      &u_{11},u_{10},u_{01},u_{11},z).
\end{aligned}
\]
Its two selected row assignments are respectively
\[
  (u_{00},u_{01},u_{10},u_{11})=(0,0,1,1)
  \quad\text{and}\quad
  (0,1,0,1),
\]
which expand back to $\mathbf x$ and $\mathbf y$.  There are four identified
row variables and one separate column variable, so the resulting arity is
$4+1=d^2+1=5$.

Assign to each new row variable the first component of its defining ordered
pair to obtain one row, and the second component to obtain the other.  The two
resulting rows of $H$ are entry-for-entry the selected rows of $G$.  Take a
legal joint printed purification
$(\widetilde G,\widetilde H)$ of $(G,H)$.  By
\cref{lem:purification-invariance}, the failed singleton test for $G$ is
preserved in $\widetilde G$.  A joint printed purification applies the same
map to equal original entries, so the selected rows of $\widetilde G$ agree
entry-for-entry with the corresponding rows of $\widetilde H$.  The same
test therefore fails in $\widetilde H$.  Applying
\cref{lem:purification-invariance} to $H$ transfers this failure to every
legal singleton purification of $H$.  Hence $H$ is also a counterexample.
\end{proof}

\section{Algebraic Certificates for Block Orthogonality}
\label{sec:bo-certificates}

This section expresses Block Orthogonality of a singleton purification $P$
directly in the original coordinates of its table $A$.
\Cref{lem:intrinsic-purification-tests} leaves support and ordinary minors
unchanged and converts purified magnitude comparisons and Hermitian block
cancellations into root-of-unity equations among the entries of $A$.
Consequently, both the certificates and the identity oracle operate directly
on the original table.

Fix the computable conjugation-stable number field
\begin{equation}
  K=\mathbb Q\left(
      \{f(\mathbf a),\overline{f(\mathbf a)}:
        f\in\cF,\ \mathbf a\in D^{r_f}\}
    \right),
  \label{eq:working-field}
\end{equation}
and let $\mu(K)$ be its finite, effectively computable group of roots of
unity~\cite{Cohen1993}.  By definition, every input value
$f(\mathbf a)$ belongs to $K$.  Equations~\eqref{eq:instance-function}
and~\eqref{eq:partial-marginal} express each entry of every
$G\in W_{\cF}$ as a finite sum of finite products of such values.  Since
$K$ is a field, every entry of every generated table therefore belongs to
$K$.  We now give an effective description of the $K$-valued tables whose
singleton purifications are block-orthogonal; all relevant torsion choices
lie in the computable finite group $\mu(K)$.

\subsection{Intrinsic torsion tests}

The following consequences of purification convert every magnitude
comparison needed by Block Orthogonality into a finite root-of-unity choice.

\Needspace{8\baselineskip}
\begin{lemma}[Intrinsic purification tests]
\label{lem:intrinsic-purification-tests}
Let $A$ be a finite $K$-valued table, let $P$ be any legal singleton
purification of $A$, and, for every nonzero value $a$ occurring in $A$,
write $\widehat a$ for its purified value.  For nonzero values $a,b,c,u$
occurring in $A$,
\begin{align}
  \abs{\widehat a}=\abs{\widehat b}
  &\quad\Longleftrightarrow\quad
  a/b\in\mu(K),
  \label{eq:intrinsic-equal-magnitude}\\
  \abs{\widehat a}\abs{\widehat u}
     =\abs{\widehat b}\abs{\widehat c}
  &\quad\Longleftrightarrow\quad
  au/(bc)\in\mu(K)\notag\\
  &\quad\Longleftrightarrow\quad
  au=\zeta bc
  \text{ for some }\zeta\in\mu(K),
  \label{eq:intrinsic-magnitude-minor}\\
  \widehat a\,\widehat u=\widehat b\,\widehat c
  &\quad\Longleftrightarrow\quad
  au=bc.
  \label{eq:intrinsic-ordinary-minor}
\end{align}
Moreover, if $a=\rho b$ for $\rho\in\mu(K)$, then
\begin{equation}
  \widehat a=\rho\,\widehat b.
  \label{eq:purification-root-covariance}
\end{equation}
\end{lemma}

\begin{proof}
If $A$ is the zero table, the assertions are vacuous.  Assume that $A$ has a
nonzero entry.
Let $\Gamma\le K^\times$ be the multiplicative group generated by the
nonzero entries of $A$.  As in the proof of
\cref{lem:purification-invariance}, equation
\eqref{eq:ambient-exponent-relation} shows that the legal generating set
induces homomorphisms
\[
  \nu:\Gamma/\Tor(\Gamma)\hookrightarrow\mathbb Z^m,
  \qquad
  \vartheta:\Gamma\longrightarrow\mu(K),
\]
where $\vartheta$ records the residual root of unity and restricts to the
identity on $\Tor(\Gamma)$.  Here $\vartheta(\Gamma)\subseteq\mu(K)$ because
this is a singleton purification of a $K$-valued table and every generator
lies in the field generated over $\mathbb Q$ by the nonzero entries of $A$,
which is a subfield of $K$.  Before denominator clearing, the purification
map is the injective homomorphism
\[
  \varphi(x)=\vartheta(x)
      \prod_{i=1}^m p_i^{\nu_i([x])}
  \qquad(x\in\Gamma).
\]
There is one positive integer $N$, independent of the entry, such that
$\widehat a=N\varphi(a)$ for every nonzero entry $a$.

Multiplicative independence of the primes gives
\[
  \abs{\varphi(a/b)}=1
  \quad\Longleftrightarrow\quad
  a/b\in\Tor(\Gamma)=\Gamma\cap\mu(K).
\]
This proves~\eqref{eq:intrinsic-equal-magnitude}; applying the same argument
to $au/(bc)$ proves~\eqref{eq:intrinsic-magnitude-minor}.  The common factor
$N^2$ cancels in~\eqref{eq:intrinsic-ordinary-minor}, after which the claim
follows from injectivity of $\varphi$.  Finally, $\varphi$ fixes torsion
pointwise as in \eqref{eq:fixes-torsion-definition}; the common
denominator-clearing factor then gives
\eqref{eq:purification-root-covariance}.
\end{proof}

In particular, ordinary row dependence is preserved by purification: all
ordinary $2\times2$ minors of two original rows vanish if and only if all
corresponding minors of their purified rows vanish.

\subsection{Finite certificate branches}

Fix an external arity $k\ge2$, put
\[
  \mathcal R_k=D^{k-1},
\]
and regard a $k$-ary table as a matrix
\[
  A=(a_{\mathbf x,c})_{\mathbf x\in\mathcal R_k,\ c\in D}.
\]
Let
\[
  \mathbf X=\{X_{\mathbf x,c}:
      \mathbf x\in\mathcal R_k,\ c\in D\}
\]
be the corresponding coordinate variables.  Fix arbitrary total orders on
$\mathcal R_k$ and $D$.

A \emph{Block Orthogonality certificate} of arity $k$ consists of the
following finite discrete data.  First choose a rectangle system
\[
  (R_1,C_1),\ldots,(R_t,C_t),
\]
where every $R_i\subseteq\mathcal R_k$ and $C_i\subseteq D$ is nonempty,
the row sets $R_i$ are pairwise disjoint, and the column sets $C_i$ are
pairwise disjoint.  The empty system $t=0$ is allowed.  Impose the support
equations
\begin{equation}
  X_{\mathbf x,c}=0
  \qquad
  \text{whenever }
  (\mathbf x,c)\notin
  \bigcup_{i=1}^t(R_i\times C_i).
  \label{eq:bo-certificate-support}
\end{equation}

For every $i$, every two distinct rows
$\mathbf x,\mathbf y\in R_i$, and every two distinct columns
$c,e\in C_i$, choose
$\tau_{\mathbf x,\mathbf y;c,e}\in\mu(K)$ and impose the twisted minor
\begin{equation}
  X_{\mathbf x,c}X_{\mathbf y,e}
  -\tau_{\mathbf x,\mathbf y;c,e}
   X_{\mathbf x,e}X_{\mathbf y,c}=0.
  \label{eq:bo-certificate-twisted-minor}
\end{equation}
For every unordered pair $\mathbf x<\mathbf y$ in one $R_i$, choose one of
the following two modes.

\paragraph{Dependent mode.}
Impose every ordinary minor
\begin{equation}
  X_{\mathbf x,c}X_{\mathbf y,e}
  -X_{\mathbf x,e}X_{\mathbf y,c}=0
  \qquad(c,e\in C_i).
  \label{eq:bo-certificate-dependent}
\end{equation}

\paragraph{Orthogonal mode.}
Choose a set partition $\mathcal T_{\mathbf x,\mathbf y}$ of $C_i$.  For
every $T\in\mathcal T_{\mathbf x,\mathbf y}$, choose an anchor $c_T\in T$
and roots
\[
  \rho_{T,c},\sigma_{T,c}\in\mu(K)
  \qquad(c\in T),
\]
with $\rho_{T,c_T}=\sigma_{T,c_T}=1$.  Retain this discrete choice only if
\begin{equation}
  \sum_{c\in T}\rho_{T,c}\overline{\sigma_{T,c}}=0
  \qquad
  \text{for every }T\in\mathcal T_{\mathbf x,\mathbf y}.
  \label{eq:bo-certificate-root-sum}
\end{equation}
For every $T\in\mathcal T_{\mathbf x,\mathbf y}$ and every
$c\in T\setminus\{c_T\}$, impose the anchor equations
\begin{align}
  X_{\mathbf x,c}-\rho_{T,c}X_{\mathbf x,c_T}&=0,
  \label{eq:bo-certificate-first-anchor}\\
  X_{\mathbf y,c}-\sigma_{T,c}X_{\mathbf y,c_T}&=0.
  \label{eq:bo-certificate-second-anchor}
\end{align}

For a retained certificate $\mathfrak c$, rewrite every coordinate equation
imposed above as $p(\mathbf X)=0$ and put the polynomial
$p(\mathbf X)$ into $E_{\mathfrak c}\subseteq K[\mathbf X]$.  Explicitly,
$E_{\mathfrak c}$ consists of all applicable polynomials of the following
types, with the indices and discrete choices specified by
$\mathfrak c$:
\begin{align*}
  &X_{\mathbf x,c},\\
  &X_{\mathbf x,c}X_{\mathbf y,e}
    -\tau_{\mathbf x,\mathbf y;c,e}
     X_{\mathbf x,e}X_{\mathbf y,c},\\
  &X_{\mathbf x,c}X_{\mathbf y,e}
    -X_{\mathbf x,e}X_{\mathbf y,c},\\
  &X_{\mathbf x,c}-\rho_{T,c}X_{\mathbf x,c_T},\\
  &X_{\mathbf y,c}-\sigma_{T,c}X_{\mathbf y,c_T}.
\end{align*}
The third type is imposed only in dependent mode, and the last two only in
orthogonal mode.
Thus ``left-hand side'' means precisely the polynomial before $=0$ in one
of~\eqref{eq:bo-certificate-support},
\eqref{eq:bo-certificate-twisted-minor},
\eqref{eq:bo-certificate-dependent}, or
\cref{eq:bo-certificate-first-anchor,eq:bo-certificate-second-anchor}.
Certificate enumeration retains exactly the discrete root choices satisfying
\eqref{eq:bo-certificate-root-sum}.  For each retained choice,
$E_{\mathfrak c}$ consists of homogeneous polynomials in the unbarred
coordinate variables: degree one for support and anchor equations, and degree
two for minor equations.  Conjugation is confined to the finite
precomputation over $\mu(K)$.

The certificate family is finite and effectively enumerable because the row
and column sets, their partitions, and $\mu(K)$ are finite.  Retention is
decided by exact number-field equality tests, and homogeneity makes the zero
table a solution of every retained branch.  Let $\mathfrak C_k$ denote the
family of retained Block Orthogonality certificates of arity $k$.

For $E\subseteq K[\mathbf X]$, write
\[
  V_K(E)=\{A\in K^{D^k}:e(A)=0\text{ for every }e\in E\}.
\]
Let $\mathcal B_k$ denote the set of $K$-valued $k$-ary tables whose
singleton purifications are block-orthogonal.  This is well-defined by
\cref{lem:purification-invariance}: whether the purification is
block-orthogonal is independent of the legal singleton purification chosen.
Although $\mathcal B_k$ is defined through a purification $P$ of an original
table $A$, every coordinate variable $X_{\mathbf x,c}$ in
$E_{\mathfrak c}$ represents $a_{\mathbf x,c}$, not
$p_{\mathbf x,c}$.  The passage from $P$ back to $A$ uses exactly
\cref{lem:intrinsic-purification-tests}:
\eqref{eq:intrinsic-magnitude-minor} gives the twisted-minor equations,
\eqref{eq:intrinsic-ordinary-minor} gives the ordinary-minor equations, and
\cref{eq:intrinsic-equal-magnitude,eq:purification-root-covariance} give
the anchor equations and translate each purified Hermitian block sum into
the prechecked root sum.  Support is unchanged by purification.  Hence the
sets $V_K(E_{\mathfrak c})$ below are loci of original tables even though
they characterize Block Orthogonality after purification.

\begin{theorem}[Algebraic description of the BO locus]
\label{thm:finite-bo-certificates}
For every $k\ge2$, the effectively enumerable certificate family
$\mathfrak C_k$ satisfies
\begin{equation}
  \mathcal B_k
  =\bigcup_{\mathfrak c\in\mathfrak C_k}V_K(E_{\mathfrak c}).
  \label{eq:bo-certificate-union}
\end{equation}
\end{theorem}

\begin{proof}
We prove the two inclusions separately.

\paragraph{Every block-orthogonal table has a certificate.}
Let $A\in\mathcal B_k$, and fix a singleton purification
$P=(p_{\mathbf x,c})$ of $A$.  The support of $P$ equals that of $A$.
Because $\abs P$ is block-rank one, the nonempty connected components of its
bipartite support graph are complete bipartite graphs
\[
  R_1\times C_1,\ldots,R_t\times C_t,
\]
with pairwise disjoint row sets and pairwise disjoint column sets.  Choose
these components as the certificate rectangles.  This gives
\eqref{eq:bo-certificate-support}.

Within one rectangle, purified magnitude rows are proportional.  Hence, for
$\mathbf x,\mathbf y\in R_i$ and $c,e\in C_i$,
\[
  \abs{p_{\mathbf x,c}}\abs{p_{\mathbf y,e}}
  =\abs{p_{\mathbf x,e}}\abs{p_{\mathbf y,c}}.
\]
All four original entries are nonzero.  By
\eqref{eq:intrinsic-magnitude-minor}, there is a root
$\tau_{\mathbf x,\mathbf y;c,e}\in\mu(K)$ for which the corresponding
twisted minor vanishes.

Fix $\mathbf x<\mathbf y$ in $R_i$.  If their purified rows are linearly
dependent, then~\eqref{eq:intrinsic-ordinary-minor} shows that every ordinary
minor of the original rows vanishes.  We choose dependent mode.

Otherwise, Block Orthogonality partitions $C_i$ into the nonempty level sets
$T$ of $c\mapsto\abs{p_{\mathbf x,c}}$.  Since the magnitude rows are
proportional, these are also the level sets for the other row.  Choose any
anchor $c_T\in T$.  For every $c\in T$,
\[
  \rho_{T,c}=\frac{a_{\mathbf x,c}}{a_{\mathbf x,c_T}},
  \qquad
  \sigma_{T,c}=\frac{a_{\mathbf y,c}}{a_{\mathbf y,c_T}}
\]
belong to $\mu(K)$ by~\eqref{eq:intrinsic-equal-magnitude}.  These roots
satisfy the anchor equations.  Root covariance and block orthogonality on
$T$ give
\begin{align*}
  0
  &=\sum_{c\in T}p_{\mathbf x,c}\overline{p_{\mathbf y,c}}\\
  &=p_{\mathbf x,c_T}\overline{p_{\mathbf y,c_T}}
    \sum_{c\in T}\rho_{T,c}\overline{\sigma_{T,c}}.
\end{align*}
The prefactor is nonzero, so the root sum vanishes.  This constructs a
retained certificate containing $A$.

\paragraph{Every certified table is block-orthogonal.}
Fix a certificate $\mathfrak c$, let
$A\in V_K(E_{\mathfrak c})$, and let $P$ be a singleton purification of
$A$.  We allow every zero degeneration admitted by the homogeneous
certificate equations.

Every nonzero row belongs to a unique certified row set $R_i$, and its
support is contained in $C_i$.  Rows belonging to different certified
rectangles have disjoint supports because the corresponding column sets are
disjoint.

Consider two nonzero rows $\mathbf x,\mathbf y\in R_i$.  If their supports
are disjoint, they already satisfy the block-rank-one alternative.  Suppose
instead that they overlap, and choose $e\in C_i$ with
$a_{\mathbf x,e}a_{\mathbf y,e}\ne0$.  For every $c\in C_i\setminus\{e\}$,
the twisted minor gives
\[
  a_{\mathbf x,c}a_{\mathbf y,e}
  =\tau_{\mathbf x,\mathbf y;c,e}
   a_{\mathbf x,e}a_{\mathbf y,c}.
\]
Consequently,
$a_{\mathbf x,c}\ne0$ if and only if $a_{\mathbf y,c}\ne0$.
Thus overlapping rows have the same support.  On that common support,
\eqref{eq:intrinsic-magnitude-minor} applied to the same twisted equations
gives
\[
  \abs{p_{\mathbf x,c}}\abs{p_{\mathbf y,e}}
  =\abs{p_{\mathbf x,e}}\abs{p_{\mathbf y,c}},
\]
so their purified magnitude rows are positively proportional.  Hence
$\abs P$ is block-rank one.  A certified rectangle may vanish completely or
split into smaller surviving support rectangles; the argument covers both
possibilities.

It remains to verify the dependent-or-block-orthogonal alternative for an
overlapping pair.  In dependent mode,
\eqref{eq:bo-certificate-dependent} and
\eqref{eq:intrinsic-ordinary-minor} show that all purified ordinary minors
vanish.  Since both rows are nonzero, they are linearly dependent.

Suppose the pair is in orthogonal mode.  The anchor equations imply that,
in each certified group $T$, all entries of either row survive together or
vanish together.  A group cannot survive in only one row while some other
group survives in both rows.  Indeed, if $T$ survives only in row
$\mathbf x$, if $U$ survives in both, and if $c\in T$ and $e\in U$, then the
twisted minor has a nonzero left product and a zero right product.  The
symmetric case is identical.  Since the two rows overlap, at least one
common surviving group exists; hence every surviving group is common to
both rows.

For every common surviving group $T$, root covariance and the prechecked
root sum yield
\begin{align}
  \sum_{c\in T}p_{\mathbf x,c}\overline{p_{\mathbf y,c}}
  &=p_{\mathbf x,c_T}\overline{p_{\mathbf y,c_T}}
    \sum_{c\in T}\rho_{T,c}\overline{\sigma_{T,c}}\notag\\
  &=0.
  \label{eq:certified-group-remains-orthogonal}
\end{align}
The anchor equations also imply that, within $T$, the purified entries in each
row have a common magnitude.  Therefore a certified group cannot split into
smaller magnitude blocks after specialization.  Distinct certified
groups may acquire the same purified magnitude.  Because the two magnitude
rows are proportional, every resulting magnitude block is a disjoint union
of whole common surviving certified groups.  Summing
\eqref{eq:certified-group-remains-orthogonal} over those groups shows that
the Hermitian sum on the merged block is still zero.

The preceding cases cover every specialization permitted by the homogeneous
certificate equations.  Thus $P$ is block-orthogonal and
$A\in\mathcal B_k$.
\end{proof}

\begin{remark}[The role of the working field]
\label{rem:bo-certificates-k-points}
Equation~\eqref{eq:bo-certificate-union} is an equality inside $K^{D^k}$,
exactly the ambient space required here: every table in $W_{\cF}$
is $K$-valued by \cref{eq:instance-function,eq:partial-marginal}, and every
relevant torsion quotient therefore lies in $\mu(K)$.
\end{remark}

\subsection{From a finite union to polynomial identities}

Enumerate the certificate branches as
\[
  \mathfrak C_k=\{\mathfrak c_1,\ldots,\mathfrak c_s\}
\]
and write $E_j=E_{\mathfrak c_j}$, after deleting syntactically zero
polynomials.  If some $E_j$ is empty, then
$V_K(E_j)=K^{D^k}$ and $\mathcal B_k=K^{D^k}$; in that case define
$Q_k=\varnothing$.  Otherwise define the finite set
\begin{equation*}
  Q_k=
  \left\{
    \prod_{j=1}^s e_j:
    (e_1,\ldots,e_s)\in E_1\times\cdots\times E_s
  \right\}
  \subseteq K[\mathbf X].
\end{equation*}

\begin{corollary}[Polynomial identities for the BO locus]
\label{cor:bo-product-identities}
For every $A\in K^{D^k}$,
\begin{equation}
  A\in\mathcal B_k
  \quad\Longleftrightarrow\quad
  q(A)=0\text{ for every }q\in Q_k.
  \label{eq:bo-locus-by-identities}
\end{equation}
The finite set $Q_k$ is effectively computable from $D$, $k$, and $\mu(K)$.
\end{corollary}

\begin{proof}
If $A\in V_K(E_j)$ for some $j$, every polynomial $q\in Q_k$ contains a
factor from $E_j$, and that factor vanishes at $A$.  Hence every $q(A)$ is
zero.

Conversely, suppose that $A\notin V_K(E_j)$ for every $j$.  For each $j$,
choose $e_j\in E_j$ with $e_j(A)\ne0$.  Since $K$ is a field,
\[
  \left(\prod_{j=1}^s e_j\right)(A)\ne0.
\]
The resulting product belongs to $Q_k$, contradicting universal vanishing.
Together with \cref{thm:finite-bo-certificates}, this proves
\eqref{eq:bo-locus-by-identities}.  All choices used to enumerate the
$E_j$ and $Q_k$ range over finite explicitly represented sets.
\end{proof}

Combining the preceding results gives the following intrinsic test.
Let $G\in W_{\cF}$ have arity $k$, and let $P$ be any legal singleton
purification of $G$.  Then
\begin{equation*}
\begin{aligned}
  &P\text{ is block-orthogonal}\\
  &\quad\Longleftrightarrow\quad G\in\mathcal B_k\\
  &\quad\Longleftrightarrow\quad
    G\in\bigcup_{\mathfrak c\in\mathfrak C_k}V_K(E_{\mathfrak c})\\
  &\quad\Longleftrightarrow\quad
    q(G)=0\text{ for every }q\in Q_k.
\end{aligned}
\end{equation*}
The first equivalence and its independence from the choice of $P$ use
\cref{lem:purification-invariance}; the middle equivalence is
\cref{thm:finite-bo-certificates}, whose proof invokes the translations in
\cref{lem:intrinsic-purification-tests}; and the last equivalence is
\cref{cor:bo-product-identities}.

For fixed $k$, let $\InstMonoid{k}$ denote the monoid of finite $k$-labelled
$\cF$-instances and let $\zvec I$ denote the pinned partition vector of
$I\in\InstMonoid{k}$.  Precise definitions appear in
\cref{subsec:labelled-instances}.  Writing $W_{\cF}^{(k)}$ for the $k$-ary
part of the generated family, the global test is
\begin{equation*}
\begin{aligned}
  W_{\cF}^{(k)}\subseteq\mathcal B_k
  &\quad\Longleftrightarrow\quad
  \begin{gathered}
    q(G)=0\text{ for every }G\in W_{\cF}^{(k)},\\[-2pt]
    \text{and for every }q\in Q_k
  \end{gathered}\\
  &\quad\Longleftrightarrow\quad
  \begin{gathered}
    q(\zvec I)=0\text{ for every }I\in\InstMonoid{k},\\[-2pt]
    \text{and for every }q\in Q_k.
  \end{gathered}
\end{aligned}
\end{equation*}
The second equivalence is the exact parametrization
$W_{\cF}^{(k)}=\{\zvec I:I\in\InstMonoid{k}\}$ proved in
\cref{lem:pinned-realization}.  Each instance supplies one joint coordinate
assignment
\[
  X_{\mathbf a}=\zvec I(\mathbf a)
  \qquad(\mathbf a\in D^k).
\]
Hence the oracle in \cref{thm:universal-identity} tests vanishing on this
parametrized subset of $K^{D^k}$.

Here $\zvec I$ is the original pinned partition vector, and the polynomials in
$Q_k$ use its unbarred coordinates.  Thus the oracle is applied directly to
the generated table $G=\zvec I$.

\section{Universal Identities on Generated Tables}
\label{sec:universal-identities}

We continue to use the field $K$ from~\eqref{eq:working-field}.

\subsection{Labelled instances and pinned partition vectors}
\label{subsec:labelled-instances}

Fix $k\ge1$.  Let $\InstMonoid{k}$ be the commutative monoid of
label-preserving isomorphism classes of finite $k$-labelled $\cF$-instances.
The $k$ labels are carried by distinct variables.  Multiplication takes the
disjoint union of two instances and identifies variables with the same
label; the unit consists of the $k$ labelled variables and no constraints.
Unlabelled isolated variables are permitted.  For $I\in\InstMonoid{k}$,
define its pinned partition vector $\zvec I\in K^{D^k}$ by
\begin{equation}
  \zvec I(\mathbf a)
  =\sum_{\substack{\sigma:V(I)\to D\\
                    \sigma\text{ assigns }\mathbf a\text{ to the labels}}}
       \prod_{\substack{(f;v_1,\ldots,v_{r_f})\\
                         \in\mathcal C_I}}
          f\bigl(\sigma(v_1),\ldots,\sigma(v_{r_f})\bigr).
  \label{eq:pinned-partition-vector}
\end{equation}

\begin{lemma}[Pinned realization of the generated family]
\label{lem:pinned-realization}
Let $W_{\cF}^{(k)}$ be the set of $k$-ary members of $W_{\cF}$.  Then
\[
  \{\zvec I:I\in\InstMonoid{k}\}=W_{\cF}^{(k)}.
\]
Moreover, if $IJ$ denotes the labelled product, then
\[
  \zvec{IJ}=\zvec I\odot\zvec J,
\]
where $\odot$ is coordinatewise multiplication.
\end{lemma}

\begin{proof}
Given $G=F_I^{[k]}\in W_{\cF}^{(k)}$, label the retained variables of the
presenting instance.  Equation~\eqref{eq:pinned-partition-vector} is exactly
$G$.  Conversely, order the labelled variables of a labelled instance first
and sum all unlabelled variables; this is a member of $W_{\cF}^{(k)}$.
The correspondence leaves repeated scopes, arbitrary sharing, isolated
variables, previous boundary identifications, and exact cancellation
unchanged.  In the labelled product, the unlabelled variables of the two
factors remain independent after the common labels are pinned, which gives
the coordinatewise product formula.
\end{proof}

\subsection{Tensor characters and pinned indistinguishability}

We recall the consequence of Young's constraint-function isomorphism theorem
needed below~\cite[Definitions~3--7, Section~4.1, and
Corollary~35]{Young2025}.  A constraint-function
family $\mathcal H$ is defined over a common finite domain, denoted
$V(\mathcal H)$.  Two finite families $\mathcal H$ and $\mathcal G$ are
\emph{similar} if they are equipped with an arity-preserving bijection between
their functions.  A domain bijection $\sigma:V(\mathcal H)\to V(\mathcal G)$
is an \emph{isomorphism} of similar families if, for every corresponding pair
$h\leftrightarrow g$ of arity $r$,
\[
  h(a_1,\ldots,a_r)=g(\sigma(a_1),\ldots,\sigma(a_r))
  \qquad(a_1,\ldots,a_r\in V(\mathcal H)).
\]
Elements $u,v\in V(\mathcal H)$ are \emph{twins in $\mathcal H$} if, for
every $h\in\mathcal H$ of arity $r$, every $j\in[r]$, and every
$(a_1,\ldots,a_{r-1})\in V(\mathcal H)^{r-1}$,
\begin{multline*}
  h(a_1,\ldots,a_{j-1},u,a_j,\ldots,a_{r-1})\\
  =h(a_1,\ldots,a_{j-1},v,a_j,\ldots,a_{r-1}).
\end{multline*}
A labelled $\mathcal H$-instance is \emph{simple} if the
variables in each constraint scope are pairwise distinct and, for every fixed
function $h\in\mathcal H$, no two occurrences of $h$ have scope tuples that
are permutations of one another.

For a $k$-labelled $\mathcal H$-instance $I$ and a pin
$\beta:[k]\to V(\mathcal H)$, write $Z_{\beta}^{\mathcal H}(I)$ for its
pinned partition value.  If $\mathcal H$ and $\mathcal G$ are similar, let
$I_{\mathcal H\to\mathcal G}$ be the $\mathcal G$-instance obtained by
replacing every function of $\mathcal H$ by its corresponding function of
$\mathcal G$, while retaining the variables, labels, and constraint scopes.
Under this notation, the cited result gives the following pinned form.

\begin{theorem}[Pinned constraint-function isomorphism]
\label{thm:young-pinned}
Let $\mathcal H$ and $\mathcal G$ be similar finite sets of constraint
functions over a field of characteristic zero.  Fix $k\ge1$ and pins
$\beta:[k]\to V(\mathcal H)$ and $\gamma:[k]\to V(\mathcal G)$.
If
\begin{equation*}
  Z_{\beta}^{\mathcal H}(I)
  =Z_{\gamma}^{\mathcal G}(I_{\mathcal H\to\mathcal G})
\end{equation*}
for every simple $k$-labelled $\mathcal H$-instance $I$, then the two
domains have the same size and there is an isomorphism
$\sigma:V(\mathcal H)\to V(\mathcal G)$ of the two families such that
$\sigma(\beta(i))$ and $\gamma(i)$ are twins in $\mathcal G$ for every
$i\in[k]$.

Conversely, such a domain bijection and twin replacements preserve pinned
partition values on every $k$-labelled instance, whether or not it is simple.
\end{theorem}

The theorem is applied at a fixed triple $(k,\beta,\gamma)$: equality is
quantified over all simple instances for these fixed data, and the resulting
isomorphism may depend on the triple.

For the converse, an isomorphism is a change of summation variables, and
replacing a pin by a twin preserves every constraint value.  In our exactly
represented algebraic-number setting, the stated condition is decidable by
enumerating the bijections between the finite domains $V(\mathcal H)$ and
$V(\mathcal G)$ and testing the corresponding function tables and twin
relations exactly.

Let $p,q\ge0$.  For $f:D^{r_f}\to K$, define the function
$T_{p,q}(f):(D^{p+q})^{r_f}\to K$ by
\begin{align*}
  T_{p,q}(f)(\mathbf u_1,\ldots,\mathbf u_{r_f})
  &={}
    \prod_{\ell=1}^{p}
      f(u_{1,\ell},\ldots,u_{r_f,\ell})\\
  &\quad\cdot
    \prod_{\ell=p+1}^{p+q}
      \overline{f(u_{1,\ell},\ldots,u_{r_f,\ell})},
\end{align*}
where $\mathbf u_j\in D^{p+q}$.  When $p=q=0$, the domain is the
singleton $D^0$ and $T_{0,0}(f)$ is the constant-one function of arity
$r_f$ on that domain.  Write
\[
  T_{p,q}(\cF)=\bigl(T_{p,q}(f):f\in\cF\bigr)
\]
for this $f$-indexed tensorized family; distinct indices may initially give
the same function table.

Fix $k\ge1$.  For each boundary assignment
$\mathbf a=(a_1,\ldots,a_k)\in D^k$, introduce two independent formal
variables $X_{\mathbf a}$ and $Y_{\mathbf a}$.  When a polynomial in these
variables is evaluated at the pinned partition vector of an instance $I$,
the substitutions are
\begin{equation}
  X_{\mathbf a}\longmapsto\zvec I(\mathbf a),
  \qquad
  Y_{\mathbf a}\longmapsto\overline{\zvec I(\mathbf a)}.
  \label{eq:formal-coordinate-substitution}
\end{equation}
A monomial of $X$-degree $p$ and $Y$-degree $q$ can be written as
\begin{equation*}
  M=\prod_{i=1}^{p}X_{\mathbf a^{(i)}}
    \prod_{j=1}^{q}Y_{\mathbf b^{(j)}}.
\end{equation*}
For $I\in\InstMonoid{k}$, define
\begin{equation*}
  \chi_M(I):=M(\zvec I,\overline{\zvec I}),
\end{equation*}
using the simultaneous substitution in
\eqref{eq:formal-coordinate-substitution}.  The next lemma proves that
$\chi_M$ is a unital multiplicative character.
Here the possibly repeated indices
\[
\begin{aligned}
  \mathbf a^{(i)}&=(a_1^{(i)},\ldots,a_k^{(i)})\in D^k
    &&(i\in[p]),\\
  \mathbf b^{(j)}&=(b_1^{(j)},\ldots,b_k^{(j)})\in D^k
    &&(j\in[q])
\end{aligned}
\]
list the factors of $M$.  The parenthesized superscript numbers a factor and
is not exponentiation; for example, $a_t^{(i)}$ is coordinate $t$ of the
tuple indexing the $i$th $X$-factor.

The pin $\beta_M:[k]\to D^{p+q}$ is obtained by transposing these $p+q$
boundary assignments.  For label $t\in[k]$, set
\[
  \beta_M(t)
  =\bigl(a_t^{(1)},\ldots,a_t^{(p)},
          b_t^{(1)},\ldots,b_t^{(q)}\bigr)\in D^{p+q}.
\]
Equivalently, regarding the values $\beta_M(t)$ as columns,
\begin{equation*}
  \bigl(\beta_M(1)\ \cdots\ \beta_M(k)\bigr)
  =
  \begin{pmatrix}
    a_1^{(1)}&\cdots&a_k^{(1)}\\
    \vdots&&\vdots\\
    a_1^{(p)}&\cdots&a_k^{(p)}\\
    b_1^{(1)}&\cdots&b_k^{(1)}\\
    \vdots&&\vdots\\
    b_1^{(q)}&\cdots&b_k^{(q)}
  \end{pmatrix}.
\end{equation*}
The block of $\mathbf a$-rows or $\mathbf b$-rows is omitted when $p=0$ or
$q=0$, respectively.
For example, let $D=\{0,1\}$, $k=2$, $p=2$, $q=1$, and
\[
  M=X_{(0,1)}X_{(1,1)}Y_{(1,0)}.
\]
Then the three row assignments are $(0,1),(1,1),(1,0)$, so
\[
  \beta_M(1)=(0,1,1),
  \qquad
  \beta_M(2)=(1,1,0).
\]

\Cref{fig:tensor-character-layers} visualizes this example and the
construction used in the next lemma.

\begin{figure}[!t]
  \centering
  \resizebox{0.80\linewidth}{!}{%
  \begin{tikzpicture}[
      x=1cm,y=1.04cm,
      title/.style={font=\small\bfseries},
      panel/.style={draw,rounded corners=2pt,line width=.45pt},
      cell/.style={draw,minimum width=7mm,minimum height=5.5mm,
                   inner sep=0pt,font=\scriptsize},
      conjcell/.style={cell,double,double distance=.6pt},
      rowlabel/.style={font=\scriptsize,anchor=east},
      inst/.style={draw,rounded corners=3pt,minimum width=4.5cm,
                   minimum height=1.55cm,align=center,font=\scriptsize},
      pin/.style={draw,circle,minimum size=6.5mm,inner sep=1pt,
                  font=\scriptsize},
      layer/.style={draw,rounded corners=2pt,minimum width=4.35cm,
                    minimum height=6.5mm,align=left,font=\scriptsize},
      conjlayer/.style={layer,double,double distance=.6pt},
      arrow/.style={-{Latex[length=2mm]},line width=.55pt}
    ]
    % Panel (a): rows indexing the factors of the monomial.
    \draw[panel] (.35,7.85) rectangle (13.65,11.40);
    \node[title] at (7.00,11.05) {(a) factors of $M$ as rows};
    \node[font=\scriptsize] at (7.15,10.62) {$t=1$};
    \node[font=\scriptsize] at (7.95,10.62) {$t=2$};
    \node[rowlabel] at (6.68,10.19) {$X_{\mathbf a^{(1)}}$};
    \node[cell] at (7.15,10.19) {$0$};
    \node[cell] at (7.95,10.19) {$1$};
    \node[rowlabel] at (6.68,9.49) {$X_{\mathbf a^{(2)}}$};
    \node[cell] at (7.15,9.49) {$1$};
    \node[cell] at (7.95,9.49) {$1$};
    \node[rowlabel] at (6.68,8.79) {$Y_{\mathbf b^{(1)}}$};
    \node[conjcell] at (7.15,8.79) {$1$};
    \node[conjcell] at (7.95,8.79) {$0$};
    \node[font=\scriptsize,align=center] at (7.00,8.23)
      {single border: ordinary factor; double border: conjugated factor};

    \draw[arrow] (7.00,7.68)--(7.00,7.05)
      node[midway,right=3pt,font=\scriptsize]
      {transpose rows into $D^3$-valued columns};

    % Panel (b): bundle the columns into D^3-valued pins.
    \draw[panel] (.35,3.50) rectangle (13.65,6.85);
    \node[title] at (7.00,6.50) {(b) transpose and bundle columns};
    \node[inst] (IM) at (7.00,5.55)
      {$I^M$ has the same variables, labels, and scopes as $I$\\[2pt]
       every occurrence $f$ is replaced by $T_{2,1}(f)$};
    \node[pin] (p1) at (5.55,4.15) {$1$};
    \node[pin] (p2) at (8.45,4.15) {$2$};
    \draw[line width=.45pt] (p1.north |- IM.south)--(p1.north);
    \draw[line width=.45pt] (p2.north |- IM.south)--(p2.north);
    \node[font=\scriptsize,align=center] at (3.55,4.15)
      {$\beta_M(1)=(0,1,1)$};
    \node[font=\scriptsize,align=center] at (10.45,4.15)
      {$\beta_M(2)=(1,1,0)$};

    \draw[arrow] (7.00,3.32)--(7.00,2.89)
      node[midway,right=3pt,font=\scriptsize]
      {split each $D^3$-assignment into its three components};

    % Panel (c): component layers of a D^3-valued assignment.
    \draw[panel] (.35,-.85) rectangle (13.65,2.68);
    \node[title] at (7.00,2.33) {(c) independent component layers};
    \node[layer,minimum width=8.0cm] (L1) at (7.00,1.75)
      {layer $1$: pin $(0,1)$; value $\zvec I(0,1)$};
    \node[layer,minimum width=8.0cm] (L2) at (7.00,.92)
      {layer $2$: pin $(1,1)$; value $\zvec I(1,1)$};
    \node[conjlayer,minimum width=8.0cm] (L3) at (7.00,.09)
      {layer $3$: pin $(1,0)$; value $\overline{\zvec I(1,0)}$};
    \node[font=\scriptsize,align=center] at (7.00,-.49)
      {multiply the three layer sums to obtain $\chi_M(I)$};
  \end{tikzpicture}%
  }
  \caption{Tensor realization for the example preceding
  \cref{lem:tensor-character}.  The rows indexing the three monomial factors
  are transposed into one $D^3$-valued pin at each label.  An assignment to
  $I^M$ is then exactly three independent $D$-valued assignments to the same
  instance template $I$; double borders mark the factor and layer that use
  conjugated constraint weights.
  Hence its pinned partition value factors as
  $\zvec I(0,1)\zvec I(1,1)\overline{\zvec I(1,0)}=\chi_M(I)$.}
  \label{fig:tensor-character-layers}
\end{figure}

\begin{lemma}[Tensor realization of monomial characters]
\label{lem:tensor-character}
For $I\in\InstMonoid{k}$, let $I^M$ be the $k$-labelled
$T_{p,q}(\cF)$-instance obtained by replacing each occurrence of
$f\in\cF$ by $T_{p,q}(f)$ while retaining all variables, labels, and scopes.
Then
\begin{equation*}
  Z_{\beta_M}^{T_{p,q}(\cF)}(I^M)
  =\chi_M(I)
  =\prod_{i=1}^{p}\zvec I(\mathbf a^{(i)})
    \prod_{j=1}^{q}\overline{\zvec I(\mathbf b^{(j)})}.
\end{equation*}
Thus $\chi_M:\InstMonoid{k}\to K$ is a unital multiplicative character.
\end{lemma}

\begin{proof}
An assignment $\eta:V(I)\to D^{p+q}$ is uniquely a tuple
$(\sigma_1,\ldots,\sigma_{p+q})$ of assignments
$\sigma_\ell:V(I)\to D$.  By the definition of $\beta_M$, the assignment
$\sigma_i$ pins the $k$ labels to $\mathbf a^{(i)}$ for $i\in[p]$, while
$\sigma_{p+j}$ pins them to $\mathbf b^{(j)}$ for $j\in[q]$.
For an assignment $\sigma:V(I)\to D$, abbreviate its constraint product by
\[
  w_I(\sigma)
  =\prod_{(f;v_1,\ldots,v_{r_f})\in\mathcal C_I}
     f\bigl(\sigma(v_1),\ldots,\sigma(v_{r_f})\bigr).
\]
Expanding the tensorized constraint functions and using this abbreviation
gives
\begin{align*}
  Z_{\beta_M}^{T_{p,q}(\cF)}(I^M)
  &=\sum_{\sigma_1,\ldots,\sigma_{p+q}}
      \prod_{i=1}^p w_I(\sigma_i)
      \prod_{j=1}^q\overline{w_I(\sigma_{p+j})}\\
  &=\prod_{i=1}^{p}\zvec I(\mathbf a^{(i)})
    \prod_{j=1}^{q}\overline{\zvec I(\mathbf b^{(j)})}.
\end{align*}
The first sum ranges over the component assignments with the pins just
specified.  The second equality holds because these component assignments
are independent, so the finite sum factors; finite sums commute with complex
conjugation.  When $p=q=0$, the domain is $D^0$ and the value is the empty
product $1$.  Finally, \cref{lem:pinned-realization} gives
$\chi_M(IJ)=\chi_M(I)\chi_M(J)$ and $\chi_M(1)=1$, so $\chi_M$ is a unital
multiplicative character.
\end{proof}

Tensorization may collapse distinct input functions.  For example, if
$g=\mathrm{i}f$, where $\mathrm{i}^2=-1$, then
$T_{1,1}(g)=T_{1,1}(f)$.  Common scalar tags restore the canonical indexing
required by \cref{thm:young-pinned}; the next lemma constructs them.

\begin{lemma}[Common scalar tags]
\label{lem:scalar-tags}
Let $(h_f:f\in\cF)$ and $(g_f:f\in\cF)$ be indexed families over finite
domains $A$ and $B$, respectively, where
$h_f:A^{r_f}\to\AlgNums$ and $g_f:B^{r_f}\to\AlgNums$ are given by exact
tables.  Suppose that, within
each family and each arity, at most one indexed table is identically zero.
One can effectively choose positive rational numbers $\lambda_f$, $f\in\cF$,
such that, within every arity, both
\[
  \{\lambda_f h_f:f\in\cF\}
  \quad\text{and}\quad
  \{\lambda_f g_f:f\in\cF\}
\]
contain no duplicate tables.
\end{lemma}

\begin{proof}
Choose the tags sequentially.  For a previously treated function $e$ of the
same arity, the equality
$\lambda h_f=\lambda_e h_e$ excludes at most one positive rational value of
$\lambda$; the all-zero case cannot occur for two distinct indices by
hypothesis.  The same statement holds for the $g$-family.  Thus only finitely
many rational values are excluded at each stage, and enumerating the positive
integers finds a valid tag.
\end{proof}

The tensorized families associated with monomial characters satisfy the
hypothesis of \cref{lem:scalar-tags}.  At positive tensor degree,
$T_{p,q}(f)$ is zero exactly when $f$ is zero, and exact deduplication of the
input leaves at most one zero function of each arity.  At degree zero every
function table is the nonzero constant-one table.

\begin{lemma}[Effective comparison of monomial characters]
\label{lem:character-comparison}
Given two monomials in the coordinates
$\{X_{\mathbf a},Y_{\mathbf a}:\mathbf a\in D^k\}$, it is decidable whether
their associated characters agree on every element of $\InstMonoid{k}$.
\end{lemma}

\begin{proof}
Let $M$ have $X$-degree $p$ and $Y$-degree $q$, and let $N$ have
$X$-degree $r$ and $Y$-degree $s$.  The decision has three stages: represent
the two monomial characters by tensorized constraint-function families,
repair possible collisions in those families with common scalar tags, and
then apply the finite isomorphism-and-twin criterion in
\cref{thm:young-pinned}.

\paragraph{Tensorized representations.}
Write
\[
\begin{aligned}
  \mathcal H_M=(h_f:f\in\cF)
    &=\bigl(T_{p,q}(f):f\in\cF\bigr),\\
  \mathcal H_N=(g_f:f\in\cF)
    &=\bigl(T_{r,s}(f):f\in\cF\bigr).
\end{aligned}
\]
These families are defined over
\[
  A_M=D^{p+q},
  \qquad
  A_N=D^{r+s},
\]
and carry the pins $\beta_M:[k]\to A_M$ and
$\beta_N:[k]\to A_N$.  For an original labelled instance template
$I\in\InstMonoid{k}$, let $I^M$ and $I^N$ be its interpretations by the two
families, as in \cref{lem:tensor-character}.  That lemma gives
\begin{equation}
  Z_{\beta_M}^{\mathcal H_M}(I^M)=\chi_M(I),
  \qquad
  Z_{\beta_N}^{\mathcal H_N}(I^N)=\chi_N(I).
  \label{eq:two-monomial-realizations}
\end{equation}
Thus $A_M$ and $A_N$ are domains for the auxiliary constraint functions;
the arguments of both characters remain the original templates
$I\in\InstMonoid{k}$.

\paragraph{Common scalar tags.}
Tensorization may make two differently indexed functions equal.  Apply
\cref{lem:scalar-tags} to choose common positive rational tags $\lambda_f$
that make both indexed families duplicate-free, and put
\[
\begin{aligned}
  \widehat h_f&=\lambda_fh_f,
  &\widehat g_f&=\lambda_fg_f,\\
  \widehat{\mathcal H}_M&=(\widehat h_f:f\in\cF),
  &\widehat{\mathcal H}_N&=(\widehat g_f:f\in\cF).
\end{aligned}
\]
Both tagged families are duplicate-free within each arity.  If $N_f(I)$ is
the number of occurrences of $f$ in the original template $I$, interpreting
that template with the tagged tables multiplies both values in
\eqref{eq:two-monomial-realizations} by the same factor
\begin{equation*}
  \tau(I)=\prod_{f\in\cF}\lambda_f^{N_f(I)}.
\end{equation*}
This factor is nonzero.  Denoting the two tagged interpretations by
$\widehat I^M$ and $\widehat I^N$, define
\begin{align*}
  \widehat Z_M(I)
  &:=Z_{\beta_M}^{\widehat{\mathcal H}_M}(\widehat I^M)
    =\tau(I)\chi_M(I),\\
  \widehat Z_N(I)
  &:=Z_{\beta_N}^{\widehat{\mathcal H}_N}(\widehat I^N)
    =\tau(I)\chi_N(I).
\end{align*}
Since the same nonzero factor occurs on both sides, we have
\begin{equation}
\begin{aligned}
  &\chi_M(I)=\chi_N(I)
    \text{ for every }I\in\InstMonoid{k}\\
  &\qquad\Longleftrightarrow\qquad
    \widehat Z_M(I)=\widehat Z_N(I)
    \text{ for every }I\in\InstMonoid{k}.
\end{aligned}
  \label{eq:character-agreement-tagged-agreement}
\end{equation}

\paragraph{Finite isomorphism-and-twin criterion.}
The tagged families are genuine sets of size $\abs{\cF}$, and
\[
  \widehat h_f\longleftrightarrow\widehat g_f
  \qquad(f\in\cF)
\]
is an arity-preserving bijection.  They are therefore similar.  Because the
tables are duplicate-free, every labelled
$\widehat{\mathcal H}_M$-instance corresponds uniquely to an original
$\cF$-instance template, and replacing $\widehat h_f$ by
$\widehat g_f$ produces its $\widehat{\mathcal H}_N$ interpretation.

If the right-hand side of
\eqref{eq:character-agreement-tagged-agreement} holds, it holds in
particular for every simple instance.  By \cref{thm:young-pinned}, the
domains have equal size and there is a bijection
\begin{equation}
  \sigma:A_M\to A_N
  \label{eq:character-comparison-domain-bijection}
\end{equation}
that is an isomorphism of the two tagged families and satisfies
\begin{equation}
  \sigma(\beta_M(t))
  \text{ and }\beta_N(t)
  \text{ are twins in }\widehat{\mathcal H}_N
  \qquad(t\in[k]).
  \label{eq:character-comparison-pin-twins}
\end{equation}
Conversely, such a bijection maps every assignment of
$\widehat I^M$ to an assignment of $\widehat I^N$, preserves each
constraint weight, and preserves the pins through the twin conditions.
The converse direction of \cref{thm:young-pinned} therefore gives equality of
the two tagged pinned values
for every $I$.  Since $\tau(I)\ne0$, division by the common factor gives
$\chi_M(I)=\chi_N(I)$ for every $I\in\InstMonoid{k}$.

The resulting decision is finite.  If $\abs{A_M}\ne\abs{A_N}$, no
bijection in \eqref{eq:character-comparison-domain-bijection} exists and
the characters are different.  Otherwise enumerate every bijection
$\sigma:A_M\to A_N$.  For each one, test exactly whether all corresponding
tagged function tables are carried to one another and whether
\eqref{eq:character-comparison-pin-twins} holds.  The characters agree if
and only if some bijection passes.  All domains and tables are finite, so
the enumeration terminates.
\end{proof}

\subsection{The universal identity oracle}

We use the elementary linear independence of distinct monoid characters.

\begin{lemma}[Linear independence of monoid characters]
\label{lem:dedekind-independence}
Let $M$ be a monoid, let $L$ be a field of characteristic zero, and let
\[
  \chi_1,\ldots,\chi_s:M\to L
\]
be pairwise distinct unital multiplicative characters.  If
$c_1,\ldots,c_s\in L$ satisfy
\begin{equation}
  \sum_{i=1}^s c_i\chi_i(x)=0
  \qquad\text{for every }x\in M,
  \label{eq:dedekind-assumed-relation}
\end{equation}
then $c_1=\cdots=c_s=0$.  Equivalently,
$\chi_1,\ldots,\chi_s$ are linearly independent in the vector space $L^M$.
\end{lemma}

\begin{proof}
Suppose the conclusion is false, and choose a nontrivial relation of the
form~\eqref{eq:dedekind-assumed-relation} involving the minimum possible
number $s$ of characters.  After deleting zero terms and relabelling, we may
assume
\begin{equation}
  c_i\ne0
  \qquad(i\in[s]).
  \label{eq:dedekind-nonzero-coefficients}
\end{equation}
If $s=1$, evaluation at the monoid identity gives
\[
  0=c_1\chi_1(1)=c_1,
\]
contrary to \eqref{eq:dedekind-nonzero-coefficients}.  Hence $s\ge2$.

Because $\chi_1$ and $\chi_s$ are distinct functions on $M$, choose and fix
$u\in M$ such that
\begin{equation}
  \chi_1(u)\ne\chi_s(u).
  \label{eq:dedekind-fixed-witness}
\end{equation}
For an arbitrary $v\in M$, evaluate the assumed relation at $uv$ and use
multiplicativity:
\begin{equation}
\begin{aligned}
  0
  &=\sum_{i=1}^s c_i\chi_i(uv)\\
  &=\sum_{i=1}^s c_i\chi_i(u)\chi_i(v).
\end{aligned}
  \label{eq:dedekind-relation-at-uv}
\end{equation}
Evaluation at $v$, followed by multiplication by the fixed scalar
$\chi_s(u)$, gives
\begin{equation}
  0
  =\chi_s(u)\sum_{i=1}^s c_i\chi_i(v)
  =\sum_{i=1}^s c_i\chi_s(u)\chi_i(v).
  \label{eq:dedekind-scaled-relation-at-v}
\end{equation}
Subtracting \eqref{eq:dedekind-scaled-relation-at-v} from
\eqref{eq:dedekind-relation-at-uv} yields
\begin{equation}
  0=\sum_{i=1}^{s-1}
      c_i\bigl(\chi_i(u)-\chi_s(u)\bigr)\chi_i(v)
  \qquad\text{for every }v\in M.
  \label{eq:dedekind-shorter-relation}
\end{equation}
Indeed, the term indexed by $s$ has coefficient zero.  Because $u$ is now
fixed, the quantities
\[
  d_i=c_i\bigl(\chi_i(u)-\chi_s(u)\bigr)
  \qquad(i\in[s-1])
\]
are fixed scalars in $L$, independent of $v$.  Moreover,
\[
  d_1=c_1\bigl(\chi_1(u)-\chi_s(u)\bigr)\ne0
\]
by \cref{eq:dedekind-nonzero-coefficients,eq:dedekind-fixed-witness}.
Consequently, \eqref{eq:dedekind-shorter-relation} is a nontrivial linear
relation involving at most $s-1$ of the characters, contradicting the
minimality of $s$.
\end{proof}

\begin{theorem}[Universal identity oracle]
\label{thm:universal-identity}
Fix $k\ge1$.  The input is an ordinary commutative polynomial
\[
  P\in\AlgNums[X_{\mathbf a},Y_{\mathbf a}:\mathbf a\in D^k],
\]
given as a finite list of monomials with exactly encoded algebraic-complex
coefficients.  The $2\abs{D}^k$ symbols
$X_{\mathbf a},Y_{\mathbf a}$ are independent formal variables; no
homogeneity assumption is imposed.  For $I\in\InstMonoid{k}$, the notation
$P(\zvec I,\overline{\zvec I})$ means the simultaneous substitution
\begin{equation*}
  X_{\mathbf a}=\zvec I(\mathbf a),
  \qquad
  Y_{\mathbf a}=\overline{\zvec I(\mathbf a)}
  \qquad(\mathbf a\in D^k),
\end{equation*}
as in \eqref{eq:formal-coordinate-substitution}.  There is a total exact
algorithm deciding whether
\begin{equation}
  P(\zvec I,\overline{\zvec I})=0
  \qquad\text{for every }I\in\InstMonoid{k}.
  \label{eq:universal-star-identity}
\end{equation}
\end{theorem}

\begin{proof}
Let $L$ be the number field generated over $K$ by the finitely many
coefficients of $P$.  It is effectively constructible from their exact
encodings.  Every monomial character takes values in $K\subseteq L$, so we
may regard all of them as $L$-valued characters.

Expand $P$ into finitely many monomials and combine identical monomials.  By
\cref{lem:character-comparison}, we can decide equality of every pair of the
resulting characters and partition them into equivalence classes.  Choosing
one representative $\chi_E$ for every class $E$ in the resulting finite set
$\mathcal E_{\mathrm{char}}$ gives, for every
$I\in\InstMonoid{k}$,
\[
  P(\zvec I,\overline{\zvec I})
  =\sum_{E\in\mathcal E_{\mathrm{char}}}
     \left(\sum_{M\in E}c_M\right)\chi_E(I),
\]
where $c_M\in L$ is the coefficient of $M$.  The characters $\chi_E$ are
pairwise distinct, so \cref{lem:dedekind-independence} makes them linearly
independent over $L$.  Therefore~\eqref{eq:universal-star-identity} holds if
and only if
\[
  \sum_{M\in E}c_M=0
  \qquad\text{for every }E\in\mathcal E_{\mathrm{char}}.
\]
These are finitely many exact algebraic-number equality tests.

All tensorized function tables, scalar-tag searches, domain bijections, table
isomorphisms, twin tests, and coefficient sums are finite and effective.  If
one class sum is nonzero, \cref{lem:dedekind-independence} says that the
displayed linear combination is not the zero function on $\InstMonoid{k}$.
Hence a
finite labelled instance $I$ satisfies
$P(\zvec I,\overline{\zvec I})\ne0$, and the algorithm rejects.
\end{proof}

\subsection{Degree-multiple generated families}
\label{subsec:degree-multiple-identities}

Fix $\delta\ge1$.  Let $\mathfrak M_{k,\delta}$ be the submonoid of
$\InstMonoid{k}$ consisting of the labelled instances in which every
variable has occurrence degree divisible by $\delta$.  It is a submonoid:
in a labelled product, the degrees of identified labels add, and zero is
divisible by $\delta$, so the constraint-free labelled unit belongs to it.
The proof of \cref{lem:pinned-realization} restricts verbatim to give
\begin{equation*}
  \{\zvec I:I\in\mathfrak M_{k,\delta}\}
  =\{G\in\WFdeg{\delta}:\arity(G)=k\}.
\end{equation*}

To compare monomial characters after restriction to
$\mathfrak M_{k,\delta}$, we use a root-of-unity filter that reduces the
problem to finitely many ordinary character comparisons on
$\InstMonoid{k}$.

\begin{lemma}[Degree-filtered comparison of monomial characters]
\label{lem:degree-filtered-character-comparison}
Fix $k,\delta\ge1$.  Given two monomials $M,N$ in the coordinates
$\{X_{\mathbf a},Y_{\mathbf a}:\mathbf a\in D^k\}$ and their associated
characters $\chi_M,\chi_N:\InstMonoid{k}\to K$, it is decidable whether
\begin{equation}
  \chi_M(I)=\chi_N(I)
  \qquad\text{for every }I\in\mathfrak M_{k,\delta}.
  \label{eq:degree-filtered-character-goal}
\end{equation}
\end{lemma}

\begin{proof}
We implement the comparison by tensorization, cyclic filtering, finite
character comparison, and the linear independence of distinct monoid
characters.

\paragraph{Tensorized representations.}
Let
\[
  \mathcal H_M=(h_f:f\in\cF)
  \quad\text{and}\quad
  \mathcal H_N=(g_f:f\in\cF)
\]
be the tensorized families associated with $M$ and $N$, defined over the
finite domains $A_M$ and $A_N$, with pins
$\beta_M:[k]\to A_M$ and $\beta_N:[k]\to A_N$, respectively.  For an
original labelled $\cF$-instance template $I\in\InstMonoid{k}$, let $I^M$
and $I^N$ denote its interpretations by these two families.  By
\cref{lem:tensor-character},
\begin{equation}
  Z_{\beta_M}^{\mathcal H_M}(I^M)=\chi_M(I),
  \qquad
  Z_{\beta_N}^{\mathcal H_N}(I^N)=\chi_N(I).
  \label{eq:degree-filter-monomial-realizations}
\end{equation}

\paragraph{Cyclic degree filter.}
Recall that
\[
  C_\delta=\mathbb Z/\delta\mathbb Z
  =\{0,1,\ldots,\delta-1\},
\]
with addition modulo $\delta$.  This is an auxiliary cyclic domain, separate
from the original CSP domain $D$, introduced solely to filter occurrence
degrees modulo $\delta$.  Let $L_\delta=K(\omega_\delta)$.  For an input
function $f\in\cF$ of arity
$r_f$, define the everywhere-nonzero function
\begin{equation*}
  u_f(s_1,\ldots,s_{r_f})
  =\omega_\delta^{s_1+\cdots+s_{r_f}}.
\end{equation*}
Augment $h_f$ and $g_f$ to functions over
$A_M\times C_\delta$ and $A_N\times C_\delta$, respectively, by
\begin{multline*}
  \widehat h_f\bigl((a_1,s_1),\ldots,(a_{r_f},s_{r_f})\bigr)\\
  =h_f(a_1,\ldots,a_{r_f})
     u_f(s_1,\ldots,s_{r_f}),
\end{multline*}
and define $\widehat g_f$ by the same formula with $g_f$ in place of $h_f$.
For $\mathbf s=(s_1,\ldots,s_k)\in C_\delta^k$, interpret an original
template $I$ using the $\widehat h_f$-family and pin label $t$ to
$(\beta_M(t),s_t)$.  Denote the resulting pinned value by
\[
  \rho^M_{\mathbf s}(I).
\]
Define $\rho^N_{\mathbf s}(I)$ analogously using the $\widehat g_f$-family
and the pins $(\beta_N(t),s_t)$.  Coordinatewise multiplication under the
labelled product shows that every
\[
  \rho^M_{\mathbf s},\rho^N_{\mathbf s}:
  \InstMonoid{k}\to L_\delta
\]
is a unital multiplicative character.

\paragraph{Formula for the filter.}
Fix $I\in\InstMonoid{k}$.  For an assignment
$\theta:V(I)\to C_\delta$ to the cyclic coordinate, the product of all
auxiliary factors is
\begin{equation}
\begin{aligned}
  &\prod_{(f;v_1,\ldots,v_{r_f})\in\mathcal C_I}
    u_f\bigl(\theta(v_1),\ldots,\theta(v_{r_f})\bigr)\\
  &\quad={}
    \omega_\delta^{
      \sum_{(f;v_1,\ldots,v_{r_f})\in\mathcal C_I}
      \sum_{j=1}^{r_f}\theta(v_j)}\\
  &\quad={}
    \omega_\delta^{
      \sum_{v\in V(I)}\deg_I(v)\theta(v)}.
  \label{eq:degree-filter-assignment-weight}
\end{aligned}
\end{equation}
Here $\deg_I(v)$ is the occurrence degree of $v$.  If a variable occurs
several times in one scope, it appears the same number of times in the first
sum in \eqref{eq:degree-filter-assignment-weight}; hence repeated scope
positions are counted correctly.

Summing $\rho^M_{\mathbf s}(I)$ over all cyclic label pins
$\mathbf s\in C_\delta^k$ is equivalent to summing $\theta$ over all cyclic
assignments to $V(I)$.  The tensorized coordinates and the cyclic coordinate
are independent, so
\cref{eq:degree-filter-monomial-realizations,eq:degree-filter-assignment-weight}
give the following identity, where $\mathbf 1_E$ denotes the indicator of an event
$E$:
\begin{equation}
\begin{aligned}
  \sum_{\mathbf s\in C_\delta^k}\rho^M_{\mathbf s}(I)
  &=\chi_M(I)
    \sum_{\theta:V(I)\to C_\delta}
      \omega_\delta^{\sum_v\deg_I(v)\theta(v)}\\
  &=\chi_M(I)
    \prod_{v\in V(I)}
      \left(
        \sum_{t\in C_\delta}
          \omega_\delta^{t\deg_I(v)}
      \right)\\
  &=\delta^{\abs{V(I)}}
    \mathbf 1_{\{I\in\mathfrak M_{k,\delta}\}}
    \chi_M(I).
  \label{eq:degree-filter-H}
\end{aligned}
\end{equation}
The last equality applies \eqref{eq:cyclic-character-filter} separately at
every variable.  The identical calculation for $N$ gives
\begin{equation}
  \sum_{\mathbf s\in C_\delta^k}\rho^N_{\mathbf s}(I)
  =\delta^{\abs{V(I)}}
    \mathbf 1_{\{I\in\mathfrak M_{k,\delta}\}}
    \chi_N(I).
  \label{eq:degree-filter-G}
\end{equation}

\paragraph{Reduction to two ordinary character sums.}
Define functions on the full monoid $\InstMonoid{k}$ by
\begin{equation}
  R_M=\sum_{\mathbf s\in C_\delta^k}\rho^M_{\mathbf s},
  \qquad
  R_N=\sum_{\mathbf s\in C_\delta^k}\rho^N_{\mathbf s}.
  \label{eq:degree-filter-character-sums}
\end{equation}
If $I\notin\mathfrak M_{k,\delta}$, both functions vanish at $I$ by
\cref{eq:degree-filter-H,eq:degree-filter-G}.  If
$I\in\mathfrak M_{k,\delta}$, the same equations give
\[
  R_M(I)=\delta^{\abs{V(I)}}\chi_M(I),
  \qquad
  R_N(I)=\delta^{\abs{V(I)}}\chi_N(I).
\]
The common factor is nonzero because $L_\delta$ has characteristic zero.
Consequently,
\begin{equation}
  \chi_M=\chi_N\text{ on }\mathfrak M_{k,\delta}
  \quad\Longleftrightarrow\quad
  R_M=R_N\text{ on }\InstMonoid{k}.
  \label{eq:degree-filter-reduction}
\end{equation}

\paragraph{Finite comparison and acceptance test.}
There are $2\delta^k$ ordinary characters in
\eqref{eq:degree-filter-character-sums}.  The scalar-tagged comparison from
the proof of \cref{lem:character-comparison} applies, for any fixed pair
of pins, to two finite arity-matched indexed families represented by exact
algebraic tables on finite domains and satisfying the zero-table hypothesis
of \cref{lem:scalar-tags}.  Positive tensor degree preserves the zero pattern
of the input family, while tensor degree zero produces the nonzero
constant-one table.  Thus each tensorized family has at most one zero table
of each arity.  Moreover, $u_f$ is everywhere nonzero, so $\widehat h_f$ is
zero exactly when $h_f$ is zero, and likewise for $\widehat g_f$.  Hence the
augmented families satisfy the hypothesis.  Common scalar tags make both
families duplicate-free, after which the comparison procedure of
\cref{lem:character-comparison} decides each pairwise equality.

Partition the $2\delta^k$ characters into equality classes, let
$\mathcal E_{\mathrm{char}}$
be the resulting finite set of classes, and write
$E\in\mathcal E_{\mathrm{char}}$ for one class.  Define its multiplicities
on the two sides by
\begin{equation*}
\begin{aligned}
  m_M(E)&=\abs{\{\mathbf s\in C_\delta^k:
                    \rho^M_{\mathbf s}\in E\}},\\
  m_N(E)&=\abs{\{\mathbf s\in C_\delta^k:
                    \rho^N_{\mathbf s}\in E\}}.
\end{aligned}
\end{equation*}
Choose one representative character $\rho_E$ from each
$E\in\mathcal E_{\mathrm{char}}$.
Then, as functions on $\InstMonoid{k}$,
\begin{equation*}
  R_M-R_N
  =\sum_{E\in\mathcal E_{\mathrm{char}}}
     \bigl(m_M(E)-m_N(E)\bigr)\rho_E.
\end{equation*}
The representatives $\rho_E$, $E\in\mathcal E_{\mathrm{char}}$, are
pairwise distinct.  By
\cref{lem:dedekind-independence},
\begin{equation}
\begin{aligned}
  R_M=R_N
  \quad\Longleftrightarrow\quad&
  m_M(E)=m_N(E)\\[-2pt]
  &\text{for every }E\in\mathcal E_{\mathrm{char}}.
\end{aligned}
  \label{eq:degree-filter-acceptance-test}
\end{equation}
Combining
\cref{eq:degree-filter-reduction,eq:degree-filter-acceptance-test}
gives a finite exact decision procedure
for \eqref{eq:degree-filtered-character-goal}.
\end{proof}

\begin{theorem}[Degree-multiple universal identity oracle]
\label{thm:degree-multiple-universal-identity}
For every $k,\delta\ge1$ and every polynomial
\[
  P\in\AlgNums[X_{\mathbf a},Y_{\mathbf a}:\mathbf a\in D^k],
\]
given as a finite list of monomials with exactly encoded algebraic-complex
coefficients, it is decidable whether
\[
  P(\zvec I,\overline{\zvec I})=0
  \qquad\text{for every }I\in\mathfrak M_{k,\delta}.
\]
\end{theorem}

\begin{proof}
Every monomial of $P$ is a unital multiplicative character after restriction
to the submonoid $\mathfrak M_{k,\delta}$.  By
\cref{lem:degree-filtered-character-comparison}, equality of any two such
restricted characters is decidable.  Partition the finitely many monomials
into restricted-character classes and sum their coefficients in the number
field generated by the input coefficients and $K$.  Applying
\cref{lem:dedekind-independence} to the monoid $\mathfrak M_{k,\delta}$ shows
that the identity
holds exactly when every class sum is zero.  All steps are finite and exact.
\end{proof}

\section{The Global Block Orthogonality Decision Procedure}
\label{sec:bo-algorithm}

Recall that, for every $k\ge2$,
\[
  \mathcal B_k
  =\left\{A\in K^{D^k}:
    \begin{array}{l}
      \text{a singleton purification of $A$}\\
      \text{is block-orthogonal}
    \end{array}
  \right\}.
\]
By \cref{cor:bo-product-identities}, one can effectively construct a finite
set
\[
  Q_k\subseteq K[X_{\mathbf a}:\mathbf a\in D^k]
\]
such that, for every $A\in K^{D^k}$,
\begin{equation}
  A\in\mathcal B_k
  \quad\Longleftrightarrow\quad
  q(A)=0\text{ for every }q\in Q_k.
  \label{eq:bo-algorithm-membership}
\end{equation}
When invoking \cref{thm:universal-identity}, we regard $q$ as a polynomial in
the full displayed ring that is independent of every variable
$Y_{\mathbf a}$.

\begin{procedure}[Global Block Orthogonality]
\label{proc:global-bo}
On input a nonempty finite domain $D$ and a finite algebraic-complex
constraint language $\cF$ whose values are given by exact encodings, perform
the following steps.
\begin{enumerate}
  \item Delete duplicate typed tables.  If $\cF=\varnothing$, return
        \YES.  Otherwise construct the field $K$
        in~\eqref{eq:working-field} and compute $\mu(K)$.
  \item For every integer
        \[
          2\le k\le d^2+1,
          \qquad d=\abs D,
        \]
        enumerate the Block Orthogonality certificates of arity $k$ and
        construct $Q_k$ as in \cref{cor:bo-product-identities}.
  \item For every tested $k$ and every $q\in Q_k$, use
        \cref{thm:universal-identity} to decide whether
        \[
          q(\zvec I)=0
          \qquad\text{for every }I\in\InstMonoid{k}.
        \]
        Return \NO{} as soon as one identity fails.
  \item If every identity at every indicated arity succeeds, return
        \YES.
\end{enumerate}
\end{procedure}

\begin{theorem}[Exact decision of global BO]
\label{thm:global-bo-decidable}
\Cref{proc:global-bo} is a total exact algorithm.  It returns \YES{}
on $(D,\cF)$ if and only if the Block Orthogonality condition from the
complex-weighted $\#\mathrm{CSP}$ dichotomy theorem of Cai and
Chen~\cite[Section~3.1]{CaiChen2017}, denoted $\BO(\cF)$ in
\cref{def:global-bo}, holds on the entire unbounded family $W_{\cF}$.
\end{theorem}

\begin{proof}
The empty-language case was verified in \cref{sec:preliminaries}; assume
$\cF\ne\varnothing$.

\paragraph{Termination.}
The arity loop is finite.  At each fixed arity, the row and column sets are
finite, $\mu(K)$ is finite, and the certificate construction ranges over
finitely many rectangle systems, row-pair modes, set partitions, anchors,
and root-of-unity labels.  Every root-sum test is an exact equality test in
$K$.  Each certificate has finitely many equations, and hence $Q_k$ is
finite and effectively computable.  A branch with no nonzero equations is
recognized syntactically and yields $Q_k=\varnothing$.

Finally, \cref{thm:universal-identity} is total and exact for each queried
polynomial.  Consequently every loop in \cref{proc:global-bo} is finite.

\paragraph{Soundness of \YES.}
Suppose the procedure returns \YES.  Fix
$2\le k\le d^2+1$ and a $k$-ary $G\in W_{\cF}$.  By
\cref{lem:pinned-realization}, there is a finite
$I\in\InstMonoid{k}$ such that $G=\zvec I$.  Every queried identity is
universal, so $q(G)=0$ for every $q\in Q_k$.  By
\eqref{eq:bo-algorithm-membership}, $G\in\mathcal B_k$.  Thus every generated
table through arity $d^2+1$ has a block-orthogonal singleton purification.

If a generated table of larger arity had a non-block-orthogonal singleton
purification, \cref{lem:bo-arity-bound} would produce such a member of
$W_{\cF}$ at a tested arity, a contradiction.  Hence every member of
$W_{\cF}$ of arity at least two has a block-orthogonal singleton
purification.  By \cref{cor:singleton-purification}, this is exactly the
original joint condition $\BO(\cF)$.

\paragraph{Completeness of \YES.}
Conversely, suppose $\BO(\cF)$ holds.  Every $G\in W_{\cF}$ of arity at
least two has a block-orthogonal singleton purification.  At every tested
arity and for every $I\in\InstMonoid{k}$, the table $\zvec I$ therefore lies
in $\mathcal B_k$.  Equation~\eqref{eq:bo-algorithm-membership} gives
\[
  q(\zvec I)=0
  \qquad
  \text{for every }q\in Q_k
  \text{ and every }I\in\InstMonoid{k}.
\]
The universal identity oracle accepts every query, and the procedure returns
\YES.

\paragraph{Soundness and completeness of \NO.}
If the procedure returns \NO, then some $q\in Q_k$ is not universal.
By \cref{thm:universal-identity}, a finite labelled instance
$I\in\InstMonoid{k}$ exists for which $q(\zvec I)\ne0$.
Equation~\eqref{eq:bo-algorithm-membership}
gives $\zvec I\notin\mathcal B_k$, so a generated table violates the
singleton condition.  By \cref{cor:singleton-purification}, $\BO(\cF)$ is
false.

Conversely, if $\BO(\cF)$ is false, then
\cref{cor:singleton-purification,lem:bo-arity-bound} give a
$G\in W_{\cF}$ of some tested arity $k$ with $G\notin\mathcal B_k$.  Write
$G=\zvec I$ using \cref{lem:pinned-realization}.  The reverse implication
in~\eqref{eq:bo-algorithm-membership} supplies $q\in Q_k$ with
$q(\zvec I)\ne0$.  That identity is not universal, so the procedure returns
\NO.

Because $\InstMonoid{k}$ ranges over all finite labelled instances and
$\zvec I$ records their exact pinned values, the argument covers every member
of $W_{\cF}$, including cancellation-created zeros; the bound $d^2+1$
concerns only retained arity.
\end{proof}

\section{Generated Supports and a Common Mal'tsev Operation}
\label{sec:generated-supports}

We now prove that global Block Orthogonality supplies one Mal'tsev operation
preserving every generated support.  Every auxiliary relation used below is
realized as the support of a member of $W_{\cF}$.

\subsection{Equality-free support realization}

We first isolate the cancellation argument used throughout this section.

\Needspace{8\baselineskip}
\begin{lemma}[Simultaneous nonvanishing of power sums]
\label{lem:simultaneous-power-sums}
Let $\mathcal A_1,\ldots,\mathcal A_t$ be nonempty finite multisets of
elements of $\AlgNums^\times$.  There exist infinitely many integers $m\ge1$
such that
\begin{equation}
  \sum_{a\in\mathcal A_i}a^m\ne0
  \qquad\text{for every }i\in[t].
  \label{eq:simultaneous-power-sums}
\end{equation}
Moreover, such an $m$ can be found by enumerating the positive integers and
performing exact algebraic zero tests.
\end{lemma}

\begin{proof}
For $i\in[t]$, put $s_i(n)=\sum_{a\in\mathcal A_i}a^n$.  This is a linear
recurrence over the number field generated by $\mathcal A_i$, with
characteristic roots among the distinct elements of $\mathcal A_i$.  By the
Skolem--Mahler--Lech theorem~\cite{Lech1953}, its zero set is the union of a
finite set and finitely many arithmetic progressions.

More explicitly, after absorbing finitely many initial terms into the finite
exceptional set, write each progression as
\[
  \{n\ge N:n\equiv r\pmod p\},
  \qquad 0\le r<p.
\]
Call such a progression a \emph{zero progression}; $p$ is its period and
$r$ its residue.  No zero progression can have residue $0$.  Otherwise, for
some $p\ge1$ and $N\ge0$,
\begin{equation}
  s_i(pn)=0
  \qquad(n\ge N).
  \label{eq:zero-progression-multiples}
\end{equation}
Group the elements of the multiset $\mathcal A_i$ according to their $p$th
powers.  Let $\beta_1,\ldots,\beta_h$ be the distinct resulting nonzero
values, and let $\nu_j\ge1$ be the multiplicity of $\beta_j$.  Then
\[
  s_i(pn)
  =\sum_{a\in\mathcal A_i}(a^p)^n
  =\sum_{j=1}^h\nu_j\beta_j^n.
\]
Applying~\eqref{eq:zero-progression-multiples} at
$n=N,N+1,\ldots,N+h-1$ gives
\begin{equation}
  \begin{pmatrix}
    1&\cdots&1\\
    \beta_1&\cdots&\beta_h\\
    \vdots&&\vdots\\
    \beta_1^{h-1}&\cdots&\beta_h^{h-1}
  \end{pmatrix}
  \begin{pmatrix}
    \nu_1\beta_1^N\\
    \vdots\\
    \nu_h\beta_h^N
  \end{pmatrix}
  =\mathbf0.
  \label{eq:zero-progression-vandermonde}
\end{equation}
The determinant of the first matrix is
\[
  \prod_{1\le j<k\le h}(\beta_k-\beta_j)\ne0,
\]
so~\eqref{eq:zero-progression-vandermonde} forces
$\nu_j\beta_j^N=0$ for every $j$.  Since every $\beta_j$ is nonzero, this
contradicts $\nu_j\ge1$.

Let $M$ be a common multiple of all progression periods occurring for all
$s_i$, taking $M=1$ if there are no such progressions.  Every multiple of
$M$ has residue zero modulo each such period and
therefore belongs to none of the zero progressions.  All sufficiently large
multiples of $M$ also avoid the finitely many exceptional zeros.  This proves
existence of infinitely many simultaneous choices.  It also proves
termination of the exact search that enumerates $m=1,2,\ldots$ and tests
\eqref{eq:simultaneous-power-sums}.
\end{proof}

Call a primitive-positive (pp) formula \emph{equality-free} if it uses
conjunction, existential quantification, and repeated occurrences of the same
logical variable, but has no equality atom between two distinct variables.

\begin{lemma}[Equality-free support realization]
\label{lem:equality-free-support-realization}
Let $h\ge1$ and $G_1,\ldots,G_h\in W_{\cF}$.  Let
$\mathbf x=(x_1,\ldots,x_\ell)$, with $\ell\ge1$, be the tuple of free
variables, and let $\mathbf y=(y_1,\ldots,y_s)$, with $s\ge0$, be the tuple
of existential variables.  Let $\mathcal E$ be a finite set indexing the
atoms.  For every $e\in\mathcal E$, choose
$i(e)\in[h]$ and an ordered scope $\mathbf z_e$ of length
$\arity(G_{i(e)})$ whose entries are variables from $\mathbf x$ and
$\mathbf y$.  The same logical variable may occur repeatedly within one
scope or across several scopes.  Define $R\subseteq D^\ell$ by
\begin{equation}
  R(\mathbf x)
  \quad\Longleftrightarrow\quad
  \exists\mathbf y\;
  \bigwedge_{e\in\mathcal E}
       \SuppRel{G_{i(e)}}(\mathbf z_e).
  \label{eq:equality-free-pp}
\end{equation}
Then there exists $H\in W_{\cF}$ with $\SuppRel H=R$.  Given presentations
of the $G_i$, such an $H$ can be obtained by a terminating exact search.
\end{lemma}

\paragraph{Interpretation and example.}
Because $\SuppRel G$ records precisely the nonzero entries of $G$, the direct
set-theoretic meaning of~\eqref{eq:equality-free-pp} is
\begin{equation*}
\begin{aligned}
  R=\{\mathbf a\in D^\ell:\;&
      \text{there exists }\mathbf b\in D^s\text{ such that}\\
    &G_{i(e)}\bigl(\mathbf z_e(\mathbf a,\mathbf b)\bigr)\ne0
      \quad\text{for every }e\in\mathcal E\}.
\end{aligned}
\end{equation*}
Thus $R$ consists of the assignments to the free variables that extend to an
assignment of the hidden variables making every selected generated table
nonzero on its indicated scope.  Equivalently, $R$ is obtained by projecting
the solution set of these support constraints onto the $\mathbf x$
coordinates.

For example, if $G_1$ and $G_2$ are binary, then
\[
  R(x_1,x_2)
  \quad\Longleftrightarrow\quad
  \exists y\;
  \SuppRel{G_1}(x_1,y)\wedge\SuppRel{G_2}(y,x_2)
\]
means directly that
\[
\begin{aligned}
  R=\{(a_1,a_2)\in D^2:\;&\exists b\in D,\\
       &G_1(a_1,b)\ne0,\\[-2pt]
       &G_2(b,a_2)\ne0\}.
\end{aligned}
\]
Literal reuse of the same logical variable, such as the two occurrences of
$y$ in this example, is permitted and is distinct from an equality atom
$u=v$ joining two distinct logical variables.

\Cref{fig:equality-free-support-construction} shows how this binary example
is implemented by one actual base-language instance.

\begin{figure}[!t]
  \centering
  \resizebox{0.94\linewidth}{!}{%
  \begin{tikzpicture}[
      x=1cm,y=1cm,
      title/.style={font=\small\bfseries},
      panel/.style={draw,rounded corners=2pt,line width=.45pt},
      var/.style={circle,draw,minimum size=6.5mm,inner sep=0pt,
                  font=\scriptsize},
      hidden/.style={var,dashed},
      factor/.style={draw,double,double distance=.45pt,rounded corners=2pt,
                     minimum width=9mm,minimum height=6.5mm,
                     font=\scriptsize},
      module/.style={draw,rounded corners=2pt,minimum width=2.2cm,
                     minimum height=6.5mm,line width=.65pt,
                     font=\scriptsize,align=center},
      edge/.style={line width=.42pt},
      arrow/.style={-{Latex[length=2mm]},line width=.55pt}
    ]
    % Panel (a): the support pattern to be implemented.
    \draw[panel] (-.20,5.95) rectangle (14.20,9.25);
    \node[title] at (7.00,8.90) {(a) equality-free support pattern};
    \node[var] (sx1) at (2.30,7.65) {$x_1$};
    \node[factor] (sG1) at (4.40,7.65) {$G_1$};
    \node[hidden] (sy) at (7.00,7.65) {$y$};
    \node[factor] (sG2) at (9.60,7.65) {$G_2$};
    \node[var] (sx2) at (11.70,7.65) {$x_2$};
    \draw[edge] (sx1)--(sG1)--(sy)--(sG2)--(sx2);
    \node[font=\scriptsize,align=center] at (7.00,6.88)
      {$R(x_1,x_2)\iff\exists y\;$
       $\SuppRel{G_1}(x_1,y)\wedge\SuppRel{G_2}(y,x_2)$};
    \node[font=\scriptsize,align=center] at (7.00,6.30)
      {the dashed $y$ node is shared variable reuse, not an equality atom};

    \draw[arrow] (7.00,5.78)--(7.00,5.15)
      node[midway,right=3pt,font=\scriptsize,align=left]
      {replace each support atom by $m$ fresh copies\\
       of a fixed base-language presentation};

    % Panel (b): m copies of presentations of G_1 and G_2.
    \draw[panel] (-.20,-.55) rectangle (14.20,4.95);
    \node[title] at (7.00,4.60)
      {(b) one $\cF$-instance whose marginal is $H_m$};

    \node[module] (I11) at (5.00,3.90) {$I_1^{(1)}\;\bullet$};
    \node[module] (I12) at (5.00,3.05) {$I_1^{(2)}\;\bullet$};
    \node[font=\scriptsize] at (5.00,2.51) {$\vdots$};
    \node[module] (I1m) at (5.00,1.93) {$I_1^{(m)}\;\bullet$};

    \node[module] (I21) at (9.00,3.90) {$I_2^{(1)}\;\bullet$};
    \node[module] (I22) at (9.00,3.05) {$I_2^{(2)}\;\bullet$};
    \node[font=\scriptsize] at (9.00,2.51) {$\vdots$};
    \node[module] (I2m) at (9.00,1.93) {$I_2^{(m)}\;\bullet$};

    \node[var] (x1) at (2.15,1.13) {$x_1$};
    \node[hidden] (y) at (7.00,1.13) {$y$};
    \node[var] (x2) at (11.85,1.13) {$x_2$};

    \foreach \u in {I11,I12,I1m}{
      \draw[edge] (x1)--(\u.west);
      \draw[edge] (y)--(\u.east);
    }
    \foreach \u in {I21,I22,I2m}{
      \draw[edge] (y)--(\u.west);
      \draw[edge] (x2)--(\u.east);
    }

    \node[font=\scriptsize,align=center] at (7.00,-.07)
      {solid boxes are fresh copies of fixed $\cF$-presentations;
       each $\bullet$ denotes fresh internal variables\\
       retain $x_1,x_2$ and sum over the shared dashed $y$ and all
       $\bullet$-variables};
  \end{tikzpicture}%
  }
  \caption{The instance construction in
  \cref{lem:equality-free-support-realization} for the binary example.
  For each atom, take $m$ fresh copies of a base-language presentation of
  the corresponding generated table.  Identify only their displayed
  boundary variables with the logical variables $x_1,y,x_2$; all internal
  variables, denoted by $\bullet$, remain fresh.  Marginalizing the internal
  variables first gives
  $G_1(x_1,y)^mG_2(y,x_2)^m$, and then marginalizing the one shared variable
  $y$ gives $H_m(x_1,x_2)$.}
  \label{fig:equality-free-support-construction}
\end{figure}

\begin{proof}
For $m\ge1$, define
\begin{equation}
  H_m(\mathbf x)
  =\sum_{\mathbf y\in D^s}
       \prod_{e\in\mathcal E}G_{i(e)}(\mathbf z_e)^m.
  \label{eq:powered-support-table}
\end{equation}
The empty product is $1$.  For fixed $\mathbf a\in D^\ell$, let
$\mathcal A_{\mathbf a}$ be the multiset of all nonzero values
\[
  \prod_{e\in\mathcal E}
      G_{i(e)}(\mathbf z_e(\mathbf a,\mathbf b)),
  \qquad \mathbf b\in D^s.
\]
Then $\mathcal A_{\mathbf a}$ is nonempty exactly when $\mathbf a\in R$,
and
\begin{equation}
  H_m(\mathbf a)=\sum_{w\in\mathcal A_{\mathbf a}}w^m,
  \label{eq:witness-power-sum}
\end{equation}
where an empty sum is zero.  Thus tuples outside $R$ give zero term by term
for every $m$.  Apply \cref{lem:simultaneous-power-sums} to the finitely many
nonempty multisets $\mathcal A_{\mathbf a}$, $\mathbf a\in R$.  The resulting
common exponent makes~\eqref{eq:witness-power-sum} nonzero for every tuple in
$R$.  If $R=\varnothing$, any $m$ works.  Hence $\SuppRel{H_m}=R$.

The table $H_m$ belongs to $W_{\cF}$ by the following instance construction.
Choose a presenting base-language instance for each $G_i$.  For every
occurrence $e$ and every
copy index in $[m]$, take a fresh copy of that presentation, identify its
displayed boundary variables with those in $\mathbf z_e$, and keep all of
its hidden variables fresh.  For fixed $(\mathbf x,\mathbf y)$,
marginalizing the copied hidden variables factorizes as
\[
  \prod_{e\in\mathcal E}G_{i(e)}(\mathbf z_e)^m.
\]
Marginalizing $\mathbf y$ then gives~\eqref{eq:powered-support-table}.
Repeated variables in $\mathbf z_e$ are legal scope identifications, and the
retained variables may be put in the displayed order.  Therefore
$H_m\in W_{\cF}$.  Zeros created by earlier cancellation remain zero, while
the simultaneous choice of $m$ removes cancellation in the outer sum.
\end{proof}

\begin{corollary}[Singleton rectangularity]
\label{cor:generated-singleton-rectangularity}
Assume $\BO(\cF)$.  Let $R\subseteq D^{q-1}\times D$, with $q\ge2$, have an
equality-free pp definition from finitely many relations $\SuppRel G$,
$G\in W_{\cF}$.  For $\mathbf u\in D^{q-1}$, put
\[
  R_{\mathbf u}=\{a\in D:(\mathbf u,a)\in R\}.
\]
Then any two nonempty sets $R_{\mathbf u}$ and $R_{\mathbf v}$ are equal or
disjoint.
\end{corollary}

\begin{proof}
By \cref{lem:equality-free-support-realization}, in the displayed coordinate
order there is an $H\in W_{\cF}$ with $\SuppRel H=R$.  Apply $\BO(\cF)$ to
the singleton set $\{H\}$ and write $\Pure(H)=\widetilde H$.  Purification
preserves support, and $\abs{\widetilde H}$ is block-rank one.  Hence two
nonzero rows that overlap have equal support.
\end{proof}

\subsection{The universal polymorphism relation}

\begin{lemma}[Finite-family support lemma]
\label{lem:finite-support-maltsev}
Assume $\BO(\cF)$.  Every finite family of generated support relations has a
common Mal'tsev polymorphism.
\end{lemma}

\begin{proof}
Let $\Delta\subseteq\{\SuppRel G:G\in W_{\cF}\}$ be finite.  Empty
relations may be removed because every operation preserves them.  If nothing
remains, fix $a_0\in D$ and use
\[
  m_0(a,b,c)=
  \begin{cases}
    c,&a=b,\\
    a,&b=c,\\
    a_0,&\text{otherwise}.
  \end{cases}
\]
Thus assume that $\Delta$ contains a nonempty relation.  For every
$R\in\Delta$, fix $G_R\in W_{\cF}$ such that $R=\SuppRel{G_R}$, and write
$r_R$ for its arity.

For every $R\in\Delta$, every ordered triple
$\mathbf u^1,\mathbf u^2,\mathbf u^3\in R$, and every $j\in[r_R]$, record
the operation-table input $(u^1_j,u^2_j,u^3_j)\in D^3$.  Let
$T\subseteq D^3$ be the set of all recorded inputs, and introduce one
logical variable $x_\tau$, taking values in $D$, for every $\tau\in T$.
Its intended meaning is
\[
  x_\tau=t(\tau)
\]
for one candidate ternary operation $t:D^3\to D$.  Consequently, an input to
the formula below is an assignment
\[
  \mathbf p=(p_\tau)_{\tau\in T}\in D^T,
\]
equivalently a partial operation table $p:T\to D$; evaluation at $\mathbf p$
sets $x_\tau=p_\tau=p(\tau)$.

For every $R\in\Delta$ and
$\mathbf U=(\mathbf u^1,\mathbf u^2,\mathbf u^3)\in R^3$, use the formal
scope abbreviation
\[
  \mathbf x[\mathbf U]
  =\bigl(x_{(u^1_j,u^2_j,u^3_j)}\bigr)_{j\in[r_R]}.
\]
Define the finite constraint formula, viewed equivalently as a predicate on
$D^T$, by
\begin{align}
  \Psi_\Delta((x_\tau)_{\tau\in T})
  &=\bigwedge_{R\in\Delta}
    \ \bigwedge_{\mathbf U\in R^3}
    \SuppRel{G_R}\bigl(\mathbf x[\mathbf U]\bigr).
  \label{eq:universal-polymorphism-formula}
\end{align}
Directly, for an input partial table $p:T\to D$,
\begin{equation*}
\begin{aligned}
  \Psi_\Delta(\mathbf p)
  \quad\Longleftrightarrow\quad
  &\bigl(p(u^1_j,u^2_j,u^3_j)\bigr)_{j\in[r_R]}\in R\\
  &\text{for every }R\in\Delta\text{ and every}\\[-2pt]
  &\qquad \mathbf u^1,\mathbf u^2,\mathbf u^3\in R.
\end{aligned}
\end{equation*}
Thus $\Psi_\Delta$ checks all coordinatewise preservation tests for all
relations in $\Delta$ simultaneously.  For example, if $D=\{0,1\}$ and one
test for a binary relation $R$ uses
\[
  \mathbf u^1=(0,0),\qquad
  \mathbf u^2=(0,1),\qquad
  \mathbf u^3=(1,1),
\]
then its conjunct is
\[
  R\bigl(x_{(0,0,1)},x_{(0,1,1)}\bigr),
\]
which requires
$(p(0,0,1),p(0,1,1))\in R$ on an input table $p$.
Identical operation-table inputs use literally the same variable.  Thus
\eqref{eq:universal-polymorphism-formula} is equality-free, even when a
variable occurs repeatedly in one atomic scope.

Let $U_\Delta\subseteq D^T$ be its solution relation.  Its elements are
exactly the restrictions to $T$ of common ternary polymorphisms of $\Delta$.
Indeed, a common polymorphism satisfies every displayed atom.  Conversely,
extend a solution arbitrarily from $T$ to $D^3$.  Every coordinate triple in
a preservation test belongs to $T$, and its atom asserts that the
coordinatewise image belongs to the tested relation.  Hence every such
extension is a common polymorphism.

Put
\begin{equation*}
  \begin{aligned}
    A&=\{(a,a,b)\in T:a,b\in D\},\\
    C&=\{(a,b,b)\in T:a,b\in D,\ a\ne b\},
  \end{aligned}
\end{equation*}
and let $P=\pr_{A\cup C}(U_\Delta)$.  The sets $A$ and $C$ are disjoint,
and every constant triple in $T$ lies in $A$.  For each $k\in\{1,2,3\}$,
let
\[
  \pi_k:D^3\to D,
  \qquad
  \pi_k(a_1,a_2,a_3)=a_k,
\]
be the $k$th coordinate-selector operation.  It preserves every relation:
for $R\in\Delta$ and
$\mathbf u^1,\mathbf u^2,\mathbf u^3\in R$,
\[
  \pi_k(\mathbf u^1,\mathbf u^2,\mathbf u^3)
  =\bigl(\pi_k(u^1_j,u^2_j,u^3_j)\bigr)_{j\in[r_R]}
  =\mathbf u^k\in R.
\]
Hence
\begin{equation}
  \pi_k|_T\in U_\Delta,
  \qquad
  \pi_k|_{A\cup C}\in P
  \qquad(k=1,2,3).
  \label{eq:projection-tables-in-P}
\end{equation}
Here $\pi_k$ selects one of three operation arguments, whereas
$\pr_{A\cup C}$ is the different map that restricts a table indexed by $T$
to the coordinates in $A\cup C$.  Moreover,
$A\ne\varnothing$: repeating any tuple of a nonempty positive-arity relation
three times produces a constant triple in $T$.

Enumerate $C=\{c_1,\ldots,c_s\}$.  For $0\le j\le s$, we inductively
construct a partial table $p_j:A\cup C\to D$ satisfying the three separate
invariants
\begin{align}
  p_j&\in P,
    \notag\\
  p_j(a,a,b)&=b
    &&\text{for every }(a,a,b)\in A,
    \notag\\
  p_j(c_i)&=\pi_1(c_i)
    &&\text{for every }i\in[j].
    \label{eq:partial-maltsev-c}
\end{align}
Membership in $P$ is maintained as an independent invariant alongside the
two value conditions.  For
$j=0$, take $p_0=\pi_3|_{A\cup C}$.  Membership follows from
\eqref{eq:projection-tables-in-P}, the first value condition follows from
$\pi_3(a,a,b)=b$, and the second is vacuous.

Suppose the three invariants hold for $p_j$.  Because
\[
  p_j\in P=\pr_{A\cup C}(U_\Delta),
\]
there exists a lift $\widehat p_j\in U_\Delta$ satisfying
\begin{equation*}
  \widehat p_j|_{A\cup C}=p_j.
\end{equation*}
This lift is effective: enumerate all maps
$T\setminus(A\cup C)\to D$, combine each with the fixed values of $p_j$,
and test~\eqref{eq:universal-polymorphism-formula}.  There are at most
$d^{\abs{T\setminus(A\cup C)}}$ candidates, and membership in $P$ guarantees
that one succeeds.  Extend $\widehat p_j$ to a map $t_j:D^3\to D$ by assigning
an arbitrary domain value on $D^3\setminus T$.  By the definition of $T$,
those cells occur in no preservation test, so $t_j$ is a common ternary
polymorphism of $\Delta$.

The minor
\begin{equation*}
  q_j(x,y,z)=t_j(x,y,y)
\end{equation*}
is again a common ternary polymorphism, because identifying and repeating
arguments of a polymorphism preserves every relation it preserves.  Let
$K_j=A\cup\{c_1,\ldots,c_j\}$.  On $K_j$, $q_j$ agrees with $\pi_1$.
This follows from~\eqref{eq:partial-maltsev-c} on the previously fixed cells.
If $(x,x,z)\in A$, then $(x,x,x)\in T$ by repeating the tuple that witnesses
the occurrence of $x$.  Since $(x,x,x)\in A$,
\[
  q_j(x,x,z)=t_j(x,x,x)=x=\pi_1(x,x,z).
\]

Write $c_{j+1}=(a,b,b)$ and set
\[
\begin{aligned}
  X&=\pi_1|_{K_j}=q_j|_{K_j},
  &Y&=p_j|_{K_j},\\
  v&=p_j(c_{j+1})=q_j(c_{j+1}).
\end{aligned}
\]
Project $P$ onto $K_j\cup\{c_{j+1}\}$, ordering $c_{j+1}$ last:
\[
  Q_j=\pr_{K_j\cup\{c_{j+1}\}}(P).
\]
This is an equality-free pp projection of
\eqref{eq:universal-polymorphism-formula}.  Since $A\ne\varnothing$, its
arity is at least two.  The minor $q_j$ is a common polymorphism, so its
restriction belongs to $P$; the same is true of $p_j$ by the induction
invariant and of $\pi_1$ by~\eqref{eq:projection-tables-in-P}.  Their further
restrictions give the three corners
\begin{equation*}
  (X,v),\qquad(Y,v),\qquad(X,a)\in Q_j.
\end{equation*}
\Cref{fig:finite-support-induction} summarizes both the finite lifting step
and the missing-corner argument at this point of the induction.

\begin{figure}[!t]
  \centering
  \resizebox{\linewidth}{!}{%
  \begin{tikzpicture}[
      x=1cm,y=1cm,
      title/.style={font=\small\bfseries},
      panel/.style={draw,rounded corners=2pt,line width=.45pt},
      rel/.style={draw,rounded corners=2pt,minimum width=2.75cm,
                  minimum height=7mm,font=\scriptsize,align=center},
      elem/.style={draw,double,double distance=.45pt,rounded corners=2pt,
                   minimum width=2.75cm,minimum height=7mm,
                   font=\scriptsize,align=center},
      cell/.style={draw,minimum width=2.25cm,minimum height=9mm,
                   font=\scriptsize,align=center},
      forcedcell/.style={cell,double,double distance=.6pt,line width=.65pt},
      arrow/.style={-{Latex[length=2mm]},line width=.55pt},
      lift/.style={-{Latex[length=2mm]},line width=.55pt,dashed}
    ]
    % Panel (a): restrictions, a lift, extension, and a minor.
    \draw[panel] (-.20,5.35) rectangle (14.20,10.75);
    \node[title] at (7.00,10.40) {(a) projection and finite lifting};

    \node[rel,minimum width=3.0cm] (U) at (1.80,9.30)
      {$U_\Delta\subseteq D^T$};
    \node[rel,minimum width=3.4cm] (P) at (7.00,9.30)
      {$P=\pr_{A\cup C}(U_\Delta)$};
    \node[rel,minimum width=3.6cm] (Q) at (12.20,9.30)
      {$Q_j=\pr_{K_j\cup\{c_{j+1}\}}(P)$};
    \draw[arrow] (U)--(P)
      node[midway,above=2pt,font=\scriptsize] {restriction};
    \draw[arrow] (P)--(Q)
      node[midway,above=2pt,font=\scriptsize] {restriction};

    \node[elem,minimum width=3.4cm] (hatp) at (1.80,7.77)
      {$\widehat p_j\in U_\Delta$\\$\widehat p_j|_{A\cup C}=p_j$};
    \node[elem,minimum width=3.0cm] (pj) at (7.00,7.77) {$p_j\in P$};
    \draw[lift] (pj)--(hatp)
      node[midway,above=2pt,font=\scriptsize] {finite lift};
    \draw[arrow] (hatp)--(U);
    \draw[arrow] (pj)--(P);

    \node[elem,minimum width=3.4cm] (tj) at (1.80,6.22)
      {$t_j:D^3\to D$\\extends $\widehat p_j$};
    \node[elem,minimum width=4.2cm] (qj) at (7.20,6.22)
      {$q_j(x,y,z)=t_j(x,y,y)$\\common polymorphism;
       $q_j|_{A\cup C}\in P$};
    \draw[arrow] (hatp)--(tj)
      node[midway,right=2pt,font=\scriptsize] {extend};
    \draw[arrow] (tj)--(qj)
      node[midway,above=2pt,font=\scriptsize] {minor};
    \node[font=\scriptsize,align=center,text width=3.4cm] at (11.95,7.77)
      {restrictions of $p_j$, $q_j$, and $\pi_1$ give the three corners};

    % Panel (b): singleton rectangularity fills the missing corner.
    \draw[panel] (-.20,-.40) rectangle (14.20,5.05);
    \node[title] at (7.00,4.70) {(b) the missing-corner step in $Q_j$};
    \node[font=\scriptsize] at (7.70,4.13)
      {last coordinate $c_{j+1}$};
    \node[font=\scriptsize] at (6.45,3.73) {$v$};
    \node[font=\scriptsize] at (9.15,3.73) {$a$};
    \node[font=\scriptsize] at (4.10,3.18) {row $X$};
    \node[font=\scriptsize] at (4.10,2.13) {row $Y$};

    \node[cell,minimum width=2.5cm] (Xv) at (6.45,3.18)
      {$(X,v)$\\from $q_j$};
    \node[cell,minimum width=2.5cm] (Xa) at (9.15,3.18)
      {$(X,a)$\\from $\pi_1$};
    \node[cell,minimum width=2.5cm] (Yv) at (6.45,2.13)
      {$(Y,v)$\\from $p_j$};
    \node[forcedcell,minimum width=2.5cm] (Ya) at (9.15,2.13)
      {$(Y,a)$\\forced into $Q_j$};

    \node[font=\scriptsize] at (6.25,1.47)
      {the two rows meet at $v$, so rectangularity fills the fourth corner};
    \node[elem,text width=7.0cm] (next) at (7.00,.48)
      {lift $(Y,a)\in Q_j$ to $p_{j+1}\in P$ with
       $p_{j+1}|_{K_j}=Y$ and $p_{j+1}(c_{j+1})=a$};
    % Route the dashed lift down the right side so it does not cover the text.
    \draw[lift] (Ya.south east) to[out=-70,in=70] (next.north east);
  \end{tikzpicture}%
  }
  \caption{One induction step in
    \cref{lem:finite-support-maltsev}.  Top: membership $p_j\in P$
    guarantees a lift $\widehat p_j\in U_\Delta$; because all sets are
    finite, it is found by exhaustive search.  Extending that lift and taking
    the displayed minor supplies the common polymorphism $q_j$.  Solid arrows
    are restrictions, extensions, or minors, whereas dashed arrows are finite
    lifts.  Bottom: in the row view of $Q_j$, the rows $X$ and $Y$ intersect at
    $v$.  Singleton rectangularity makes their supports equal, so the
    double-bordered fourth corner $(Y,a)$ must exist.  Lifting that corner
    produces the next partial table $p_{j+1}$.}
  \label{fig:finite-support-induction}
\end{figure}

By \cref{cor:generated-singleton-rectangularity}, the two nonempty rows
indexed by $X$ and $Y$ have equal support because they intersect at $v$.
Consequently $(Y,a)\in Q_j$.  By the definition of the projection $Q_j$,
this means precisely that there exists $p_{j+1}\in P$ such that
\begin{equation*}
\begin{aligned}
  p_{j+1}|_{K_j}&=Y=p_j|_{K_j},\\
  p_{j+1}(c_{j+1})&=a=\pi_1(c_{j+1}).
\end{aligned}
\end{equation*}
Thus $p_{j+1}$ retains the required values on $A,c_1,\ldots,c_j$ and fixes
the next cell $c_{j+1}$, proving all three invariants at the next step.
This second lift is also a finite exact construction: enumerate
$\widehat p\in D^T$ subject to
\[
  \widehat p|_{K_j}=Y,
  \qquad
  \widehat p(c_{j+1})=a,
  \qquad
  \Psi_\Delta(\widehat p),
\]
and take $p_{j+1}=\widehat p|_{A\cup C}$.  Membership
$(Y,a)\in Q_j$ guarantees that this search terminates.

At $j=s$, choose a lift $\widehat p_s\in U_\Delta$ with
$\widehat p_s|_{A\cup C}=p_s$, using the same finite search.  On every
Mal'tsev-form cell in $T$, this solution has the required value.  Extend it
to $D^3$ by
assigning the required Mal'tsev value to every Mal'tsev-form cell outside
$T$, and assign all remaining cells arbitrarily.  The two prescriptions
overlap only on constant triples and agree there.  Cells outside $T$ occur
in no preservation test, so the resulting operation still preserves every
relation in $\Delta$.  It satisfies
\[
  m(a,b,b)=m(b,b,a)=a
  \qquad(a,b\in D),
\]
and is therefore a common Mal'tsev polymorphism of $\Delta$.
\end{proof}

\subsection{One operation for the entire generated family}

\begin{theorem}[Common Mal'tsev operation for generated supports]
\label{thm:common-support-maltsev}
If $\BO(\cF)$ holds, then
\begin{equation}
  \begin{aligned}
  &\exists m:D^3\to D\quad\forall G\in W_{\cF},\\
  &\qquad
    m\text{ is a Mal'tsev polymorphism of }\SuppRel G.
  \end{aligned}
  \label{eq:global-support-maltsev}
\end{equation}
\end{theorem}

\begin{proof}
Let $\mathcal M_D$ be the finite set of all Mal'tsev operation tables on
$D$.  By \cref{lem:finite-support-maltsev}, every finite family of generated
supports is preserved by some member of $\mathcal M_D$.

Suppose, toward a contradiction, that no member of $\mathcal M_D$ preserves
all generated supports.  For every $m\in\mathcal M_D$, choose
$G_m\in W_{\cF}$ such that $m$ fails to preserve $\SuppRel{G_m}$.  Such a
support is necessarily nonempty.  The finite family
\[
  \Delta=\{\SuppRel{G_m}:m\in\mathcal M_D\}
\]
has a common Mal'tsev polymorphism $m^\star\in\mathcal M_D$ by
\cref{lem:finite-support-maltsev}.  But $\Delta$ contains
$\SuppRel{G_{m^\star}}$, which was chosen not to be preserved by
$m^\star$, a contradiction.  This proves~\eqref{eq:global-support-maltsev}.
\end{proof}

\section{Realizing Row Equivalence as a Generated Support}
\label{sec:omega-realization}

Under global Block Orthogonality, we realize every derived relation
$\RowEq G$ as the support of a table in $W_{\cF}$.  The construction
preserves the zeros forced by Block Orthogonality while avoiding cancellation
on proportional rows.

Let $e_K$ be the exponent of the finite group $\mu(K)$, and define
\begin{equation*}
  K_\mu=\lcm(2,e_K),
  \qquad
  L_D=\lcm(1,\ldots,d).
\end{equation*}
These integers are computable from the input, $K_\mu\ge2$, and
\begin{equation}
  \xi^{K_\mu}=1
  \qquad\text{for every }\xi\in\mu(K).
  \label{eq:Kmu-kills-torsion}
\end{equation}
For a table or vector $H$ and $h\ge1$, let $\Pow{H}{h}$ denote its entrywise
$h$th power: every coordinate $H(\mathbf a)$ is replaced by
$H(\mathbf a)^h$.

\begin{lemma}[Generated entrywise powers]
\label{lem:generated-entrywise-powers}
If $G\in W_{\cF}$ and $h\ge1$, then $\Pow{G}{h}\in W_{\cF}$.
\end{lemma}

\begin{proof}
Choose a presentation
\[
  G(\mathbf a)=\sum_{\mathbf w\in D^s}P(\mathbf a,\mathbf w),
\]
where $P$ is the product of the constraints in a finite $\cF$-instance.
Take $h$ copies of that instance, identify their retained variables
$\mathbf a$, and make their internal tuples
$\mathbf w^{(1)},\ldots,\mathbf w^{(h)}$ pairwise fresh.  Marginalizing them
gives
\[
  \sum_{\mathbf w^{(1)},\ldots,\mathbf w^{(h)}}
     \prod_{j=1}^hP(\mathbf a,\mathbf w^{(j)})
  =G(\mathbf a)^h.
\]
Thus $\Pow{G}{h}\in W_{\cF}$; the construction preserves every zero of $G$,
including those created by cancellation.
\end{proof}

\begin{lemma}[Finite relative row phases]
\label{lem:finite-row-phases}
Assume $\BO(\cF)$.  Let $G:D^r\to K$, $r\ge2$, belong to $W_{\cF}$, and
let $\mathbf u=G(\mathbf x,\ast)$ and
$\mathbf v=G(\mathbf y,\ast)$ be nonzero rows with overlapping supports.
Then their supports are equal.  Fix any anchor $z_0$ in their common support
$S$, and define
\begin{equation}
  \lambda:=\frac{v_{z_0}}{u_{z_0}},
  \qquad
  \zeta_z:=
    \frac{v_z/u_z}{v_{z_0}/u_{z_0}}
    =\frac{v_zu_{z_0}}{u_zv_{z_0}}
  \qquad(z\in S).
  \label{eq:anchored-row-phases}
\end{equation}
Then $\lambda\in K^\times$ and
$(\zeta_z)_{z\in S}\in\mu(K)^S$.  For every $z\in S$,
\begin{equation}
  v_z=\lambda\zeta_z u_z,
  \qquad
  \zeta_{z_0}=1,
  \qquad
  \zeta_z^{L_D}=1.
  \label{eq:finite-row-phases}
\end{equation}
\end{lemma}

The anchor separates one scalar $\lambda$, independent of $z$, from the
coordinate-dependent torsion factors $(\zeta_z)_{z\in S}$.  Although
$\lambda\zeta_z=v_z/u_z\in K^\times$, absorbing this product would produce
a scalar depending on $z$ unless all $\zeta_z$ are equal.  The common
$\lambda$ is what can later be factored out of the detector sum.

\begin{proof}
By \cref{cor:singleton-purification} and Corollary~1 of the Purification Lemma
of Cai and Chen~\cite[Lemma~7 and Corollary~1]{CaiChen2017}, the original
table $G$ is block-orthogonal.
Hence overlapping nonzero rows have equal
supports and proportional magnitude vectors.

By \cref{lem:generated-entrywise-powers}, $\Pow{G}{h}\in W_{\cF}$ for every
$h\ge1$.  Applying the same corollary again shows that each original table
$\Pow{G}{h}$ is block-orthogonal.  Suppose, toward a contradiction,
that $\Pow{\mathbf u}{h}$ and $\Pow{\mathbf v}{h}$ are linearly independent
for every $h\in[d]$.  Fix a nonempty magnitude block $B$ of $\mathbf u$.
Positive powering is injective on positive real numbers, so $B$ remains the
same magnitude block in every $\Pow{\mathbf u}{h}$.  Block orthogonality of
$\Pow{G}{h}$ gives the following power sums.  Put
\[
  b=\abs B,
  \qquad
  x_z=u_z\overline{v_z}\quad(z\in B).
\]
Because $B$ lies in the common support, every $x_z$ is nonzero.  Also
$b\le d$, and hence
\begin{equation*}
  p_h:=\sum_{z\in B}x_z^h
  =\sum_{z\in B}u_z^h\overline{v_z^h}=0,
  \qquad h=1,\ldots,b.
\end{equation*}
Enumerate the $x_z$, $z\in B$, as $x_1,\ldots,x_b$, and let
\[
  e_k=\sum_{1\le i_1<\cdots<i_k\le b}
       x_{i_1}\cdots x_{i_k},
  \qquad e_0=1,
\]
be their elementary symmetric polynomials.  Newton's identities state
\begin{equation}
  k e_k
  =\sum_{j=1}^k(-1)^{j-1}e_{k-j}p_j
  \qquad(k=1,\ldots,b).
  \label{eq:newton-block-identities}
\end{equation}
Since the ambient field has characteristic zero, each integer
$k\in[b]$ is invertible.  As $p_1=\cdots=p_b=0$, induction on $k$ in
\eqref{eq:newton-block-identities} gives
$e_1=\cdots=e_b=0$.  In particular,
\[
  0=e_b=x_1\cdots x_b
  =\prod_{z\in B}u_z\overline{v_z},
\]
contradicting the fact that every factor is nonzero.

Consequently, $\Pow{\mathbf u}{h}$ and $\Pow{\mathbf v}{h}$ are dependent
for some $h\in[d]$.  Thus $(v_z/u_z)^h$ is independent of $z\in S$.  By
\eqref{eq:anchored-row-phases}, for every $z\in S$,
\[
  \zeta_z^h
  =\frac{(v_z/u_z)^h}{(v_{z_0}/u_{z_0})^h}=1.
\]
Since $h$ divides $L_D$, \eqref{eq:finite-row-phases} follows.  All ratios
in~\eqref{eq:anchored-row-phases} belong to $K^\times$, so each
$\zeta_z\in\mu(K)$; the same display gives
$\lambda\in K^\times$ and $\zeta_{z_0}=1$.
\end{proof}

\begin{lemma}[Nondegenerate exponential sums]
\label{lem:nondegenerate-exponential-sums}
Let $s\ge1$, let $a_1,\ldots,a_s,b_1,\ldots,b_s$ be nonzero algebraic
numbers, and
suppose $b_i/b_j$ is not a root of unity whenever $i\ne j$.  Then
\[
  p(t)=\sum_{i=1}^s a_ib_i^t
\]
vanishes for only finitely many integers $t\ge0$.
\end{lemma}

\begin{proof}
The sequence $p(t)$ is a linear recurrence with characteristic roots among
$b_1,\ldots,b_s$.  The Skolem--Mahler--Lech theorem~\cite{Lech1953} says that
its zero set is a finite set together with finitely many arithmetic
progressions.  Suppose one such progression is
$\{r+mn:n\ge N\}$, where $m\ge1$.  Then
\[
  0=p(r+mn)=\sum_{i=1}^s(a_i b_i^r)(b_i^m)^n
\]
for $n=N,N+1,\ldots,N+s-1$.  The bases $b_i^m$ are pairwise distinct because
no $b_i/b_j$ is a root of unity.  The resulting Vandermonde system is
invertible and forces
$(a_i b_i^r)(b_i^m)^N=0$ for every $i$, contradicting the nonzero hypotheses.
Hence no infinite progression occurs.
\end{proof}

\begin{theorem}[Row-equivalence realization]
\label{thm:omega-support-realization}
Assume $\BO(\cF)$.  Let $G:D^r\to K$, $r\ge2$, be any member of
$W_{\cF}$.  There is an integer $t\ge0$ such that, for $q=1+tL_D$, the
$2(r-1)$-ary function
\begin{equation}
  C_q(\mathbf x,\mathbf y)
  =\sum_{z\in D}
       G(\mathbf x,z)^{(K_\mu-1)q}G(\mathbf y,z)^q
  \label{eq:omega-detector}
\end{equation}
belongs to $W_{\cF}$ and satisfies
\begin{equation}
  \SuppRel{C_q}=\RowEq G.
  \label{eq:omega-is-support}
\end{equation}
We call $C_q$ the \emph{row-equivalence detector}; the displayed equality
says that it is nonzero exactly when the rows indexed by $\mathbf x$ and
$\mathbf y$ are both nonzero and linearly dependent.
Given the exact table of $G$ and an instance presentation of it, a terminating
exact search finds such a $t$, the corresponding exponent $q=1+tL_D$, and an
instance presentation of $C_q$.
\end{theorem}

\begin{proof}
Fix an allowed exponent $q=1+tL_D$.  Both exponents
in~\eqref{eq:omega-detector} are positive.  Let
$\mathbf u=G(\mathbf x,\ast)$ and $\mathbf v=G(\mathbf y,\ast)$.  If one row
is zero, or if their supports are disjoint, every summand is zero.
Suppose henceforth that the rows are nonzero and overlap.  Global Block
Orthogonality and support invariance under purification imply that their
supports agree and that their purified magnitude rows are proportional.
Write their common support as
\[
  S=\supp(\mathbf u)=\supp(\mathbf v).
\]

First suppose that $\mathbf u$ and $\mathbf v$ are linearly independent.
Fix a legal singleton purification of $G$, and let $B$ be one magnitude block
of the purified row corresponding to $\mathbf u$.  The two purified
magnitude rows are proportional, so $B$ is simultaneously a magnitude block
of the purified row corresponding to $\mathbf v$.  Choose $z_B\in B$.  By
\cref{lem:intrinsic-purification-tests}, the anchor ratios
\begin{equation}
  \eta_z=u_z/u_{z_B},
  \qquad
  \theta_z=v_z/v_{z_B}
  \qquad(z\in B)
  \label{eq:block-anchor-ratios}
\end{equation}
belong to $\mu(K)$.  Purification preserves complex row dependence, so the
two purified rows remain independent and are block-orthogonal.  Root
covariance then turns the Hermitian sum on $B$ into
\begin{equation}
  \sum_{z\in B}\eta_z\overline{\theta_z}=0.
  \label{eq:purified-root-sum}
\end{equation}

Apply \cref{lem:finite-row-phases} with the same anchor $z_B$:
$v_z=\lambda\zeta_z u_z$, where $\zeta_{z_B}=1$ and
$\zeta_z^{L_D}=1$.  Taking ratios gives $\theta_z=\zeta_z\eta_z$.
Substitution in~\eqref{eq:purified-root-sum}, followed by conjugation, yields
\begin{equation*}
  \sum_{z\in B}\zeta_z=0.
\end{equation*}
Equation~\eqref{eq:block-anchor-ratios} and
\eqref{eq:Kmu-kills-torsion} give
$u_z^{K_\mu}=u_{z_B}^{K_\mu}$.  Since $q\equiv1\pmod{L_D}$, also
$\zeta_z^q=\zeta_z$.  The contribution of $B$ to
$C_q(\mathbf x,\mathbf y)$ is therefore
\begin{align}
  \sum_{z\in B}u_z^{(K_\mu-1)q}v_z^q
  &=\lambda^q\sum_{z\in B}\zeta_z^q u_z^{K_\mu q}\notag\\
  &=\lambda^q u_{z_B}^{K_\mu q}\sum_{z\in B}\zeta_z
   =0.
  \label{eq:independent-block-stays-zero}
\end{align}
The purified magnitude blocks partition the common support; write
\[
  S=B_1\mathbin{\dot\cup}\cdots\mathbin{\dot\cup}B_N.
\]
Terms outside $S$ are zero, and
\eqref{eq:independent-block-stays-zero} applies to every $B_i$.  Therefore
\begin{equation*}
\begin{aligned}
  C_q(\mathbf x,\mathbf y)
  &=\sum_{z\in S}u_z^{(K_\mu-1)q}v_z^q\\
  &=\sum_{i=1}^N
      \sum_{z\in B_i}u_z^{(K_\mu-1)q}v_z^q
    =0.
\end{aligned}
\end{equation*}
Thus the detector vanishes on every overlapping independent ordered row pair
for every allowed $q$.

It remains to treat a nonzero dependent pair, say
$\mathbf v=\lambda\mathbf u$ with $\lambda\ne0$.  Using the common support
$S$, we have
\begin{equation*}
  C_q(\mathbf x,\mathbf y)
   =\lambda^q\sum_{z\in S}u_z^{K_\mu q}.
\end{equation*}
Along $q=1+tL_D$, put $a_z=u_z^{K_\mu}$ and
$b_z=u_z^{K_\mu L_D}$.  The remaining sum is
\begin{equation}
  p_{\mathbf x,\mathbf y}(t)=\sum_{z\in S}a_zb_z^t.
  \label{eq:dependent-exponential-sum}
\end{equation}
Group coordinates having the same base $b_z$.  If $b_z=b_w$, then
$u_z/u_w\in\mu(K)$, so $a_z=a_w$ by
\eqref{eq:Kmu-kills-torsion}.  Thus each grouped coefficient is a positive
integer multiple of one nonzero $a_z$.  Moreover, the quotient of two
distinct grouped bases is not a root of unity: otherwise a power of
$u_z/u_w$ would be one, so $u_z/u_w\in\mu(K)$ and the two bases would have
been equal.  After grouping,
\cref{lem:nondegenerate-exponential-sums} shows that
\eqref{eq:dependent-exponential-sum} has only finitely many zeros.

There are finitely many ordered dependent nonzero row pairs in $G$.  Choose
$t$ outside the union of their finite exceptional sets.  Then the detector
is nonzero on every member of $\RowEq G$, whereas the zero-row,
disjoint-support, and independent cases show that it vanishes on every
nonmember.  This proves~\eqref{eq:omega-is-support}.

The witness $t$ is found by exact search.  Compute $\RowEq G$ exactly: a row
is nonzero if one entry is nonzero, and two
nonzero rows are dependent exactly when
\begin{equation*}
  G(\mathbf x,z)G(\mathbf y,w)
  -G(\mathbf x,w)G(\mathbf y,z)=0
  \qquad(z,w\in D).
\end{equation*}
Enumerate $t=0,1,2,\ldots$, evaluate~\eqref{eq:omega-detector} by exact
algebraic arithmetic, and stop when its support equals $\RowEq G$.  The
finite-exception argument proves termination.

Finally, choose a presentation
\[
  G(\mathbf x,z)
   =\sum_{\mathbf h\in D^s}P(\mathbf x,z,\mathbf h)
\]
by a finite $\cF$-instance, and set $A=(K_\mu-1)q$.  Take $A$
fresh-internal copies of this presentation sharing only $(\mathbf x,z)$ and
$q$ fresh-internal copies sharing only $(\mathbf y,z)$.  Make every copied
internal tuple fresh.  For fixed $(\mathbf x,\mathbf y,z)$, marginalizing
the two collections of copied hidden variables gives
\begin{align*}
  \left(\sum_{\mathbf h}P(\mathbf x,z,\mathbf h)\right)^A
    &=G(\mathbf x,z)^A,\\
  \left(\sum_{\mathbf k}P(\mathbf y,z,\mathbf k)\right)^q
    &=G(\mathbf y,z)^q.
\end{align*}
Because the collections are disjoint, their contributions multiply.
Summing over $z$ now gives exactly $C_q(\mathbf x,\mathbf y)$.
Thus $C_q\in W_{\cF}$.  The construction preserves repeated scopes and
previous boundary identifications, and it treats every earlier marginal by
fresh copies of its hidden variables.
\end{proof}

\begin{corollary}[Support preservation extends to all row relations]
\label{cor:omega-reduces-to-supports}
Assume $\BO(\cF)$.  If one operation $m:D^3\to D$ preserves $\SuppRel H$
for every $H\in W_{\cF}$, then the same operation preserves $\RowEq G$ for
every $G\in W_{\cF}$ of arity at least two.
\end{corollary}

\begin{proof}
Fix the one operation $m$ from the hypothesis, and let
$G\in W_{\cF}$ have arity at least two.  By
\cref{thm:omega-support-realization}, there is a witness $t_G\ge0$, with
$q_G=1+t_GL_D$, for which the detector table constructed from $G$ belongs
to $W_{\cF}$.  Denote this table by $C_G$.  Then
\[
  \SuppRel{C_G}=\RowEq G.
\]
The hypothesis says that the fixed operation $m$ preserves the support of
every member of $W_{\cF}$, so it preserves $\SuppRel{C_G}=\RowEq G$.
Since $G$ was arbitrary, this same $m$ preserves every generated row
relation.

Although $t_G$ and $C_G$ may depend on $G$, the operation $m$ is fixed.
Hence the quantifier order is $\exists m\,\forall G$.
\end{proof}

\section{Type Partition and the Structural Collapse}
\label{sec:type-collapse}

For equal-length tuples $\mathbf u,\mathbf v,\mathbf w$, we write
$m(\mathbf u,\mathbf v,\mathbf w)$ for coordinatewise application of $m$.

\begin{lemma}[Row equivalence forces Type Partition]
\label{lem:omega-implies-type}
Let $G:D^r\to\AlgNums$, $r\ge2$, and let $m:D^3\to D$ be a Mal'tsev
operation preserving $\RowEq G$.  Then the type map of $G$ is a
type-partition map.
\end{lemma}

\begin{proof}
Write
\[
  \Row(G)=\{(S_1,\mathbf v_1),\ldots,(S_s,\mathbf v_s)\}.
\]
Fix $\ell\in[r-1]$ and
$\boldsymbol\alpha,\boldsymbol\beta\in D^\ell$.  Suppose their types
intersect.  Choose
\[
  j\in\typeop_G(\boldsymbol\alpha)
       \cap\typeop_G(\boldsymbol\beta)
\]
and rows $\mathbf a,\mathbf b\in S_j$ extending
$\boldsymbol\alpha,\boldsymbol\beta$, respectively.

Let $h\in\typeop_G(\boldsymbol\alpha)$ be arbitrary, and choose
$\mathbf c\in S_h$ extending $\boldsymbol\alpha$.  The three
$2(r-1)$-tuples
\[
  (\mathbf c,\mathbf c),\qquad
  (\mathbf a,\mathbf a),\qquad
  (\mathbf b,\mathbf a)
\]
belong to $\RowEq G$: the first two are diagonal pairs of nonzero rows, and
the third belongs because $\mathbf a$ and $\mathbf b$ lie in the same row
class $S_j$.  Put $\mathbf d=m(\mathbf c,\mathbf a,\mathbf b)$.  Applying
$m$ coordinatewise to these three members of $\RowEq G$ gives
\[
\begin{aligned}
  &m\bigl((\mathbf c,\mathbf c),
          (\mathbf a,\mathbf a),
          (\mathbf b,\mathbf a)\bigr)\\
  &\qquad=
  \bigl(m(\mathbf c,\mathbf a,\mathbf b),
        m(\mathbf c,\mathbf a,\mathbf a)\bigr)
  =(\mathbf d,\mathbf c),
\end{aligned}
\]
where the last equality uses the Mal'tsev identity $m(x,y,y)=x$ in every
coordinate.  Preservation therefore gives
\begin{equation*}
  (\mathbf d,\mathbf c)\in\RowEq G.
\end{equation*}
For every $i\in[\ell]$, the $i$th coordinate of the first component is
\[
  m(\alpha_i,\alpha_i,\beta_i)=\beta_i.
\]
Thus the row indexed by $\mathbf d$ extends $\boldsymbol\beta$.
Membership in $\RowEq G$ also says that this row is
nonzero and belongs to the same original complex row class $S_h$ as
$\mathbf c$.  Therefore $h\in\typeop_G(\boldsymbol\beta)$.  Since $h$ was
arbitrary,
\[
  \typeop_G(\boldsymbol\alpha)
  \subseteq\typeop_G(\boldsymbol\beta).
\]
Interchanging the prefixes proves the reverse inclusion.  Thus any two types
are equal or disjoint.  Empty types and the zero table are covered as well.
\end{proof}

\begin{theorem}[Structural collapse]
\label{thm:structural-collapse}
For every nonempty finite domain $D$ and every finite algebraic-complex
constraint language $\cF$,
\[
  \BO(\cF)\quad\Longrightarrow\quad\TP(\cF)\wedge\Mal(\cF).
\]
More explicitly, under $\BO(\cF)$ there is one Mal'tsev operation
$m:D^3\to D$ such that, for every $G\in W_{\cF}$,
\begin{equation*}
  \begin{aligned}
  &m\text{ preserves }\SuppRel G,\\
  &\arity(G)\ge2
    \quad\Longrightarrow\quad
    m\text{ preserves }\RowEq G.
  \end{aligned}
\end{equation*}
Consequently,
\[
  \CC(\cF)\quad\Longleftrightarrow\quad\BO(\cF).
\]
\end{theorem}

\begin{proof}
The empty-language case was handled in \cref{sec:preliminaries}; assume
$\cF\ne\varnothing$.
Assume $\BO(\cF)$.  By \cref{thm:common-support-maltsev}, there is one
Mal'tsev operation $m:D^3\to D$ preserving $\SuppRel H$ for every
$H\in W_{\cF}$.

Fix $G\in W_{\cF}$ of arity $r\ge2$.  By
\cref{thm:omega-support-realization}, there exists a generated table
\[
  C_G:D^{2(r-1)}\to\AlgNums,
  \qquad C_G\in W_{\cF},
\]
such that $\SuppRel{C_G}=\RowEq G$.  Since $m$ preserves every generated
support, it preserves $\RowEq G=\SuppRel{C_G}$.  Hence the same $m$ preserves
every relation in $\Lambda_{\cF}$, proving $\Mal(\cF)$.

Applying \cref{lem:omega-implies-type} to every generated table of arity at
least two proves $\TP(\cF)$.  The argument applies in every generated
coordinate ordering; relabelling the retained variables preserves membership
in $W_{\cF}$.

We have proved $\BO(\cF)\Rightarrow\TP(\cF)\wedge\Mal(\cF)$.  The reverse
implication in the final equivalence is immediate because $\BO(\cF)$ is a
conjunct of $\CC(\cF)$.
\end{proof}

\subsection{Proof of uniform decidability}

\begin{proof}[Proof of \cref{thm:main}]
On input $(D,\cF)$, run \cref{proc:global-bo}.  By
\cref{thm:global-bo-decidable}, this is a total exact procedure and returns
\YES{} if and only if $\BO(\cF)$ holds.  By
\cref{thm:structural-collapse},
\[
  \BO(\cF)\quad\Longleftrightarrow\quad\CC(\cF).
\]
Consequently, the same procedure halts on every input and returns
\YES{} exactly when $\CC(\cF)$ holds, proving both uniform
decidability and~\eqref{eq:main-equivalence}.
\end{proof}

\begin{proof}[Proof of \cref{cor:degree-multiple-main}]
Fix $\delta\ge1$.  The proof of \cref{cor:singleton-purification} restricts
verbatim to $\WFdeg{\delta}$: $\BO^\delta(\cF)$ holds if and only if every
member of $\WFdeg{\delta}$ of arity at least two has a block-orthogonal
singleton purification.  At arity $k$, use the labelled submonoid
$\mathfrak M_{k,\delta}$ and the oracle in
\cref{thm:degree-multiple-universal-identity}.  The purification-invariance
and finite-certificate arguments concern one finite table and are unchanged.
The external-arity compression also remains valid: identifying retained
variables replaces their occurrence degrees by their sum, which is still
divisible by $\delta$.  Hence \cref{proc:global-bo}, using the
degree-multiple oracle, is a total exact decision procedure for
$\BO^\delta(\cF)$.

The structural proof likewise remains inside $\WFdeg{\delta}$.  More
explicitly, every instance substitution used in
\cref{lem:equality-free-support-realization,lem:generated-entrywise-powers,thm:omega-support-realization}
takes copies of degree-multiple presentations, keeps their hidden variables
fresh, and identifies some boundary variables.  A fresh variable retains a
degree divisible by $\delta$, while an identified variable receives a sum
of such degrees.  Labelled products, diagonal minors, coordinate
permutations, and repeated boundary positions obey the same rule.  Thus all
support tables and row-equivalence detectors used in
\crefrange{sec:generated-supports}{sec:type-collapse} are members of
$\WFdeg{\delta}$.  Repeating the proofs gives
\[
  \BO^\delta(\cF)
  \Longrightarrow
  \TP^\delta(\cF)\wedge\Mal^\delta(\cF),
\]
with one Mal'tsev operation preserving the entire degree-multiple generated
family.  Since $\BO^\delta$ is a conjunct of $\CC^\delta$, this proves the
claimed equivalence and decidability.  Together with the dichotomy theorem
for $\#\mathrm{CSP}^{\delta}$ proved by Lin~\cite[Theorem~3]{Lin2021}, this
decides the polynomial-time versus $\#\mathsf P$-hard alternative for
$\CSPdeg{\delta}{\cF}$.
\end{proof}

\section{Conclusion}
\label{sec:conclusion}

We have given a total exact decision procedure for the three tractability
conditions characterizing the polynomial-time side of the complex-weighted
$\#\mathrm{CSP}$ dichotomy theorem of Cai and
Chen~\cite[Theorem~1 and Section~3.1]{CaiChen2017}.  The procedure decides
these conditions on the full unbounded family $W_{\cF}$.
Its finite reduction combines the external-arity bound $d^2+1$ with universal
identity testing over all finite labelled instances.

Global Block Orthogonality moreover forces the Type Partition and common
Mal'tsev conditions from that dichotomy theorem, with a single operation
satisfying
$\exists m\,\forall G$.  The support-amplification and row-equivalence
detector constructions remain within $W_{\cF}$ and handle
cancellation-created zeros exactly.

A cyclic root-of-unity degree filter extends both recognition and the
structural collapse to the degree-divisibility restriction
$\#\mathrm{CSP}^{\delta}$ studied by Lin~\cite[Definition~2 and
Theorem~3]{Lin2021}.  Thus the three infinitary tractability conditions on
$W_{\cF}^{\delta}$ in that dichotomy theorem are uniformly decidable for every
positive degree modulus $\delta$.

\bibliographystyle{alpha}
\bibliography{cai-chen-complex-csp-decidability}
\end{document}
